% Recovered from the previously completed PDF; no new mathematical content.
% Prose is editable; formula boxes retain original symbol positions.
% Build from repository root: xelatex -no-shell-escape -halt-on-error -output-directory=... print/fudan/NAME.tex
\documentclass[11pt,a4paper]{article}
\usepackage[margin=23mm,headheight=14pt,headsep=8.5mm,footskip=10mm]{geometry}
\usepackage{amsmath,amssymb}
\DeclareMathAlphabet{\mathsf}{OT1}{cmss}{m}{n}
\usepackage[no-math]{fontspec}
\usepackage{xeCJK}
\setmainfont{Latin Modern Roman}
\setCJKmainfont{FandolSong-Regular.otf}[Path=assets/fonts/fandol/,BoldFont=FandolSong-Bold.otf,ItalicFont=FandolSong-Regular.otf,Script=CustomDefault]
\xeCJKsetup{PunctStyle=quanjiao,CheckSingle=true}
\usepackage{setspace}\setstretch{1.12}
\setlength{\parindent}{2em}\setlength{\parskip}{3pt}
\emergencystretch=6em
\hfuzz=9pt
\usepackage{titlesec}
\titleformat{\section}{\centering\normalfont\Large\bfseries}{}{0pt}{}
\titlespacing*{\section}{0pt}{1.8\baselineskip}{1.0\baselineskip}
\usepackage{fancyhdr}\pagestyle{fancy}\fancyhf{}
\fancyhead[L]{\small 数学手记·高等代数}\fancyhead[R]{\small 19S15–19S15、19A01–19A06}
\fancyfoot[C]{\thepage}\renewcommand{\headrulewidth}{0.4pt}
\usepackage[hidelinks,unicode]{hyperref}
\hypersetup{pdftitle={19S15–19S15、19A01–19A06：题目、解答与小结},pdfauthor={Vincent Arqain}}
\newcommand{\RB}[4]{\raisebox{0pt}[#2bp][#3bp]{\hbox to #1bp{#4\hss}}}
\newcommand{\RG}[5]{\hbox to 0pt{\kern#1bp\raisebox{#2bp}[0pt][0pt]{\fontsize{#3bp}{#3bp}\csname font#4\endcsname\char#5}\hss}}
\newcommand{\RT}[4]{\hbox to 0pt{\kern#1bp\raisebox{#2bp}[0pt][0pt]{\fontsize{#3bp}{#3bp}\selectfont#4}\hss}}
\newcommand{\RS}[4]{\hbox to 0pt{\kern#1bp\raisebox{#2bp}[0pt][0pt]{\fontsize{#3bp}{#3bp}\selectfont\ensuremath{#4}}\hss}}
\newcommand{\RR}[4]{\hbox to 0pt{\kern#1bp\raisebox{#2bp}[0pt][0pt]{\rule{#3bp}{#4bp}}\hss}}
\newcommand{\fontmi}{\usefont{OML}{cmm}{m}{it}}
\newcommand{\fontsy}{\usefont{OMS}{cmsy}{m}{n}}
\newcommand{\fontex}{\usefont{OMX}{cmex}{m}{n}}
\newcommand{\fontrm}{\usefont{OT1}{cmr}{m}{n}}
\newcommand{\fontbf}{\usefont{OT1}{cmr}{bx}{n}}
\newcommand{\fontam}{\usefont{U}{msa}{m}{n}}
\newcommand{\fontbb}{\usefont{U}{msb}{m}{n}}
\widowpenalty=10000\clubpenalty=10000\raggedbottom
\begin{document}
\begin{center}{\LARGE\bfseries 19S15–19S15、19A01–19A06\par}\vspace{8pt}复旦大学高等代数习题集·题目与解答\end{center}
% Original page 1.
\section*{2019S15}
设\allowbreak{}\hskip.1em \RB{8.889}{7.951}{0.740}{%
\RG{0.000}{-0.000}{10.9091}{mi}{86}%
}\hskip.1em\allowbreak{} 是\allowbreak{}\hskip.1em \RB{6.729}{5.322}{0.620}{%
\RG{0.000}{-0.000}{10.9091}{mi}{110}%
}\hskip.1em\allowbreak{} 维欧氏空间，向量组\allowbreak{}\hskip.1em \RB{61.534}{8.682}{3.227}{%
\RG{0.000}{-0.000}{10.9091}{sy}{102}\RG{5.455}{-0.000}{10.9091}{mi}{11}\RG{12.433}{-1.636}{7.9701}{rm}{49}\RG{17.166}{-0.000}{10.9091}{mi}{59}\RG{21.999}{-0.000}{10.9091}{mi}{58}\RG{26.843}{-0.000}{10.9091}{mi}{58}\RG{31.675}{-0.000}{10.9091}{mi}{58}\RG{36.519}{-0.000}{10.9091}{mi}{59}\RG{41.363}{-0.000}{10.9091}{mi}{11}\RG{48.387}{-1.636}{7.9701}{mi}{109}\RG{56.375}{-0.000}{10.9091}{sy}{103}%
}\hskip.1em\allowbreak{} 满足：若非负实数\allowbreak{}\hskip.1em \RB{49.431}{8.071}{2.605}{%
\RG{0.000}{-0.000}{10.9091}{mi}{21}\RG{6.364}{-1.636}{7.9701}{rm}{49}\RG{11.096}{-0.000}{10.9091}{mi}{59}\RG{15.929}{-0.000}{10.9091}{mi}{58}\RG{20.773}{-0.000}{10.9091}{mi}{58}\RG{25.605}{-0.000}{10.9091}{mi}{58}\RG{30.449}{-0.000}{10.9091}{mi}{59}\RG{35.293}{-0.000}{10.9091}{mi}{21}\RG{41.702}{-1.636}{7.9701}{mi}{109}%
}\hskip.1em\allowbreak{} 使得

\[
\RB{114.235}{8.067}{2.609}{%
\RG{0.000}{-0.004}{10.9091}{mi}{21}\RG{11.096}{-0.004}{10.9091}{mi}{11}\RG{25.231}{-0.004}{10.9091}{rm}{43}\RG{36.140}{-0.004}{10.9091}{sy}{1}\RG{40.973}{-0.004}{10.9091}{sy}{1}\RG{45.817}{-0.004}{10.9091}{sy}{1}\RG{51.291}{-0.004}{10.9091}{rm}{43}\RG{62.201}{-0.004}{10.9091}{mi}{21}\RG{76.554}{-0.004}{10.9091}{mi}{11}\RG{94.551}{-0.004}{10.9091}{rm}{61}\RG{106.049}{-0.004}{10.9091}{rm}{48}\RG{111.521}{-0.004}{10.9091}{mi}{59}\RG{6.363}{-1.640}{7.9701}{rm}{49}\RG{18.074}{-1.640}{7.9701}{rm}{49}\RG{68.565}{-1.640}{7.9701}{mi}{109}\RG{83.532}{-1.640}{7.9701}{mi}{109}%
}
\]
则必有\allowbreak{}\hskip.1em \RB{87.313}{8.073}{2.214}{%
\RG{0.000}{0.002}{10.9091}{mi}{21}\RG{6.364}{-1.634}{7.9701}{rm}{49}\RG{14.126}{0.002}{10.9091}{rm}{61}\RG{25.641}{0.002}{10.9091}{sy}{1}\RG{30.474}{0.002}{10.9091}{sy}{1}\RG{35.318}{0.002}{10.9091}{sy}{1}\RG{41.399}{0.002}{10.9091}{rm}{61}\RG{52.914}{0.002}{10.9091}{mi}{21}\RG{59.277}{-1.634}{7.9701}{mi}{109}\RG{70.297}{0.002}{10.9091}{rm}{61}\RG{81.795}{0.002}{10.9091}{rm}{48}%
}\hskip.1em\allowbreak{} 。证明：存在\allowbreak{}\hskip.1em \RB{29.242}{7.953}{0.934}{%
\RG{0.000}{0.002}{10.9091}{mi}{11}\RG{10.050}{0.002}{10.9091}{sy}{50}\RG{20.353}{0.002}{10.9091}{mi}{86}%
}\hskip.1em\allowbreak{} ，使得

\[
\RB{131.493}{8.680}{3.229}{%
\RG{0.000}{-0.002}{10.9091}{rm}{40}\RG{4.243}{-0.002}{10.9091}{mi}{11}\RG{11.247}{-0.002}{10.9091}{mi}{59}\RG{16.090}{-0.002}{10.9091}{mi}{11}\RG{26.470}{-0.002}{10.9091}{rm}{41}\RG{33.743}{-0.002}{10.9091}{mi}{62}\RG{45.258}{-0.002}{10.9091}{rm}{48}\RG{72.520}{-0.002}{10.9091}{rm}{40}\RG{76.753}{-0.002}{10.9091}{rm}{49}\RG{85.259}{-0.002}{10.9091}{sy}{20}\RG{96.774}{-0.002}{10.9091}{mi}{105}\RG{103.562}{-0.002}{10.9091}{sy}{20}\RG{115.078}{-0.002}{10.9091}{mi}{109}\RG{124.656}{-0.002}{10.9091}{rm}{41}\RG{128.898}{-0.002}{10.9091}{mi}{58}\RG{23.089}{-1.638}{7.9701}{mi}{105}%
}
\]
证明：考虑这些向量的凸包

\[
\RB{181.067}{19.589}{14.128}{%
\RG{23.123}{18.652}{10.9092}{ex}{40}\RG{167.866}{18.652}{10.9092}{ex}{41}\RG{36.044}{13.634}{7.9701}{mi}{109}\RG{127.103}{13.634}{7.9701}{mi}{109}\RG{31.911}{10.361}{10.9092}{ex}{88}\RG{122.970}{10.361}{10.9092}{ex}{88}\RG{0.000}{-0.003}{10.9091}{mi}{67}\RG{11.608}{-0.003}{10.9091}{rm}{61}\RG{49.487}{-0.003}{10.9091}{mi}{116}\RG{56.807}{-0.003}{10.9091}{mi}{11}\RG{70.198}{-0.003}{10.9091}{rm}{58}\RG{79.892}{-0.003}{10.9091}{mi}{116}\RG{90.243}{-0.003}{10.9091}{sy}{21}\RG{101.758}{-0.003}{10.9091}{rm}{48}\RG{107.212}{-0.003}{10.9091}{mi}{59}\RG{140.546}{-0.003}{10.9091}{mi}{116}\RG{150.896}{-0.003}{10.9091}{rm}{61}\RG{162.395}{-0.003}{10.9091}{rm}{49}\RG{178.472}{-0.003}{10.9091}{mi}{58}\RG{53.426}{-1.639}{7.9701}{mi}{105}\RG{63.787}{-1.639}{7.9701}{mi}{105}\RG{83.831}{-1.639}{7.9701}{mi}{105}\RG{144.485}{-1.639}{7.9701}{mi}{105}\RG{32.938}{-13.090}{7.9701}{mi}{105}\RG{35.821}{-13.090}{7.9701}{rm}{61}\RG{42.405}{-13.090}{7.9701}{rm}{49}\RG{123.997}{-13.090}{7.9701}{mi}{105}\RG{126.880}{-13.090}{7.9701}{rm}{61}\RG{133.463}{-13.090}{7.9701}{rm}{49}%
}
\]
由题设，\allowbreak{}\hskip.1em \RB{27.579}{8.680}{3.229}{%
\RG{0.000}{-0.002}{10.9091}{rm}{48}\RG{8.485}{-0.002}{10.9091}{sy}{50}\RG{9.697}{-0.002}{10.9091}{mi}{61}\RG{18.788}{-0.002}{10.9091}{mi}{67}%
}\hskip.1em\allowbreak{} 。集合

\[
\RB{223.363}{19.588}{14.129}{%
\RG{23.636}{18.651}{10.9092}{ex}{40}\RG{215.652}{18.651}{10.9092}{ex}{41}\RG{174.888}{13.633}{7.9701}{mi}{109}\RG{170.755}{10.360}{10.9092}{ex}{88}\RG{106.964}{4.501}{7.9701}{mi}{109}\RG{0.000}{-0.003}{10.9091}{rm}{1}\RG{12.109}{-0.003}{10.9091}{rm}{61}\RG{32.424}{-0.003}{10.9091}{rm}{40}\RG{36.667}{-0.003}{10.9091}{mi}{116}\RG{45.338}{-0.003}{10.9091}{mi}{59}\RG{50.171}{-0.003}{10.9091}{mi}{58}\RG{55.015}{-0.003}{10.9091}{mi}{58}\RG{59.858}{-0.003}{10.9091}{mi}{58}\RG{64.691}{-0.003}{10.9091}{mi}{59}\RG{69.535}{-0.003}{10.9091}{mi}{116}\RG{81.510}{-0.003}{10.9091}{rm}{41}\RG{88.782}{-0.003}{10.9091}{sy}{50}\RG{99.085}{-0.003}{10.9091}{bb}{82}\RG{117.983}{-0.003}{10.9091}{rm}{58}\RG{127.676}{-0.003}{10.9091}{mi}{116}\RG{138.027}{-0.003}{10.9091}{sy}{21}\RG{149.542}{-0.003}{10.9091}{rm}{48}\RG{154.997}{-0.003}{10.9091}{mi}{59}\RG{188.331}{-0.003}{10.9091}{mi}{116}\RG{198.682}{-0.003}{10.9091}{rm}{61}\RG{210.180}{-0.003}{10.9091}{rm}{49}\RG{40.606}{-1.640}{7.9701}{rm}{49}\RG{73.521}{-1.640}{7.9701}{mi}{109}\RG{131.616}{-1.640}{7.9701}{mi}{105}\RG{192.271}{-1.640}{7.9701}{mi}{105}\RG{171.781}{-13.090}{7.9701}{mi}{105}\RG{174.665}{-13.090}{7.9701}{rm}{61}\RG{181.249}{-13.090}{7.9701}{rm}{49}%
}
\]
闭且有界，因而紧致；\allowbreak{}\hskip.1em \RB{8.791}{8.195}{0.736}{%
\RG{0.000}{0.004}{10.9091}{mi}{67}%
}\hskip.1em\allowbreak{} 是\allowbreak{}\hskip.1em \RB{9.064}{8.315}{0.500}{%
\RG{0.000}{0.004}{10.9091}{rm}{1}%
}\hskip.1em\allowbreak{} 在连续映射\allowbreak{}\hskip.1em \RB{106.324}{8.686}{3.801}{%
\RG{0.000}{0.004}{10.9091}{rm}{40}\RG{4.242}{0.004}{10.9091}{mi}{116}\RG{8.182}{-1.632}{7.9701}{rm}{49}\RG{12.914}{0.004}{10.9091}{mi}{59}\RG{17.747}{0.004}{10.9091}{mi}{58}\RG{22.591}{0.004}{10.9091}{mi}{58}\RG{27.423}{0.004}{10.9091}{mi}{58}\RG{32.267}{0.004}{10.9091}{mi}{59}\RG{37.111}{0.004}{10.9091}{mi}{116}\RG{41.096}{-1.632}{7.9701}{mi}{109}\RG{49.084}{0.004}{10.9091}{rm}{41}\RG{57.350}{0.004}{10.9091}{sy}{55}\RG{57.350}{0.004}{10.9091}{sy}{33}\RG{72.284}{8.186}{10.9092}{ex}{80}\RG{83.799}{-3.221}{7.9701}{mi}{105}\RG{88.999}{0.004}{10.9091}{mi}{116}\RG{92.938}{-1.632}{7.9701}{mi}{105}\RG{96.319}{0.004}{10.9091}{mi}{11}\RG{103.298}{-1.632}{7.9701}{mi}{105}%
}\hskip.1em\allowbreak{} 下的像，故也紧致。连续函数\allowbreak{}\hskip.1em \RB{44.650}{9.758}{3.228}{%
\RG{0.000}{-0.001}{10.9091}{mi}{120}\RG{9.265}{-0.001}{10.9091}{sy}{55}\RG{9.265}{-0.001}{10.9091}{sy}{33}\RG{23.196}{-0.001}{10.9091}{sy}{107}\RG{28.659}{-0.001}{10.9091}{mi}{120}\RG{34.894}{-0.001}{10.9091}{sy}{107}\RG{40.348}{3.958}{7.9701}{rm}{50}%
}\hskip.1em\allowbreak{} 在\allowbreak{}\hskip.1em \RB{8.791}{8.190}{0.741}{%
\RG{0.000}{-0.001}{10.9091}{mi}{67}%
}\hskip.1em\allowbreak{} 上取得最小值。取极小点\allowbreak{}\hskip.1em \RB{29.143}{8.190}{0.937}{%
\RG{0.000}{-0.001}{10.9091}{mi}{11}\RG{10.049}{-0.001}{10.9091}{sy}{50}\RG{20.352}{-0.001}{10.9091}{mi}{67}%
}\hskip.1em\allowbreak{} ，则

\[
\RB{45.255}{10.300}{3.232}{%
\RG{17.928}{4.500}{7.9701}{rm}{50}\RG{0.000}{-0.005}{10.9091}{sy}{107}\RG{5.454}{-0.005}{10.9091}{mi}{11}\RG{12.473}{-0.005}{10.9091}{sy}{107}\RG{25.690}{-0.005}{10.9091}{mi}{62}\RG{37.205}{-0.005}{10.9091}{rm}{48}\RG{42.660}{-0.005}{10.9091}{mi}{58}%
}
\]
固定\allowbreak{}\hskip.1em \RB{66.980}{8.682}{3.227}{%
\RG{0.000}{0.000}{10.9091}{mi}{105}\RG{6.789}{0.000}{10.9091}{sy}{50}\RG{17.076}{0.000}{10.9091}{sy}{102}\RG{22.546}{0.000}{10.9091}{rm}{49}\RG{28.001}{0.000}{10.9091}{mi}{59}\RG{32.834}{0.000}{10.9091}{mi}{58}\RG{37.678}{0.000}{10.9091}{mi}{58}\RG{42.510}{0.000}{10.9091}{mi}{58}\RG{47.354}{0.000}{10.9091}{mi}{59}\RG{52.198}{0.000}{10.9091}{mi}{109}\RG{61.822}{0.000}{10.9091}{sy}{103}%
}\hskip.1em\allowbreak{} 。对任意\allowbreak{}\hskip.1em \RB{38.295}{8.682}{3.227}{%
\RG{0.000}{0.000}{10.9091}{mi}{116}\RG{6.970}{0.000}{10.9091}{sy}{50}\RG{17.273}{0.000}{10.9091}{rm}{91}\RG{20.295}{0.000}{10.9091}{rm}{48}\RG{25.757}{0.000}{10.9091}{mi}{59}\RG{30.606}{0.000}{10.9091}{rm}{49}\RG{36.061}{0.000}{10.9091}{rm}{93}%
}\hskip.1em\allowbreak{} ，凸性给出

\[
\RB{167.780}{8.678}{3.231}{%
\RG{0.000}{-0.004}{10.9091}{mi}{11}\RG{9.443}{-0.004}{10.9091}{rm}{43}\RG{20.352}{-0.004}{10.9091}{mi}{116}\RG{24.291}{-0.004}{10.9091}{rm}{40}\RG{28.534}{-0.004}{10.9091}{mi}{11}\RG{41.318}{-0.004}{10.9091}{sy}{0}\RG{52.227}{-0.004}{10.9091}{mi}{11}\RG{59.246}{-0.004}{10.9091}{rm}{41}\RG{66.501}{-0.004}{10.9091}{rm}{61}\RG{78.010}{-0.004}{10.9091}{rm}{40}\RG{82.243}{-0.004}{10.9091}{rm}{49}\RG{90.156}{-0.004}{10.9091}{sy}{0}\RG{101.065}{-0.004}{10.9091}{mi}{116}\RG{105.004}{-0.004}{10.9091}{rm}{41}\RG{109.247}{-0.004}{10.9091}{mi}{11}\RG{118.690}{-0.004}{10.9091}{rm}{43}\RG{129.599}{-0.004}{10.9091}{mi}{116}\RG{133.537}{-0.004}{10.9091}{mi}{11}\RG{146.929}{-0.004}{10.9091}{sy}{50}\RG{157.232}{-0.004}{10.9091}{mi}{67}\RG{165.185}{-0.004}{10.9091}{mi}{58}\RG{35.512}{-1.641}{7.9701}{mi}{105}\RG{140.517}{-1.641}{7.9701}{mi}{105}%
}
\]
由\allowbreak{}\hskip.1em \RB{7.067}{5.324}{0.618}{%
\RG{0.000}{0.002}{10.9091}{mi}{11}%
}\hskip.1em\allowbreak{} 的极小性，

\[
\RB{151.413}{28.472}{3.223}{%
\RG{94.399}{22.672}{7.9701}{rm}{50}\RG{130.393}{22.672}{7.9701}{rm}{50}\RG{0.000}{18.168}{10.9091}{rm}{48}\RG{8.485}{18.168}{10.9091}{sy}{20}\RG{19.983}{18.168}{10.9091}{sy}{107}\RG{25.455}{18.168}{10.9091}{mi}{11}\RG{34.898}{18.168}{10.9091}{rm}{43}\RG{45.807}{18.168}{10.9091}{mi}{116}\RG{49.747}{18.168}{10.9091}{rm}{40}\RG{53.989}{18.168}{10.9091}{mi}{11}\RG{66.774}{18.168}{10.9091}{sy}{0}\RG{77.683}{18.168}{10.9091}{mi}{11}\RG{84.702}{18.168}{10.9091}{rm}{41}\RG{88.945}{18.168}{10.9091}{sy}{107}\RG{101.556}{18.168}{10.9091}{sy}{0}\RG{112.454}{18.168}{10.9091}{sy}{107}\RG{117.919}{18.168}{10.9091}{mi}{11}\RG{124.938}{18.168}{10.9091}{sy}{107}\RG{60.968}{16.532}{7.9701}{mi}{105}\RG{97.732}{4.508}{7.9701}{rm}{50}\RG{144.086}{4.508}{7.9701}{rm}{50}\RG{8.485}{0.004}{10.9091}{rm}{61}\RG{19.983}{0.004}{10.9091}{rm}{50}\RG{25.455}{0.004}{10.9091}{mi}{116}\RG{29.394}{0.004}{10.9091}{rm}{40}\RG{33.637}{0.004}{10.9091}{mi}{11}\RG{40.641}{0.004}{10.9091}{mi}{59}\RG{45.485}{0.004}{10.9091}{mi}{11}\RG{58.288}{0.004}{10.9091}{sy}{0}\RG{69.198}{0.004}{10.9091}{mi}{11}\RG{76.217}{0.004}{10.9091}{rm}{41}\RG{82.871}{0.004}{10.9091}{rm}{43}\RG{93.793}{0.004}{10.9091}{mi}{116}\RG{102.465}{0.004}{10.9091}{sy}{107}\RG{107.919}{0.004}{10.9091}{mi}{11}\RG{120.703}{0.004}{10.9091}{sy}{0}\RG{131.612}{0.004}{10.9091}{mi}{11}\RG{138.631}{0.004}{10.9091}{sy}{107}\RG{148.818}{0.004}{10.9091}{mi}{58}\RG{52.483}{-1.632}{7.9701}{mi}{105}\RG{114.898}{-1.632}{7.9701}{mi}{105}%
}
\]
当\allowbreak{}\hskip.1em \RB{24.003}{7.766}{0.936}{%
\RG{0.000}{0.000}{10.9091}{mi}{116}\RG{6.960}{0.000}{10.9091}{mi}{62}\RG{18.485}{0.000}{10.9091}{rm}{48}%
}\hskip.1em\allowbreak{} 时除以\allowbreak{}\hskip.1em \RB{4.100}{7.329}{0.620}{%
\RG{0.000}{0.000}{10.9091}{mi}{116}%
}\hskip.1em\allowbreak{} ，再令\allowbreak{}\hskip.1em \RB{20.973}{8.071}{2.605}{%
\RG{0.000}{0.000}{10.9091}{mi}{116}\RG{6.970}{0.000}{10.9091}{sy}{35}\RG{15.455}{0.000}{10.9091}{rm}{48}%
}\hskip.1em\allowbreak{} ，得到

\[
\RB{73.659}{8.677}{3.232}{%
\RG{0.000}{-0.005}{10.9091}{rm}{40}\RG{4.242}{-0.005}{10.9091}{mi}{11}\RG{11.246}{-0.005}{10.9091}{mi}{59}\RG{16.089}{-0.005}{10.9091}{mi}{11}\RG{28.894}{-0.005}{10.9091}{sy}{0}\RG{39.803}{-0.005}{10.9091}{mi}{11}\RG{46.822}{-0.005}{10.9091}{rm}{41}\RG{54.094}{-0.005}{10.9091}{sy}{21}\RG{65.610}{-0.005}{10.9091}{rm}{48}\RG{71.065}{-0.005}{10.9091}{mi}{58}\RG{23.088}{-1.641}{7.9701}{mi}{105}%
}
\]
因此

\[
\RB{168.819}{10.304}{3.228}{%
\RG{63.187}{4.504}{7.9701}{rm}{50}\RG{0.000}{-0.001}{10.9091}{rm}{40}\RG{4.243}{-0.001}{10.9091}{mi}{11}\RG{11.247}{-0.001}{10.9091}{mi}{59}\RG{16.090}{-0.001}{10.9091}{mi}{11}\RG{26.470}{-0.001}{10.9091}{rm}{41}\RG{33.743}{-0.001}{10.9091}{sy}{21}\RG{45.242}{-0.001}{10.9091}{sy}{107}\RG{50.713}{-0.001}{10.9091}{mi}{11}\RG{57.731}{-0.001}{10.9091}{sy}{107}\RG{70.950}{-0.001}{10.9091}{mi}{62}\RG{82.465}{-0.001}{10.9091}{rm}{48}\RG{109.727}{-0.001}{10.9091}{rm}{40}\RG{113.959}{-0.001}{10.9091}{rm}{49}\RG{122.465}{-0.001}{10.9091}{sy}{20}\RG{133.980}{-0.001}{10.9091}{mi}{105}\RG{140.768}{-0.001}{10.9091}{sy}{20}\RG{152.284}{-0.001}{10.9091}{mi}{109}\RG{161.862}{-0.001}{10.9091}{rm}{41}\RG{166.104}{-0.001}{10.9091}{mi}{59}\RG{23.089}{-1.637}{7.9701}{mi}{105}%
}
\]
即得所证结论。\allowbreak{}\hskip.1em \RB{8.366}{7.994}{0.501}{%
\RG{0.000}{-0.001}{10.9091}{am}{3}%
}\hskip.1em\allowbreak{}

小结：在不含原点的有限点凸包中，取距离原点最近的点；沿凸包内线段比较范数，即可得到严格为正的内积。

题目来源：谢启鸿《复旦大学高等代数习题集（2023版）》，第76页。

% Original page 2.
\section*{2019A01}
请用行列式的基本方法，求下列\allowbreak{}\hskip.1em \RB{6.729}{5.322}{0.620}{%
\RG{0.000}{-0.000}{10.9091}{mi}{110}%
}\hskip.1em\allowbreak{} 阶矩阵的行列式：

\[
\RB{322.665}{24.487}{15.764}{%
\RG{240.886}{23.998}{10.9092}{ex}{56}\RG{240.885}{14.179}{10.9092}{ex}{60}\RG{250.582}{9.824}{10.9091}{rm}{48}\RG{256.037}{9.824}{10.9091}{mi}{59}\RG{297.151}{9.824}{10.9091}{mi}{105}\RG{303.941}{9.824}{10.9091}{rm}{61}\RG{315.456}{9.824}{10.9091}{mi}{106}\RG{319.950}{9.824}{10.9091}{mi}{59}\RG{0.000}{1.634}{10.9091}{mi}{65}\RG{11.212}{1.634}{10.9091}{rm}{61}\RG{22.710}{1.634}{10.9091}{rm}{40}\RG{26.969}{1.634}{10.9091}{mi}{97}\RG{38.541}{1.634}{10.9091}{rm}{43}\RG{49.450}{1.634}{10.9091}{mi}{97}\RG{59.600}{1.634}{10.9091}{rm}{41}\RG{96.006}{1.634}{10.9091}{mi}{59}\RG{122.657}{1.634}{10.9091}{mi}{66}\RG{134.524}{1.634}{10.9091}{rm}{61}\RG{146.022}{1.634}{10.9091}{rm}{40}\RG{150.282}{1.634}{10.9091}{mi}{98}\RG{162.229}{1.634}{10.9091}{rm}{41}\RG{198.635}{1.634}{10.9091}{mi}{59}\RG{214.377}{1.634}{10.9091}{mi}{98}\RG{229.370}{1.634}{10.9091}{rm}{61}\RG{32.736}{-0.002}{7.9701}{mi}{105}\RG{55.217}{-0.002}{7.9701}{mi}{106}\RG{63.842}{-0.002}{7.9701}{rm}{49}\RG{68.077}{-0.002}{7.9701}{sy}{20}\RG{74.663}{-0.002}{7.9701}{mi}{105}\RG{77.540}{-0.002}{7.9701}{mi}{59}\RG{79.891}{-0.002}{7.9701}{mi}{106}\RG{83.783}{-0.002}{7.9701}{sy}{20}\RG{90.369}{-0.002}{7.9701}{mi}{110}\RG{154.964}{-0.002}{7.9701}{mi}{105}\RG{157.842}{-0.002}{7.9701}{mi}{106}\RG{166.471}{-0.002}{7.9701}{rm}{49}\RG{170.706}{-0.002}{7.9701}{sy}{20}\RG{177.292}{-0.002}{7.9701}{mi}{105}\RG{180.170}{-0.002}{7.9701}{mi}{59}\RG{182.521}{-0.002}{7.9701}{mi}{106}\RG{186.413}{-0.002}{7.9701}{sy}{20}\RG{192.999}{-0.002}{7.9701}{mi}{110}\RG{219.075}{-0.002}{7.9701}{mi}{105}\RG{221.952}{-0.002}{7.9701}{mi}{106}\RG{240.886}{-5.457}{10.9092}{ex}{58}\RG{250.582}{-8.386}{10.9091}{mi}{97}\RG{262.154}{-8.386}{10.9091}{rm}{43}\RG{273.063}{-8.386}{10.9091}{mi}{97}\RG{283.212}{-8.386}{10.9091}{mi}{59}\RG{297.143}{-8.386}{10.9091}{mi}{105}\RG{303.941}{-8.386}{10.9091}{sy}{54}\RG{303.941}{-8.386}{10.9091}{rm}{61}\RG{315.456}{-8.386}{10.9091}{mi}{106}\RG{319.950}{-8.386}{10.9091}{mi}{58}\RG{256.349}{-10.022}{7.9701}{mi}{105}\RG{278.830}{-10.022}{7.9701}{mi}{106}%
}
\]
解：记\allowbreak{}\hskip.1em \RB{81.395}{9.937}{3.225}{%
\RG{0.000}{0.002}{10.9091}{mi}{117}\RG{9.275}{0.002}{10.9091}{rm}{61}\RG{20.773}{0.002}{10.9091}{rm}{40}\RG{25.033}{0.002}{10.9091}{mi}{97}\RG{30.799}{-1.634}{7.9701}{rm}{49}\RG{35.531}{0.002}{10.9091}{mi}{59}\RG{40.364}{0.002}{10.9091}{mi}{58}\RG{45.208}{0.002}{10.9091}{mi}{58}\RG{50.051}{0.002}{10.9091}{mi}{58}\RG{54.884}{0.002}{10.9091}{mi}{59}\RG{59.728}{0.002}{10.9091}{mi}{97}\RG{65.541}{-1.634}{7.9701}{mi}{110}\RG{71.177}{0.002}{10.9091}{rm}{41}\RG{75.420}{3.961}{7.9701}{rm}{84}%
}\hskip.1em\allowbreak{} ，\allowbreak{}\hskip.1em \RB{69.237}{9.937}{3.225}{%
\RG{0.000}{0.002}{10.9091}{mi}{101}\RG{8.110}{0.002}{10.9091}{rm}{61}\RG{19.608}{0.002}{10.9091}{rm}{40}\RG{23.841}{0.002}{10.9091}{rm}{49}\RG{29.322}{0.002}{10.9091}{mi}{59}\RG{34.155}{0.002}{10.9091}{mi}{58}\RG{38.999}{0.002}{10.9091}{mi}{58}\RG{43.831}{0.002}{10.9091}{mi}{58}\RG{48.675}{0.002}{10.9091}{mi}{59}\RG{53.564}{0.002}{10.9091}{rm}{49}\RG{59.019}{0.002}{10.9091}{rm}{41}\RG{63.262}{3.961}{7.9701}{rm}{84}%
}\hskip.1em\allowbreak{} 。

(1)\allowbreak{}\hskip.1em \RB{8.366}{8.307}{0.504}{%
\RG{0.000}{-0.004}{10.9091}{mi}{65}%
}\hskip.1em\allowbreak{} 的第\allowbreak{}\hskip.1em \RB{4.831}{7.707}{2.740}{%
\RG{0.000}{-0.004}{10.9091}{mi}{106}%
}\hskip.1em\allowbreak{} 列为\allowbreak{}\hskip.1em \RB{35.053}{6.856}{3.767}{%
\RG{0.000}{-0.004}{10.9091}{mi}{117}\RG{8.736}{-0.004}{10.9091}{rm}{43}\RG{19.713}{-0.004}{10.9091}{mi}{97}\RG{25.479}{-1.641}{7.9701}{mi}{106}\RG{29.862}{-0.004}{10.9091}{mi}{101}%
}\hskip.1em\allowbreak{} 。当\allowbreak{}\hskip.1em \RB{26.916}{7.762}{1.999}{%
\RG{0.000}{-0.004}{10.9091}{mi}{110}\RG{9.746}{-0.004}{10.9091}{sy}{21}\RG{21.430}{-0.004}{10.9091}{rm}{51}%
}\hskip.1em\allowbreak{} 时，将各列按行列式的多线性展开，每一项中至少有两列同为\allowbreak{}\hskip.1em \RB{6.424}{5.323}{0.619}{%
\RG{0.000}{0.001}{10.9091}{mi}{117}%
}\hskip.1em\allowbreak{} 的倍数或同为\allowbreak{}\hskip.1em \RB{5.191}{5.323}{0.619}{%
\RG{0.000}{0.001}{10.9091}{mi}{101}%
}\hskip.1em\allowbreak{} 的倍数，故各项均为零。结合\allowbreak{}\hskip.1em \RB{36.794}{7.767}{2.604}{%
\RG{0.000}{0.001}{10.9091}{mi}{110}\RG{9.578}{0.001}{10.9091}{rm}{61}\RG{21.076}{0.001}{10.9091}{rm}{49}\RG{26.548}{0.001}{10.9091}{mi}{59}\RG{31.396}{0.001}{10.9091}{rm}{50}%
}\hskip.1em\allowbreak{} 时的直接计算，得到

\[
\RB{148.630}{30.309}{23.033}{%
\RG{39.697}{29.820}{10.9092}{ex}{56}\RG{39.697}{20.001}{10.9092}{ex}{62}\RG{49.394}{18.206}{10.9091}{rm}{50}\RG{54.848}{18.206}{10.9091}{mi}{97}\RG{65.348}{18.206}{10.9091}{mi}{59}\RG{119.367}{18.206}{10.9091}{mi}{110}\RG{128.945}{18.206}{10.9091}{rm}{61}\RG{140.443}{18.206}{10.9091}{rm}{49}\RG{145.915}{18.206}{10.9091}{mi}{59}\RG{39.697}{16.729}{10.9092}{ex}{62}\RG{60.616}{16.569}{7.9701}{rm}{49}\RG{39.697}{13.456}{10.9092}{ex}{60}\RG{100.695}{3.955}{7.9701}{rm}{50}\RG{0.000}{0.910}{10.9091}{rm}{100}\RG{6.055}{0.910}{10.9091}{rm}{101}\RG{10.898}{0.910}{10.9091}{rm}{116}\RG{16.970}{0.910}{10.9091}{mi}{65}\RG{28.182}{0.910}{10.9091}{rm}{61}\RG{49.394}{-0.004}{10.9091}{sy}{0}\RG{57.880}{-0.004}{10.9091}{rm}{40}\RG{62.122}{-0.004}{10.9091}{mi}{97}\RG{75.045}{-0.004}{10.9091}{sy}{0}\RG{85.954}{-0.004}{10.9091}{mi}{97}\RG{96.453}{-0.004}{10.9091}{rm}{41}\RG{105.428}{-0.004}{10.9091}{mi}{59}\RG{119.348}{-0.004}{10.9091}{mi}{110}\RG{128.945}{-0.004}{10.9091}{rm}{61}\RG{140.443}{-0.004}{10.9091}{rm}{50}\RG{145.915}{-0.004}{10.9091}{mi}{59}\RG{67.889}{-1.641}{7.9701}{rm}{49}\RG{91.721}{-1.641}{7.9701}{rm}{50}\RG{39.697}{-6.181}{10.9092}{ex}{62}\RG{39.697}{-9.453}{10.9092}{ex}{62}\RG{39.697}{-12.726}{10.9092}{ex}{58}\RG{49.394}{-18.214}{10.9091}{rm}{48}\RG{54.848}{-18.214}{10.9091}{mi}{59}\RG{119.367}{-18.214}{10.9091}{mi}{110}\RG{128.945}{-18.214}{10.9091}{sy}{21}\RG{140.460}{-18.214}{10.9091}{rm}{51}\RG{145.915}{-18.214}{10.9091}{mi}{58}%
}
\]
(2)当\allowbreak{}\hskip.1em \RB{26.148}{7.766}{0.620}{%
\RG{0.000}{0.000}{10.9091}{mi}{110}\RG{9.579}{0.000}{10.9091}{rm}{61}\RG{21.077}{0.000}{10.9091}{rm}{49}%
}\hskip.1em\allowbreak{} 时，\allowbreak{}\hskip.1em \RB{45.839}{8.071}{0.740}{%
\RG{0.000}{0.000}{10.9091}{rm}{100}\RG{6.055}{0.000}{10.9091}{rm}{101}\RG{10.898}{0.000}{10.9091}{rm}{116}\RG{16.970}{0.000}{10.9091}{mi}{66}\RG{28.822}{0.000}{10.9091}{rm}{61}\RG{40.321}{0.000}{10.9091}{rm}{48}%
}\hskip.1em\allowbreak{} 。以下设\allowbreak{}\hskip.1em \RB{26.491}{7.766}{1.995}{%
\RG{0.000}{0.000}{10.9091}{mi}{110}\RG{9.578}{0.000}{10.9091}{sy}{21}\RG{21.093}{0.000}{10.9091}{rm}{50}%
}\hskip.1em\allowbreak{} ，并约定空乘积为\allowbreak{}\hskip.1em \RB{5.071}{7.766}{0.500}{%
\RG{0.000}{0.000}{10.9091}{rm}{49}%
}\hskip.1em\allowbreak{} 。写成

\[
\RB{284.656}{17.648}{13.671}{%
\RG{261.557}{13.633}{7.9701}{mi}{110}\RG{257.156}{10.360}{10.9092}{ex}{89}\RG{0.000}{-0.003}{10.9091}{mi}{66}\RG{11.852}{-0.003}{10.9091}{rm}{61}\RG{23.367}{-0.003}{10.9091}{mi}{68}\RG{35.126}{-0.003}{10.9091}{rm}{43}\RG{46.035}{-0.003}{10.9091}{mi}{65}\RG{54.217}{-0.003}{10.9091}{mi}{59}\RG{80.868}{-0.003}{10.9091}{mi}{68}\RG{93.250}{-0.003}{10.9091}{rm}{61}\RG{104.748}{-0.003}{10.9091}{rm}{100}\RG{110.803}{-0.003}{10.9091}{rm}{105}\RG{113.825}{-0.003}{10.9091}{rm}{97}\RG{119.279}{-0.003}{10.9091}{rm}{103}\RG{124.886}{-0.003}{10.9091}{rm}{40}\RG{129.159}{-0.003}{10.9091}{sy}{0}\RG{137.644}{-0.003}{10.9091}{rm}{50}\RG{143.099}{-0.003}{10.9091}{mi}{97}\RG{153.597}{-0.003}{10.9091}{mi}{59}\RG{158.430}{-0.003}{10.9091}{mi}{58}\RG{163.274}{-0.003}{10.9091}{mi}{58}\RG{168.117}{-0.003}{10.9091}{mi}{58}\RG{172.950}{-0.003}{10.9091}{mi}{59}\RG{177.840}{-0.003}{10.9091}{sy}{0}\RG{186.325}{-0.003}{10.9091}{rm}{50}\RG{191.780}{-0.003}{10.9091}{mi}{97}\RG{203.183}{-0.003}{10.9091}{rm}{41}\RG{207.425}{-0.003}{10.9091}{mi}{59}\RG{234.076}{-0.003}{10.9091}{mi}{80}\RG{245.641}{-0.003}{10.9091}{rm}{61}\RG{272.914}{-0.003}{10.9091}{mi}{97}\RG{282.061}{-0.003}{10.9091}{mi}{58}\RG{148.865}{-1.640}{7.9701}{rm}{49}\RG{197.546}{-1.640}{7.9701}{mi}{110}\RG{278.680}{-1.640}{7.9701}{mi}{105}\RG{257.274}{-13.091}{7.9701}{mi}{105}\RG{260.157}{-13.091}{7.9701}{rm}{61}\RG{266.741}{-13.091}{7.9701}{rm}{49}%
}
\]
按列的多线性展开\allowbreak{}\hskip.1em \RB{54.209}{8.686}{3.223}{%
\RG{0.000}{0.004}{10.9091}{rm}{100}\RG{6.055}{0.004}{10.9091}{rm}{101}\RG{10.898}{0.004}{10.9091}{rm}{116}\RG{15.131}{0.004}{10.9091}{rm}{40}\RG{19.394}{0.004}{10.9091}{mi}{68}\RG{31.309}{0.004}{10.9091}{rm}{43}\RG{42.374}{0.004}{10.9091}{mi}{65}\RG{50.556}{0.004}{10.9091}{rm}{41}%
}\hskip.1em\allowbreak{} 。若有至少三列选自\allowbreak{}\hskip.1em \RB{8.366}{8.315}{0.500}{%
\RG{0.000}{0.004}{10.9091}{mi}{65}%
}\hskip.1em\allowbreak{} ，由(1)的同一理由，该项为零。只需计算以下三类项。

没有列选自\allowbreak{}\hskip.1em \RB{8.366}{8.314}{0.500}{%
\RG{0.000}{0.003}{10.9091}{mi}{65}%
}\hskip.1em\allowbreak{} 时，所得项为\allowbreak{}\hskip.1em \RB{36.787}{8.685}{3.224}{%
\RG{0.000}{0.003}{10.9091}{rm}{40}\RG{4.243}{0.003}{10.9091}{sy}{0}\RG{12.727}{0.003}{10.9091}{rm}{50}\RG{18.182}{0.003}{10.9091}{rm}{41}\RG{22.424}{3.961}{7.9701}{mi}{110}\RG{28.061}{0.003}{10.9091}{mi}{80}%
}\hskip.1em\allowbreak{} 。恰有第\allowbreak{}\hskip.1em \RB{3.696}{7.707}{0.624}{%
\RG{0.000}{-0.004}{10.9091}{mi}{105}%
}\hskip.1em\allowbreak{} 列选自\allowbreak{}\hskip.1em \RB{8.366}{8.307}{0.504}{%
\RG{0.000}{-0.004}{10.9091}{mi}{65}%
}\hskip.1em\allowbreak{} 时，沿其余对角列展开，所得项为

\[
\RB{124.807}{10.864}{15.567}{%
\RG{16.496}{10.364}{10.9092}{ex}{89}\RG{109.278}{4.505}{7.9701}{mi}{110}\RG{0.000}{0.000}{10.9091}{rm}{50}\RG{5.454}{0.000}{10.9091}{mi}{97}\RG{30.511}{0.000}{10.9091}{rm}{40}\RG{34.754}{0.000}{10.9091}{sy}{0}\RG{43.239}{0.000}{10.9091}{rm}{50}\RG{48.693}{0.000}{10.9091}{mi}{97}\RG{59.580}{0.000}{10.9091}{rm}{41}\RG{66.835}{0.000}{10.9091}{rm}{61}\RG{78.368}{0.000}{10.9091}{sy}{0}\RG{86.853}{0.000}{10.9091}{rm}{40}\RG{91.096}{0.000}{10.9091}{sy}{0}\RG{99.581}{0.000}{10.9091}{rm}{50}\RG{105.035}{0.000}{10.9091}{rm}{41}\RG{114.914}{0.000}{10.9091}{mi}{80}\RG{122.212}{0.000}{10.9091}{mi}{58}\RG{11.221}{-1.636}{7.9701}{mi}{105}\RG{54.460}{-1.777}{7.9701}{mi}{107}\RG{16.420}{-13.353}{7.9701}{mi}{107}\RG{21.042}{-13.353}{7.9701}{sy}{54}\RG{21.042}{-13.353}{7.9701}{rm}{61}\RG{27.628}{-13.353}{7.9701}{mi}{105}%
}
\]
此类共\allowbreak{}\hskip.1em \RB{6.729}{5.323}{0.619}{%
\RG{0.000}{0.001}{10.9091}{mi}{110}%
}\hskip.1em\allowbreak{} 项。

恰有第\allowbreak{}\hskip.1em \RB{13.416}{7.715}{2.732}{%
\RG{0.000}{0.004}{10.9091}{mi}{105}\RG{3.753}{0.004}{10.9091}{mi}{59}\RG{8.586}{0.004}{10.9091}{mi}{106}%
}\hskip.1em\allowbreak{} 列选自\allowbreak{}\hskip.1em \RB{8.366}{8.315}{0.500}{%
\RG{0.000}{0.004}{10.9091}{mi}{65}%
}\hskip.1em\allowbreak{} （\allowbreak{}\hskip.1em \RB{23.115}{7.715}{2.732}{%
\RG{0.000}{0.004}{10.9091}{mi}{105}\RG{6.775}{0.004}{10.9091}{mi}{60}\RG{18.284}{0.004}{10.9091}{mi}{106}%
}\hskip.1em\allowbreak{} ）时，沿其余对角列展开，所得项为

\[
\RB{283.156}{19.824}{15.563}{%
\RG{0.000}{19.095}{10.9092}{ex}{12}\RG{78.859}{19.095}{10.9092}{ex}{12}\RG{0.000}{12.550}{10.9092}{ex}{12}\RG{78.859}{12.550}{10.9092}{ex}{12}\RG{87.508}{10.368}{10.9092}{ex}{89}\RG{250.723}{10.368}{10.9092}{ex}{89}\RG{12.650}{7.284}{10.9091}{rm}{50}\RG{18.105}{7.284}{10.9091}{mi}{97}\RG{46.229}{7.284}{10.9091}{mi}{97}\RG{57.801}{7.284}{10.9091}{rm}{43}\RG{68.711}{7.284}{10.9091}{mi}{97}\RG{0.000}{6.004}{10.9092}{ex}{12}\RG{78.859}{6.004}{10.9092}{ex}{12}\RG{23.871}{5.648}{7.9701}{mi}{105}\RG{51.995}{5.648}{7.9701}{mi}{105}\RG{74.477}{5.648}{7.9701}{mi}{106}\RG{183.407}{4.509}{7.9701}{mi}{110}\RG{188.546}{4.509}{7.9701}{sy}{0}\RG{195.132}{4.509}{7.9701}{rm}{50}\RG{240.979}{4.509}{7.9701}{rm}{50}\RG{104.641}{0.004}{10.9091}{rm}{40}\RG{108.884}{0.004}{10.9091}{sy}{0}\RG{117.369}{0.004}{10.9091}{rm}{50}\RG{122.823}{0.004}{10.9091}{mi}{97}\RG{133.709}{0.004}{10.9091}{rm}{41}\RG{140.964}{0.004}{10.9091}{rm}{61}\RG{152.497}{0.004}{10.9091}{sy}{0}\RG{160.982}{0.004}{10.9091}{rm}{40}\RG{165.225}{0.004}{10.9091}{sy}{0}\RG{173.709}{0.004}{10.9091}{rm}{50}\RG{179.164}{0.004}{10.9091}{rm}{41}\RG{199.864}{0.004}{10.9091}{rm}{40}\RG{204.107}{0.004}{10.9091}{mi}{97}\RG{215.679}{0.004}{10.9091}{sy}{0}\RG{226.588}{0.004}{10.9091}{mi}{97}\RG{236.736}{0.004}{10.9091}{rm}{41}\RG{269.675}{0.004}{10.9091}{mi}{97}\RG{280.561}{0.004}{10.9091}{mi}{58}\RG{0.000}{-0.541}{10.9092}{ex}{12}\RG{78.859}{-0.541}{10.9092}{ex}{12}\RG{209.873}{-1.632}{7.9701}{mi}{105}\RG{232.354}{-1.632}{7.9701}{mi}{106}\RG{128.590}{-1.773}{7.9701}{mi}{107}\RG{275.442}{-1.773}{7.9701}{mi}{107}\RG{0.000}{-7.087}{10.9092}{ex}{12}\RG{78.859}{-7.087}{10.9092}{ex}{12}\RG{3.636}{-7.891}{10.9091}{mi}{97}\RG{15.208}{-7.891}{10.9091}{rm}{43}\RG{26.117}{-7.891}{10.9091}{mi}{97}\RG{54.742}{-7.891}{10.9091}{rm}{50}\RG{60.197}{-7.891}{10.9091}{mi}{97}\RG{9.403}{-9.527}{7.9701}{mi}{105}\RG{31.884}{-9.527}{7.9701}{mi}{106}\RG{65.964}{-9.527}{7.9701}{mi}{106}\RG{84.314}{-13.349}{7.9701}{mi}{107}\RG{88.935}{-13.349}{7.9701}{sy}{54}\RG{88.935}{-13.349}{7.9701}{rm}{61}\RG{95.522}{-13.349}{7.9701}{mi}{105}\RG{98.399}{-13.349}{7.9701}{mi}{59}\RG{100.750}{-13.349}{7.9701}{mi}{106}\RG{247.529}{-13.349}{7.9701}{mi}{107}\RG{252.151}{-13.349}{7.9701}{sy}{54}\RG{252.151}{-13.349}{7.9701}{rm}{61}\RG{258.737}{-13.349}{7.9701}{mi}{105}\RG{261.614}{-13.349}{7.9701}{mi}{59}\RG{263.966}{-13.349}{7.9701}{mi}{106}%
}
\]
\noindent \RB{365.257}{34.416}{17.403}{%
\RT{0.000}{25.188}{10.9091}{于}\RT{10.909}{25.188}{10.9091}{是}\RG{180.230}{21.923}{10.9092}{ex}{50}\RG{353.571}{21.923}{10.9092}{ex}{51}\RG{257.116}{10.359}{10.9092}{ex}{88}\RG{323.734}{10.359}{10.9092}{ex}{89}\RG{161.955}{4.500}{7.9701}{mi}{110}\RG{167.093}{4.500}{7.9701}{sy}{0}\RG{173.679}{4.500}{7.9701}{rm}{50}\RG{313.989}{4.500}{7.9701}{rm}{50}\RG{180.230}{2.286}{10.9092}{ex}{52}\RG{353.571}{2.286}{10.9092}{ex}{53}\RG{99.193}{-0.004}{10.9091}{rm}{100}\RG{105.247}{-0.004}{10.9091}{rm}{101}\RG{110.091}{-0.004}{10.9091}{rm}{116}\RG{116.163}{-0.004}{10.9091}{mi}{66}\RG{128.015}{-0.004}{10.9091}{rm}{61}\RG{139.513}{-0.004}{10.9091}{rm}{40}\RG{143.773}{-0.004}{10.9091}{sy}{0}\RG{152.258}{-0.004}{10.9091}{rm}{50}\RG{157.713}{-0.004}{10.9091}{rm}{41}\RG{187.504}{-0.004}{10.9091}{sy}{0}\RG{195.988}{-0.004}{10.9091}{rm}{52}\RG{201.442}{-0.004}{10.9091}{rm}{40}\RG{205.685}{-0.004}{10.9091}{mi}{110}\RG{214.658}{-0.004}{10.9091}{sy}{0}\RG{225.567}{-0.004}{10.9091}{rm}{49}\RG{231.021}{-0.004}{10.9091}{rm}{41}\RG{235.264}{-0.004}{10.9091}{mi}{80}\RG{246.207}{-0.004}{10.9091}{sy}{0}\RG{272.874}{-0.004}{10.9091}{rm}{40}\RG{277.116}{-0.004}{10.9091}{mi}{97}\RG{288.688}{-0.004}{10.9091}{sy}{0}\RG{299.597}{-0.004}{10.9091}{mi}{97}\RG{309.747}{-0.004}{10.9091}{rm}{41}\RG{342.685}{-0.004}{10.9091}{mi}{97}\RG{362.662}{-0.004}{10.9091}{mi}{58}\RG{282.882}{-1.641}{7.9701}{mi}{105}\RG{305.363}{-1.641}{7.9701}{mi}{106}\RG{348.452}{-1.782}{7.9701}{mi}{107}\RG{258.318}{-13.092}{7.9701}{mi}{105}\RG{261.195}{-13.092}{7.9701}{mi}{60}\RG{267.779}{-13.092}{7.9701}{mi}{106}\RG{320.539}{-13.358}{7.9701}{mi}{107}\RG{325.161}{-13.358}{7.9701}{sy}{54}\RG{325.161}{-13.358}{7.9701}{rm}{61}\RG{331.748}{-13.358}{7.9701}{mi}{105}\RG{334.625}{-13.358}{7.9701}{mi}{59}\RG{336.976}{-13.358}{7.9701}{mi}{106}%
}

也可将其写为

\[
\RB{287.884}{22.851}{17.401}{%
\RG{81.038}{21.925}{10.9092}{ex}{50}\RG{212.736}{21.925}{10.9092}{ex}{48}\RG{266.653}{21.925}{10.9092}{ex}{49}\RG{276.198}{21.925}{10.9092}{ex}{51}\RG{166.920}{18.652}{10.9092}{ex}{32}\RG{202.281}{18.652}{10.9092}{ex}{33}\RG{133.083}{13.634}{7.9701}{mi}{110}\RG{180.868}{13.634}{7.9701}{mi}{110}\RG{227.591}{13.634}{7.9701}{mi}{110}\RG{128.682}{10.361}{10.9092}{ex}{89}\RG{175.557}{10.361}{10.9092}{ex}{88}\RG{222.281}{10.361}{10.9092}{ex}{88}\RG{239.933}{10.361}{10.9092}{ex}{89}\RG{62.763}{4.502}{7.9701}{mi}{110}\RG{67.901}{4.502}{7.9701}{sy}{0}\RG{74.488}{4.502}{7.9701}{rm}{50}\RG{122.131}{4.502}{7.9701}{rm}{50}\RG{81.038}{2.288}{10.9092}{ex}{52}\RG{212.736}{2.288}{10.9092}{ex}{64}\RG{266.653}{2.288}{10.9092}{ex}{65}\RG{276.198}{2.288}{10.9092}{ex}{53}\RG{0.000}{-0.003}{10.9091}{rm}{100}\RG{6.055}{-0.003}{10.9091}{rm}{101}\RG{10.898}{-0.003}{10.9091}{rm}{116}\RG{16.970}{-0.003}{10.9091}{mi}{66}\RG{28.822}{-0.003}{10.9091}{rm}{61}\RG{40.321}{-0.003}{10.9091}{rm}{40}\RG{44.580}{-0.003}{10.9091}{sy}{0}\RG{53.065}{-0.003}{10.9091}{rm}{50}\RG{58.520}{-0.003}{10.9091}{rm}{41}\RG{88.311}{-0.003}{10.9091}{rm}{40}\RG{92.553}{-0.003}{10.9091}{mi}{110}\RG{101.525}{-0.003}{10.9091}{sy}{0}\RG{112.434}{-0.003}{10.9091}{rm}{50}\RG{117.889}{-0.003}{10.9091}{rm}{41}\RG{144.440}{-0.003}{10.9091}{mi}{97}\RG{156.011}{-0.003}{10.9091}{sy}{0}\RG{193.134}{-0.003}{10.9091}{mi}{97}\RG{255.767}{-0.003}{10.9091}{mi}{97}\RG{285.289}{-0.003}{10.9091}{mi}{58}\RG{150.206}{-1.639}{7.9701}{mi}{105}\RG{198.900}{-1.639}{7.9701}{mi}{105}\RG{261.533}{-1.780}{7.9701}{mi}{107}\RG{128.800}{-13.090}{7.9701}{mi}{105}\RG{131.683}{-13.090}{7.9701}{rm}{61}\RG{138.266}{-13.090}{7.9701}{rm}{49}\RG{176.584}{-13.090}{7.9701}{mi}{105}\RG{179.467}{-13.090}{7.9701}{rm}{61}\RG{186.051}{-13.090}{7.9701}{rm}{49}\RG{223.308}{-13.090}{7.9701}{mi}{105}\RG{226.191}{-13.090}{7.9701}{rm}{61}\RG{232.774}{-13.090}{7.9701}{rm}{49}\RG{239.857}{-13.356}{7.9701}{mi}{107}\RG{244.479}{-13.356}{7.9701}{sy}{54}\RG{244.479}{-13.356}{7.9701}{rm}{61}\RG{251.065}{-13.356}{7.9701}{mi}{105}%
}
\]
两式等价，因为

\[
\RB{254.785}{22.857}{17.394}{%
\RG{156.815}{21.931}{10.9092}{ex}{48}\RG{210.733}{21.931}{10.9092}{ex}{49}\RG{111.001}{18.659}{10.9092}{ex}{32}\RG{146.361}{18.659}{10.9092}{ex}{33}\RG{0.000}{10.368}{10.9092}{ex}{88}\RG{66.618}{10.368}{10.9092}{ex}{89}\RG{119.637}{10.368}{10.9092}{ex}{88}\RG{166.361}{10.368}{10.9092}{ex}{88}\RG{184.014}{10.368}{10.9092}{ex}{89}\RG{56.872}{4.508}{7.9701}{rm}{50}\RG{240.160}{4.508}{7.9701}{rm}{50}\RG{156.815}{2.295}{10.9092}{ex}{64}\RG{210.733}{2.295}{10.9092}{ex}{65}\RG{15.758}{0.004}{10.9091}{rm}{40}\RG{20.000}{0.004}{10.9091}{mi}{97}\RG{31.572}{0.004}{10.9091}{sy}{0}\RG{42.481}{0.004}{10.9091}{mi}{97}\RG{52.630}{0.004}{10.9091}{rm}{41}\RG{85.570}{0.004}{10.9091}{mi}{97}\RG{99.486}{0.004}{10.9091}{rm}{61}\RG{137.213}{0.004}{10.9091}{mi}{97}\RG{199.847}{0.004}{10.9091}{mi}{97}\RG{222.703}{0.004}{10.9091}{sy}{0}\RG{233.612}{0.004}{10.9091}{mi}{110}\RG{244.892}{0.004}{10.9091}{mi}{80}\RG{252.190}{0.004}{10.9091}{mi}{58}\RG{25.767}{-1.632}{7.9701}{mi}{105}\RG{48.248}{-1.632}{7.9701}{mi}{106}\RG{142.980}{-1.632}{7.9701}{mi}{105}\RG{91.336}{-1.773}{7.9701}{mi}{107}\RG{205.614}{-1.773}{7.9701}{mi}{107}\RG{1.202}{-13.083}{7.9701}{mi}{105}\RG{4.079}{-13.083}{7.9701}{mi}{60}\RG{10.663}{-13.083}{7.9701}{mi}{106}\RG{126.075}{-13.083}{7.9701}{mi}{105}\RG{172.798}{-13.083}{7.9701}{mi}{105}\RG{63.424}{-13.349}{7.9701}{mi}{107}\RG{68.045}{-13.349}{7.9701}{sy}{54}\RG{68.045}{-13.349}{7.9701}{rm}{61}\RG{74.632}{-13.349}{7.9701}{mi}{105}\RG{77.509}{-13.349}{7.9701}{mi}{59}\RG{79.860}{-13.349}{7.9701}{mi}{106}\RG{183.938}{-13.349}{7.9701}{mi}{107}\RG{188.559}{-13.349}{7.9701}{sy}{54}\RG{188.559}{-13.349}{7.9701}{rm}{61}\RG{195.146}{-13.349}{7.9701}{mi}{105}%
}
\]
小结：利用列向量只含两种共同成分，将多线性展开截断到至多两列。上述公式没有除以\allowbreak{}\hskip.1em \RB{8.793}{5.322}{2.217}{%
\RG{0.000}{0.000}{10.9091}{mi}{97}\RG{5.766}{-1.637}{7.9701}{mi}{105}%
}\hskip.1em\allowbreak{} ，故允许任意参数为零。

题目来源：谢启鸿《复旦大学高等代数习题集（2023版）》，第80页。

% Original page 3.
\section*{2019A02}
设\allowbreak{}\hskip.1em \RB{96.247}{11.104}{3.840}{%
\RG{0.000}{-0.003}{10.9091}{mi}{65}\RG{11.212}{-0.003}{10.9091}{rm}{61}\RG{22.710}{-0.003}{10.9091}{rm}{40}\RG{26.969}{-0.003}{10.9091}{mi}{120}\RG{33.204}{5.304}{7.9701}{mi}{106}\RG{33.204}{-3.260}{7.9701}{mi}{105}\RG{37.088}{5.304}{7.9701}{sy}{0}\RG{43.675}{5.304}{7.9701}{rm}{49}\RG{48.407}{-0.003}{10.9091}{rm}{41}\RG{52.649}{-1.639}{7.9701}{rm}{49}\RG{56.884}{-1.639}{7.9701}{sy}{20}\RG{63.471}{-1.639}{7.9701}{mi}{105}\RG{66.348}{-1.639}{7.9701}{mi}{59}\RG{68.699}{-1.639}{7.9701}{mi}{106}\RG{72.591}{-1.639}{7.9701}{sy}{20}\RG{79.177}{-1.639}{7.9701}{rm}{50}\RG{83.409}{-1.639}{7.9701}{rm}{48}\RG{87.641}{-1.639}{7.9701}{rm}{49}\RG{91.873}{-1.639}{7.9701}{rm}{57}%
}\hskip.1em\allowbreak{} 为\allowbreak{}\hskip.1em \RB{21.849}{7.763}{0.743}{%
\RG{0.000}{-0.003}{10.9091}{rm}{50}\RG{5.455}{-0.003}{10.9091}{rm}{48}\RG{10.909}{-0.003}{10.9091}{rm}{49}\RG{16.364}{-0.003}{10.9091}{rm}{57}%
}\hskip.1em\allowbreak{} 阶矩阵，\allowbreak{}\hskip.1em \RB{14.851}{8.308}{3.765}{%
\RG{0.000}{-0.003}{10.9091}{mi}{65}\RG{8.182}{-1.639}{7.9701}{mi}{105}\RG{11.059}{-1.639}{7.9701}{mi}{106}%
}\hskip.1em\allowbreak{} 为其第\allowbreak{}\hskip.1em \RB{21.619}{8.679}{3.230}{%
\RG{0.000}{-0.003}{10.9091}{rm}{40}\RG{4.242}{-0.003}{10.9091}{mi}{105}\RG{7.995}{-0.003}{10.9091}{mi}{59}\RG{12.828}{-0.003}{10.9091}{mi}{106}\RG{17.966}{-0.003}{10.9091}{rm}{41}%
}\hskip.1em\allowbreak{} 元的代数余子式。求

\[
\RB{97.554}{19.433}{15.217}{%
\RG{1.502}{13.633}{7.9701}{rm}{50}\RG{5.734}{13.633}{7.9701}{rm}{48}\RG{9.966}{13.633}{7.9701}{rm}{49}\RG{14.198}{13.633}{7.9701}{rm}{57}\RG{2.092}{10.360}{10.9092}{ex}{88}\RG{30.418}{4.501}{7.9701}{rm}{50}\RG{34.650}{4.501}{7.9701}{rm}{48}\RG{38.883}{4.501}{7.9701}{rm}{49}\RG{43.115}{4.501}{7.9701}{rm}{57}\RG{66.304}{4.501}{7.9701}{rm}{55}\RG{70.536}{4.501}{7.9701}{rm}{48}\RG{19.941}{-0.004}{10.9091}{rm}{40}\RG{24.183}{-0.004}{10.9091}{mi}{120}\RG{50.277}{-0.004}{10.9091}{rm}{43}\RG{61.186}{-0.004}{10.9091}{mi}{106}\RG{75.270}{-0.004}{10.9091}{rm}{41}\RG{79.513}{-0.004}{10.9091}{mi}{65}\RG{94.960}{-0.004}{10.9091}{mi}{58}\RG{87.695}{-1.640}{7.9701}{mi}{105}\RG{90.572}{-1.640}{7.9701}{mi}{106}\RG{30.418}{-2.701}{7.9701}{mi}{105}\RG{0.000}{-13.091}{7.9701}{mi}{105}\RG{2.877}{-13.091}{7.9701}{mi}{59}\RG{5.228}{-13.091}{7.9701}{mi}{106}\RG{9.120}{-13.091}{7.9701}{rm}{61}\RG{15.703}{-13.091}{7.9701}{rm}{49}%
}
\]
解：记\allowbreak{}\hskip.1em \RB{46.333}{7.951}{0.740}{%
\RG{0.000}{-0.000}{10.9091}{mi}{78}\RG{12.985}{-0.000}{10.9091}{rm}{61}\RG{24.483}{-0.000}{10.9091}{rm}{50}\RG{29.938}{-0.000}{10.9091}{rm}{48}\RG{35.393}{-0.000}{10.9091}{rm}{49}\RG{40.847}{-0.000}{10.9091}{rm}{57}%
}\hskip.1em\allowbreak{} ，\allowbreak{}\hskip.1em \RB{49.216}{8.311}{0.620}{%
\RG{0.000}{-0.000}{10.9091}{mi}{68}\RG{12.365}{-0.000}{10.9091}{rm}{61}\RG{23.863}{-0.000}{10.9091}{rm}{100}\RG{29.918}{-0.000}{10.9091}{rm}{101}\RG{34.762}{-0.000}{10.9091}{rm}{116}\RG{40.850}{-0.000}{10.9091}{mi}{65}%
}\hskip.1em\allowbreak{} ，\allowbreak{}\hskip.1em \RB{70.177}{11.108}{3.837}{%
\RG{0.000}{-0.000}{10.9091}{mi}{67}\RG{7.797}{-1.636}{7.9701}{mi}{106}\RG{15.209}{-0.000}{10.9091}{rm}{61}\RG{26.707}{-0.000}{10.9091}{rm}{40}\RG{30.967}{-0.000}{10.9091}{mi}{120}\RG{37.202}{5.308}{7.9701}{mi}{106}\RG{37.202}{-3.257}{7.9701}{mi}{105}\RG{41.086}{5.308}{7.9701}{sy}{0}\RG{47.672}{5.308}{7.9701}{rm}{49}\RG{52.405}{-0.000}{10.9091}{rm}{41}\RG{56.647}{3.959}{7.9701}{mi}{78}\RG{56.647}{-3.055}{7.9701}{mi}{105}\RG{59.530}{-3.055}{7.9701}{rm}{61}\RG{66.114}{-3.055}{7.9701}{rm}{49}%
}\hskip.1em\allowbreak{} 为\allowbreak{}\hskip.1em \RB{8.366}{8.311}{0.500}{%
\RG{0.000}{-0.000}{10.9091}{mi}{65}%
}\hskip.1em\allowbreak{} 的第\allowbreak{}\hskip.1em \RB{4.831}{7.711}{2.736}{%
\RG{0.000}{-0.000}{10.9091}{mi}{106}%
}\hskip.1em\allowbreak{} 列。由Vandermonde行列式公式，

\[
\RB{106.858}{10.865}{15.389}{%
\RG{36.077}{10.365}{10.9092}{ex}{89}\RG{0.000}{0.001}{10.9091}{mi}{68}\RG{12.366}{0.001}{10.9091}{rm}{61}\RG{62.213}{0.001}{10.9091}{rm}{40}\RG{66.455}{0.001}{10.9091}{mi}{120}\RG{79.496}{0.001}{10.9091}{sy}{0}\RG{90.406}{0.001}{10.9091}{mi}{120}\RG{100.022}{0.001}{10.9091}{rm}{41}\RG{104.264}{0.001}{10.9091}{mi}{58}\RG{72.690}{-1.635}{7.9701}{mi}{106}\RG{96.640}{-1.635}{7.9701}{mi}{105}\RG{23.881}{-13.263}{7.9701}{rm}{49}\RG{28.115}{-13.263}{7.9701}{sy}{20}\RG{34.701}{-13.263}{7.9701}{mi}{105}\RG{37.579}{-13.263}{7.9701}{mi}{60}\RG{44.162}{-13.263}{7.9701}{mi}{106}\RG{48.055}{-13.263}{7.9701}{sy}{20}\RG{54.641}{-13.263}{7.9701}{mi}{78}%
}
\]
这一公式可由归纳法得到：将行列式视为\allowbreak{}\hskip.1em \RB{14.036}{5.321}{2.190}{%
\RG{0.000}{-0.001}{10.9091}{mi}{120}\RG{6.235}{-1.690}{7.9701}{mi}{78}%
}\hskip.1em\allowbreak{} 的至多\allowbreak{}\hskip.1em \RB{28.176}{7.950}{0.501}{%
\RG{0.000}{-0.001}{10.9091}{mi}{78}\RG{12.288}{-0.001}{10.9091}{sy}{0}\RG{23.105}{-0.001}{10.9091}{rm}{49}%
}\hskip.1em\allowbreak{} 次多项式，代入\allowbreak{}\hskip.1em \RB{38.111}{5.321}{2.218}{%
\RG{0.000}{-0.001}{10.9091}{mi}{120}\RG{6.235}{-1.690}{7.9701}{mi}{78}\RG{17.334}{-0.001}{10.9091}{rm}{61}\RG{28.849}{-0.001}{10.9091}{mi}{120}\RG{35.084}{-1.638}{7.9701}{mi}{105}%
}\hskip.1em\allowbreak{} （\allowbreak{}\hskip.1em \RB{28.395}{7.950}{0.937}{%
\RG{0.000}{-0.001}{10.9091}{mi}{105}\RG{6.775}{-0.001}{10.9091}{mi}{60}\RG{18.284}{-0.001}{10.9091}{mi}{78}%
}\hskip.1em\allowbreak{} ）时有两行相同，而最高次项系数为前\allowbreak{}\hskip.1em \RB{28.359}{7.955}{0.500}{%
\RG{0.000}{0.004}{10.9091}{mi}{78}\RG{12.379}{0.004}{10.9091}{sy}{0}\RG{23.288}{0.004}{10.9091}{rm}{49}%
}\hskip.1em\allowbreak{} 阶Vandermonde行列式。

先计算含\allowbreak{}\hskip.1em \RB{13.747}{9.844}{2.739}{%
\RG{0.000}{-0.003}{10.9091}{mi}{106}\RG{5.117}{3.956}{7.9701}{rm}{55}\RG{9.349}{3.956}{7.9701}{rm}{48}%
}\hskip.1em\allowbreak{} 的部分。将第\allowbreak{}\hskip.1em \RB{4.831}{7.708}{2.739}{%
\RG{0.000}{-0.003}{10.9091}{mi}{106}%
}\hskip.1em\allowbreak{} 列替换为\allowbreak{}\hskip.1em \RB{76.687}{9.932}{3.230}{%
\RG{0.000}{-0.003}{10.9091}{mi}{67}\RG{7.797}{-1.640}{7.9701}{rm}{49}\RG{15.560}{-0.003}{10.9091}{rm}{61}\RG{27.058}{-0.003}{10.9091}{rm}{40}\RG{31.291}{-0.003}{10.9091}{rm}{49}\RG{36.772}{-0.003}{10.9091}{mi}{59}\RG{41.605}{-0.003}{10.9091}{mi}{58}\RG{46.449}{-0.003}{10.9091}{mi}{58}\RG{51.282}{-0.003}{10.9091}{mi}{58}\RG{56.125}{-0.003}{10.9091}{mi}{59}\RG{61.015}{-0.003}{10.9091}{rm}{49}\RG{66.469}{-0.003}{10.9091}{rm}{41}\RG{70.712}{3.956}{7.9701}{rm}{84}%
}\hskip.1em\allowbreak{} ，沿该列展开，得到

\[
\RB{108.371}{32.868}{7.382}{%
\RG{47.568}{32.379}{10.9092}{ex}{56}\RG{4.094}{23.652}{7.9701}{mi}{78}\RG{47.568}{22.561}{10.9092}{ex}{60}\RG{0.000}{20.379}{10.9092}{ex}{88}\RG{57.265}{18.206}{10.9091}{mi}{68}\RG{66.582}{18.206}{10.9091}{mi}{59}\RG{80.513}{18.206}{10.9091}{mi}{106}\RG{88.688}{18.206}{10.9091}{rm}{61}\RG{100.186}{18.206}{10.9091}{rm}{49}\RG{105.657}{18.206}{10.9091}{mi}{59}\RG{17.576}{10.016}{10.9091}{mi}{65}\RG{36.053}{10.016}{10.9091}{rm}{61}\RG{25.757}{8.379}{7.9701}{mi}{105}\RG{28.634}{8.379}{7.9701}{mi}{106}\RG{47.568}{2.925}{10.9092}{ex}{58}\RG{57.265}{-0.004}{10.9091}{rm}{48}\RG{62.721}{-0.004}{10.9091}{mi}{59}\RG{80.524}{-0.004}{10.9091}{mi}{106}\RG{88.687}{-0.004}{10.9091}{sy}{54}\RG{88.688}{-0.004}{10.9091}{rm}{61}\RG{100.186}{-0.004}{10.9091}{rm}{49}\RG{105.657}{-0.004}{10.9091}{mi}{58}\RG{1.027}{-3.072}{7.9701}{mi}{105}\RG{3.910}{-3.072}{7.9701}{rm}{61}\RG{10.493}{-3.072}{7.9701}{rm}{49}%
}
\]
因此

\[
\RB{77.746}{19.582}{15.211}{%
\RG{6.185}{13.638}{7.9701}{mi}{78}\RG{2.092}{10.365}{10.9092}{ex}{88}\RG{26.875}{4.506}{7.9701}{rm}{55}\RG{31.107}{4.506}{7.9701}{rm}{48}\RG{21.759}{0.002}{10.9091}{mi}{106}\RG{35.842}{0.002}{10.9091}{mi}{65}\RG{54.319}{0.002}{10.9091}{rm}{61}\RG{65.835}{0.002}{10.9091}{mi}{68}\RG{75.151}{0.002}{10.9091}{mi}{58}\RG{44.024}{-1.635}{7.9701}{mi}{105}\RG{46.902}{-1.635}{7.9701}{mi}{106}\RG{0.000}{-13.085}{7.9701}{mi}{105}\RG{2.877}{-13.085}{7.9701}{mi}{59}\RG{5.228}{-13.085}{7.9701}{mi}{106}\RG{9.120}{-13.086}{7.9701}{rm}{61}\RG{15.703}{-13.086}{7.9701}{rm}{49}%
}
\]
再令

\[
\RB{246.887}{19.580}{15.214}{%
\RG{35.643}{13.636}{7.9701}{mi}{78}\RG{202.564}{13.636}{7.9701}{mi}{78}\RG{32.458}{10.363}{10.9092}{ex}{89}\RG{198.470}{10.363}{10.9092}{ex}{88}\RG{144.609}{4.504}{7.9701}{mi}{78}\RG{229.089}{4.504}{7.9701}{mi}{106}\RG{232.973}{4.504}{7.9701}{sy}{0}\RG{239.559}{4.504}{7.9701}{rm}{49}\RG{0.000}{0.000}{10.9091}{mi}{112}\RG{5.489}{0.000}{10.9091}{rm}{40}\RG{9.731}{0.000}{10.9091}{mi}{116}\RG{13.670}{0.000}{10.9091}{rm}{41}\RG{20.925}{0.000}{10.9091}{rm}{61}\RG{46.398}{0.000}{10.9091}{rm}{40}\RG{50.640}{0.000}{10.9091}{mi}{116}\RG{57.004}{0.000}{10.9091}{sy}{0}\RG{67.914}{0.000}{10.9091}{mi}{120}\RG{77.530}{0.000}{10.9091}{rm}{41}\RG{81.772}{0.000}{10.9091}{mi}{59}\RG{108.423}{0.000}{10.9091}{mi}{113}\RG{113.700}{0.000}{10.9091}{rm}{40}\RG{117.943}{0.000}{10.9091}{mi}{116}\RG{121.882}{0.000}{10.9091}{rm}{41}\RG{129.137}{0.000}{10.9091}{rm}{61}\RG{140.670}{0.000}{10.9091}{mi}{116}\RG{155.102}{0.000}{10.9091}{sy}{0}\RG{166.011}{0.000}{10.9091}{mi}{112}\RG{171.500}{0.000}{10.9091}{rm}{40}\RG{175.742}{0.000}{10.9091}{mi}{116}\RG{179.682}{0.000}{10.9091}{rm}{41}\RG{186.937}{0.000}{10.9091}{rm}{61}\RG{216.046}{0.000}{10.9091}{mi}{99}\RG{225.149}{0.000}{10.9091}{mi}{116}\RG{244.292}{0.000}{10.9091}{mi}{58}\RG{74.149}{-1.637}{7.9701}{mi}{105}\RG{220.767}{-1.637}{7.9701}{mi}{106}\RG{32.576}{-13.088}{7.9701}{mi}{105}\RG{35.459}{-13.088}{7.9701}{rm}{61}\RG{42.043}{-13.088}{7.9701}{rm}{49}\RG{198.997}{-13.088}{7.9701}{mi}{106}\RG{202.881}{-13.088}{7.9701}{rm}{61}\RG{209.464}{-13.088}{7.9701}{rm}{49}%
}
\]
因为\allowbreak{}\hskip.1em \RB{5.846}{5.324}{2.614}{%
\RG{0.000}{0.002}{10.9091}{mi}{112}%
}\hskip.1em\allowbreak{} 为首一\allowbreak{}\hskip.1em \RB{10.111}{7.953}{0.500}{%
\RG{0.000}{0.002}{10.9091}{mi}{78}%
}\hskip.1em\allowbreak{} 次多项式，所以\allowbreak{}\hskip.1em \RB{66.499}{8.073}{2.745}{%
\RG{0.000}{0.002}{10.9091}{rm}{100}\RG{6.055}{0.002}{10.9091}{rm}{101}\RG{10.898}{0.002}{10.9091}{rm}{103}\RG{18.333}{0.002}{10.9091}{mi}{113}\RG{26.625}{0.002}{10.9091}{sy}{20}\RG{38.140}{0.002}{10.9091}{mi}{78}\RG{50.519}{0.002}{10.9091}{sy}{0}\RG{61.429}{0.002}{10.9091}{rm}{49}%
}\hskip.1em\allowbreak{} ；又因\allowbreak{}\hskip.1em \RB{43.629}{8.684}{3.225}{%
\RG{0.000}{0.002}{10.9091}{mi}{112}\RG{5.489}{0.002}{10.9091}{rm}{40}\RG{9.731}{0.002}{10.9091}{mi}{120}\RG{15.966}{-1.635}{7.9701}{mi}{105}\RG{19.347}{0.002}{10.9091}{rm}{41}\RG{26.602}{0.002}{10.9091}{rm}{61}\RG{38.111}{0.002}{10.9091}{rm}{48}%
}\hskip.1em\allowbreak{} ，有

\[
\RB{146.870}{19.581}{15.212}{%
\RG{109.511}{13.637}{7.9701}{mi}{78}\RG{105.417}{10.364}{10.9092}{ex}{88}\RG{33.732}{4.505}{7.9701}{mi}{78}\RG{72.279}{4.505}{7.9701}{mi}{78}\RG{84.589}{4.505}{7.9701}{rm}{84}\RG{0.000}{0.001}{10.9091}{mi}{118}\RG{8.709}{0.001}{10.9091}{rm}{58}\RG{11.731}{0.001}{10.9091}{rm}{61}\RG{23.229}{0.001}{10.9091}{rm}{40}\RG{27.497}{0.001}{10.9091}{mi}{120}\RG{41.800}{0.001}{10.9091}{mi}{59}\RG{46.633}{0.001}{10.9091}{mi}{58}\RG{51.477}{0.001}{10.9091}{mi}{58}\RG{56.320}{0.001}{10.9091}{mi}{58}\RG{61.153}{0.001}{10.9091}{mi}{59}\RG{65.997}{0.001}{10.9091}{mi}{120}\RG{80.347}{0.001}{10.9091}{rm}{41}\RG{93.902}{0.001}{10.9091}{rm}{61}\RG{122.993}{0.001}{10.9091}{mi}{99}\RG{132.096}{0.001}{10.9091}{mi}{67}\RG{144.275}{0.001}{10.9091}{mi}{58}\RG{127.714}{-1.636}{7.9701}{mi}{106}\RG{139.893}{-1.636}{7.9701}{mi}{106}\RG{33.732}{-2.696}{7.9701}{rm}{49}\RG{72.279}{-2.696}{7.9701}{mi}{78}\RG{105.943}{-13.086}{7.9701}{mi}{106}\RG{109.827}{-13.086}{7.9701}{rm}{61}\RG{116.410}{-13.086}{7.9701}{rm}{49}%
}
\]
将\allowbreak{}\hskip.1em \RB{8.366}{8.309}{0.502}{%
\RG{0.000}{-0.002}{10.9091}{mi}{65}%
}\hskip.1em\allowbreak{} 的第\allowbreak{}\hskip.1em \RB{4.831}{7.709}{2.738}{%
\RG{0.000}{-0.002}{10.9091}{mi}{106}%
}\hskip.1em\allowbreak{} 列替换为\allowbreak{}\hskip.1em \RB{5.606}{5.320}{0.622}{%
\RG{0.000}{-0.002}{10.9091}{mi}{118}%
}\hskip.1em\allowbreak{} 。一方面，沿该列展开得到\allowbreak{}\hskip.1em \RB{45.869}{9.900}{3.808}{%
\RG{0.000}{8.180}{10.9092}{ex}{80}\RG{11.515}{-3.228}{7.9701}{mi}{105}\RG{16.715}{-0.002}{10.9091}{mi}{120}\RG{22.950}{3.956}{7.9701}{mi}{78}\RG{22.950}{-3.058}{7.9701}{mi}{105}\RG{31.018}{-0.002}{10.9091}{mi}{65}\RG{39.200}{-1.639}{7.9701}{mi}{105}\RG{42.078}{-1.639}{7.9701}{mi}{106}%
}\hskip.1em\allowbreak{} ；另一方面，按该列的线性性展开，只有\allowbreak{}\hskip.1em \RB{20.692}{8.194}{3.760}{%
\RG{0.000}{0.003}{10.9091}{mi}{99}\RG{4.721}{-1.634}{7.9701}{mi}{106}\RG{9.103}{0.003}{10.9091}{mi}{67}\RG{16.900}{-1.634}{7.9701}{mi}{106}%
}\hskip.1em\allowbreak{} 对应的项不含重复列。故

\[
\RB{82.887}{19.583}{13.664}{%
\RG{4.094}{13.639}{7.9701}{mi}{78}\RG{0.000}{10.366}{10.9092}{ex}{88}\RG{23.811}{4.507}{7.9701}{mi}{78}\RG{17.576}{0.003}{10.9091}{mi}{120}\RG{31.879}{0.003}{10.9091}{mi}{65}\RG{50.357}{0.003}{10.9091}{rm}{61}\RG{61.873}{0.003}{10.9091}{mi}{99}\RG{70.976}{0.003}{10.9091}{mi}{68}\RG{80.292}{0.003}{10.9091}{mi}{58}\RG{40.061}{-1.634}{7.9701}{mi}{105}\RG{42.939}{-1.634}{7.9701}{mi}{106}\RG{66.594}{-1.634}{7.9701}{mi}{106}\RG{23.811}{-2.694}{7.9701}{mi}{105}\RG{1.027}{-13.084}{7.9701}{mi}{105}\RG{3.910}{-13.084}{7.9701}{rm}{61}\RG{10.493}{-13.084}{7.9701}{rm}{49}%
}
\]
对\allowbreak{}\hskip.1em \RB{4.831}{7.709}{2.738}{%
\RG{0.000}{-0.002}{10.9091}{mi}{106}%
}\hskip.1em\allowbreak{} 求和即得

\[
\RB{205.210}{19.589}{15.216}{%
\RG{109.136}{18.652}{10.9092}{ex}{34}\RG{185.116}{18.652}{10.9092}{ex}{35}\RG{6.185}{13.633}{7.9701}{mi}{78}\RG{137.472}{13.633}{7.9701}{mi}{78}\RG{2.091}{10.361}{10.9092}{ex}{88}\RG{134.288}{10.361}{10.9092}{ex}{89}\RG{27.993}{4.501}{7.9701}{mi}{78}\RG{21.758}{-0.003}{10.9091}{mi}{120}\RG{36.062}{-0.003}{10.9091}{mi}{65}\RG{54.539}{-0.003}{10.9091}{rm}{61}\RG{66.055}{-0.003}{10.9091}{mi}{113}\RG{71.317}{-0.003}{10.9091}{rm}{40}\RG{75.549}{-0.003}{10.9091}{rm}{49}\RG{81.004}{-0.003}{10.9091}{rm}{41}\RG{85.256}{-0.003}{10.9091}{mi}{68}\RG{97.621}{-0.003}{10.9091}{rm}{61}\RG{115.500}{-0.003}{10.9091}{rm}{49}\RG{123.379}{-0.003}{10.9091}{sy}{0}\RG{148.227}{-0.003}{10.9091}{rm}{40}\RG{152.460}{-0.003}{10.9091}{rm}{49}\RG{160.348}{-0.003}{10.9091}{sy}{0}\RG{171.257}{-0.003}{10.9091}{mi}{120}\RG{180.874}{-0.003}{10.9091}{rm}{41}\RG{193.299}{-0.003}{10.9091}{mi}{68}\RG{202.615}{-0.003}{10.9091}{mi}{58}\RG{44.244}{-1.639}{7.9701}{mi}{105}\RG{47.122}{-1.639}{7.9701}{mi}{106}\RG{177.492}{-1.639}{7.9701}{mi}{105}\RG{27.993}{-2.700}{7.9701}{mi}{105}\RG{0.000}{-13.090}{7.9701}{mi}{105}\RG{2.877}{-13.090}{7.9701}{mi}{59}\RG{5.228}{-13.090}{7.9701}{mi}{106}\RG{9.120}{-13.090}{7.9701}{rm}{61}\RG{15.703}{-13.090}{7.9701}{rm}{49}\RG{134.406}{-13.090}{7.9701}{mi}{105}\RG{137.289}{-13.090}{7.9701}{rm}{61}\RG{143.872}{-13.090}{7.9701}{rm}{49}%
}
\]
结合\allowbreak{}\hskip.1em \RB{46.333}{7.948}{0.743}{%
\RG{0.000}{-0.003}{10.9091}{mi}{78}\RG{12.985}{-0.003}{10.9091}{rm}{61}\RG{24.483}{-0.003}{10.9091}{rm}{50}\RG{29.938}{-0.003}{10.9091}{rm}{48}\RG{35.393}{-0.003}{10.9091}{rm}{49}\RG{40.847}{-0.003}{10.9091}{rm}{57}%
}\hskip.1em\allowbreak{} 为奇数，所求结果为

\[
\RB{179.504}{19.590}{15.275}{%
\RG{0.000}{18.653}{10.9092}{ex}{34}\RG{78.978}{18.653}{10.9092}{ex}{35}\RG{25.152}{13.635}{7.9701}{rm}{50}\RG{29.384}{13.635}{7.9701}{rm}{48}\RG{33.617}{13.635}{7.9701}{rm}{49}\RG{37.849}{13.635}{7.9701}{rm}{57}\RG{26.650}{10.362}{10.9092}{ex}{89}\RG{104.039}{10.362}{10.9092}{ex}{89}\RG{6.364}{-0.002}{10.9091}{rm}{50}\RG{14.240}{-0.002}{10.9091}{rm}{43}\RG{42.088}{-0.002}{10.9091}{rm}{40}\RG{46.331}{-0.002}{10.9091}{mi}{120}\RG{58.371}{-0.002}{10.9091}{sy}{0}\RG{69.281}{-0.002}{10.9091}{rm}{49}\RG{74.735}{-0.002}{10.9091}{rm}{41}\RG{134.857}{-0.002}{10.9091}{rm}{40}\RG{139.100}{-0.002}{10.9091}{mi}{120}\RG{152.141}{-0.002}{10.9091}{sy}{0}\RG{163.050}{-0.002}{10.9091}{mi}{120}\RG{172.666}{-0.002}{10.9091}{rm}{41}\RG{176.909}{-0.002}{10.9091}{mi}{58}\RG{52.566}{-1.638}{7.9701}{mi}{105}\RG{145.335}{-1.638}{7.9701}{mi}{106}\RG{169.285}{-1.638}{7.9701}{mi}{105}\RG{26.768}{-13.089}{7.9701}{mi}{105}\RG{29.651}{-13.089}{7.9701}{rm}{61}\RG{36.235}{-13.089}{7.9701}{rm}{49}\RG{87.160}{-13.149}{7.9701}{rm}{49}\RG{91.394}{-13.149}{7.9701}{sy}{20}\RG{97.981}{-13.149}{7.9701}{mi}{105}\RG{100.858}{-13.149}{7.9701}{mi}{60}\RG{107.441}{-13.149}{7.9701}{mi}{106}\RG{111.334}{-13.149}{7.9701}{sy}{20}\RG{117.921}{-13.149}{7.9701}{rm}{50}\RG{122.153}{-13.149}{7.9701}{rm}{48}\RG{126.385}{-13.149}{7.9701}{rm}{49}\RG{130.617}{-13.149}{7.9701}{rm}{57}%
}
\]
小结：把代数余子式的线性组合解释为“替换一列后的行列式”，再用\allowbreak{}\hskip.1em \RB{11.740}{9.903}{0.620}{%
\RG{0.000}{0.000}{10.9091}{mi}{116}\RG{3.939}{3.959}{7.9701}{mi}{78}%
}\hskip.1em\allowbreak{} 除以\allowbreak{}\hskip.1em \RB{48.643}{8.682}{3.806}{%
\RG{0.000}{8.182}{10.9092}{ex}{81}\RG{10.303}{-3.226}{7.9701}{mi}{105}\RG{13.684}{0.000}{10.9091}{rm}{40}\RG{17.927}{0.000}{10.9091}{mi}{116}\RG{24.378}{0.000}{10.9091}{sy}{0}\RG{35.374}{0.000}{10.9091}{mi}{120}\RG{41.609}{-1.637}{7.9701}{mi}{105}\RG{44.990}{0.000}{10.9091}{rm}{41}%
}\hskip.1em\allowbreak{} 所得的余式降次。证明不要求\allowbreak{}\hskip.1em \RB{9.262}{5.317}{2.222}{%
\RG{0.000}{-0.005}{10.9091}{mi}{120}\RG{6.235}{-1.642}{7.9701}{mi}{105}%
}\hskip.1em\allowbreak{} 两两不同。

题目来源：谢启鸿《复旦大学高等代数习题集（2023版）》，第80页。

% Original page 4.
\section*{2019A03}
设\allowbreak{}\hskip.1em \RB{12.067}{8.191}{2.216}{%
\RG{0.000}{-0.000}{10.9091}{mi}{83}\RG{6.689}{-1.636}{7.9701}{mi}{110}%
}\hskip.1em\allowbreak{} 为\allowbreak{}\hskip.1em \RB{46.859}{8.682}{3.227}{%
\RG{0.000}{-0.000}{10.9091}{sy}{102}\RG{5.455}{-0.000}{10.9091}{rm}{49}\RG{10.909}{-0.000}{10.9091}{mi}{59}\RG{15.742}{-0.000}{10.9091}{mi}{58}\RG{20.586}{-0.000}{10.9091}{mi}{58}\RG{25.429}{-0.000}{10.9091}{mi}{58}\RG{30.262}{-0.000}{10.9091}{mi}{59}\RG{35.106}{-0.000}{10.9091}{mi}{110}\RG{41.700}{-0.000}{10.9091}{sy}{103}%
}\hskip.1em\allowbreak{} 上的置换全体，乘积约定为\allowbreak{}\hskip.1em \RB{80.661}{8.682}{3.227}{%
\RG{0.000}{-0.000}{10.9091}{rm}{40}\RG{4.243}{-0.000}{10.9091}{mi}{28}\RG{10.243}{-0.000}{10.9091}{mi}{27}\RG{16.872}{-0.000}{10.9091}{rm}{41}\RG{21.105}{-0.000}{10.9091}{rm}{40}\RG{25.356}{-0.000}{10.9091}{mi}{105}\RG{29.115}{-0.000}{10.9091}{rm}{41}\RG{36.370}{-0.000}{10.9091}{rm}{61}\RG{47.903}{-0.000}{10.9091}{mi}{28}\RG{53.906}{-0.000}{10.9091}{rm}{40}\RG{58.149}{-0.000}{10.9091}{mi}{27}\RG{64.775}{-0.000}{10.9091}{rm}{40}\RG{69.017}{-0.000}{10.9091}{mi}{105}\RG{72.776}{-0.000}{10.9091}{rm}{41}\RG{77.008}{-0.000}{10.9091}{rm}{41}%
}\hskip.1em\allowbreak{} ，\allowbreak{}\hskip.1em \RB{5.191}{5.322}{0.620}{%
\RG{0.000}{-0.000}{10.9091}{mi}{101}%
}\hskip.1em\allowbreak{} 为恒等置换。设\allowbreak{}\hskip.1em \RB{44.512}{5.322}{2.605}{%
\RG{0.000}{-0.000}{10.9091}{mi}{101}\RG{5.079}{-1.636}{7.9701}{rm}{49}\RG{9.812}{-0.000}{10.9091}{mi}{59}\RG{14.645}{-0.000}{10.9091}{mi}{58}\RG{19.489}{-0.000}{10.9091}{mi}{58}\RG{24.321}{-0.000}{10.9091}{mi}{58}\RG{29.165}{-0.000}{10.9091}{mi}{59}\RG{34.009}{-0.000}{10.9091}{mi}{101}\RG{39.134}{-1.636}{7.9701}{mi}{110}%
}\hskip.1em\allowbreak{} 为标准单位列向量，并定义

\[
\RB{106.446}{8.682}{4.713}{%
\RG{0.000}{0.000}{10.9091}{mi}{80}\RG{15.658}{0.000}{10.9091}{rm}{61}\RG{27.156}{0.000}{10.9091}{rm}{40}\RG{31.415}{0.000}{10.9091}{mi}{101}\RG{52.939}{0.000}{10.9091}{mi}{59}\RG{57.772}{0.000}{10.9091}{mi}{58}\RG{62.616}{0.000}{10.9091}{mi}{58}\RG{67.449}{0.000}{10.9091}{mi}{58}\RG{72.292}{0.000}{10.9091}{mi}{59}\RG{77.136}{0.000}{10.9091}{mi}{101}\RG{99.610}{0.000}{10.9091}{rm}{41}\RG{103.852}{0.000}{10.9091}{mi}{58}\RG{7.004}{-1.637}{7.9701}{mi}{27}\RG{36.495}{-2.220}{7.9701}{mi}{27}\RG{41.620}{-2.220}{7.9701}{rm}{40}\RG{44.912}{-2.220}{7.9701}{rm}{49}\RG{49.144}{-2.220}{7.9701}{rm}{41}\RG{82.262}{-2.220}{7.9701}{mi}{27}\RG{87.387}{-2.220}{7.9701}{rm}{40}\RG{90.680}{-2.220}{7.9701}{mi}{110}\RG{95.818}{-2.220}{7.9701}{rm}{41}%
}
\]
设\allowbreak{}\hskip.1em \RB{8.366}{8.310}{0.501}{%
\RG{0.000}{-0.001}{10.9091}{mi}{65}%
}\hskip.1em\allowbreak{} 为\allowbreak{}\hskip.1em \RB{6.729}{5.321}{0.621}{%
\RG{0.000}{-0.001}{10.9091}{mi}{110}%
}\hskip.1em\allowbreak{} 阶方阵。证明：

(1)第一类初等矩阵\allowbreak{}\hskip.1em \RB{13.673}{7.955}{3.758}{%
\RG{0.000}{0.004}{10.9091}{mi}{80}\RG{7.004}{-1.632}{7.9701}{mi}{105}\RG{9.881}{-1.632}{7.9701}{mi}{106}%
}\hskip.1em\allowbreak{} 、基础循环矩阵\allowbreak{}\hskip.1em \RB{7.406}{7.955}{0.736}{%
\RG{0.000}{0.004}{10.9091}{mi}{74}%
}\hskip.1em\allowbreak{} 和反单位阵都是置换矩阵。(2)\allowbreak{}\hskip.1em \RB{85.518}{10.440}{3.224}{%
\RG{0.000}{0.003}{10.9091}{rm}{100}\RG{6.055}{0.003}{10.9091}{rm}{101}\RG{10.898}{0.003}{10.9091}{rm}{116}\RG{16.969}{0.003}{10.9091}{mi}{80}\RG{23.973}{-1.634}{7.9701}{mi}{27}\RG{32.627}{0.003}{10.9091}{rm}{61}\RG{44.126}{0.003}{10.9091}{rm}{40}\RG{48.385}{0.003}{10.9091}{sy}{0}\RG{56.869}{0.003}{10.9091}{rm}{49}\RG{62.324}{0.003}{10.9091}{rm}{41}\RG{66.567}{3.962}{7.9701}{mi}{78}\RG{74.138}{3.962}{7.9701}{rm}{40}\RG{77.431}{3.962}{7.9701}{mi}{27}\RG{82.556}{3.962}{7.9701}{rm}{41}%
}\hskip.1em\allowbreak{} ，其中\allowbreak{}\hskip.1em \RB{24.475}{8.685}{3.224}{%
\RG{0.000}{0.003}{10.9091}{mi}{78}\RG{9.954}{0.003}{10.9091}{rm}{40}\RG{14.197}{0.003}{10.9091}{mi}{27}\RG{20.822}{0.003}{10.9091}{rm}{41}%
}\hskip.1em\allowbreak{} 为排列\allowbreak{}\hskip.1em \RB{74.350}{8.685}{3.224}{%
\RG{0.000}{0.003}{10.9091}{rm}{40}\RG{4.242}{0.003}{10.9091}{mi}{27}\RG{10.867}{0.003}{10.9091}{rm}{40}\RG{15.100}{0.003}{10.9091}{rm}{49}\RG{20.554}{0.003}{10.9091}{rm}{41}\RG{24.807}{0.003}{10.9091}{mi}{59}\RG{29.640}{0.003}{10.9091}{mi}{58}\RG{34.484}{0.003}{10.9091}{mi}{58}\RG{39.316}{0.003}{10.9091}{mi}{58}\RG{44.160}{0.003}{10.9091}{mi}{59}\RG{49.004}{0.003}{10.9091}{mi}{27}\RG{55.674}{0.003}{10.9091}{rm}{40}\RG{59.916}{0.003}{10.9091}{mi}{110}\RG{66.465}{0.003}{10.9091}{rm}{41}\RG{70.697}{0.003}{10.9091}{rm}{41}%
}\hskip.1em\allowbreak{} 的逆序数。

\[
\RB{231.956}{9.936}{3.275}{%
\RG{179.881}{3.960}{7.9701}{sy}{0}\RG{186.468}{3.960}{7.9701}{rm}{49}\RG{214.265}{3.960}{7.9701}{rm}{84}\RT{0.000}{0.001}{10.9091}{(}\RT{4.244}{0.001}{10.9091}{3}\RT{9.698}{0.001}{10.9091}{)}\RG{19.397}{0.001}{10.9091}{mi}{80}\RG{39.733}{0.001}{10.9091}{rm}{61}\RG{51.248}{0.001}{10.9091}{mi}{80}\RG{63.429}{0.001}{10.9091}{mi}{80}\RT{76.056}{0.001}{10.9091}{，}\RG{86.965}{0.001}{10.9091}{mi}{80}\RG{101.426}{0.001}{10.9091}{rm}{61}\RG{112.941}{0.001}{10.9091}{mi}{73}\RT{123.373}{0.001}{10.9091}{，}\RG{134.282}{0.001}{10.9091}{mi}{80}\RG{159.846}{0.001}{10.9091}{rm}{61}\RG{171.362}{0.001}{10.9091}{mi}{80}\RG{194.231}{0.001}{10.9091}{rm}{61}\RG{205.746}{0.001}{10.9091}{mi}{80}\RT{220.547}{0.001}{10.9091}{。}\RG{146.411}{-0.093}{5.9776}{sy}{0}\RG{152.167}{-0.093}{5.9776}{rm}{49}\RG{26.401}{-1.635}{7.9701}{mi}{28}\RG{31.072}{-1.635}{7.9701}{mi}{27}\RG{58.252}{-1.635}{7.9701}{mi}{28}\RG{70.433}{-1.635}{7.9701}{mi}{27}\RG{93.969}{-1.635}{7.9701}{mi}{101}\RG{117.736}{-1.635}{7.9701}{mi}{110}\RG{141.286}{-2.362}{7.9701}{mi}{27}\RG{178.365}{-2.695}{7.9701}{mi}{27}\RG{212.750}{-2.695}{7.9701}{mi}{27}%
}
\]
(4)\allowbreak{}\hskip.1em \RB{20.475}{8.311}{2.216}{%
\RG{0.000}{-0.000}{10.9091}{mi}{65}\RG{8.182}{-0.000}{10.9091}{mi}{80}\RG{15.185}{-1.636}{7.9701}{mi}{27}%
}\hskip.1em\allowbreak{} 的第\allowbreak{}\hskip.1em \RB{3.696}{7.711}{0.620}{%
\RG{0.000}{-0.000}{10.9091}{mi}{105}%
}\hskip.1em\allowbreak{} 列是\allowbreak{}\hskip.1em \RB{8.366}{8.311}{0.500}{%
\RG{0.000}{-0.000}{10.9091}{mi}{65}%
}\hskip.1em\allowbreak{} 的第\allowbreak{}\hskip.1em \RB{18.279}{8.682}{3.227}{%
\RG{0.000}{-0.000}{10.9091}{mi}{27}\RG{6.625}{-0.000}{10.9091}{rm}{40}\RG{10.867}{-0.000}{10.9091}{mi}{105}\RG{14.626}{-0.000}{10.9091}{rm}{41}%
}\hskip.1em\allowbreak{} 列；\allowbreak{}\hskip.1em \RB{23.167}{9.935}{3.277}{%
\RG{0.000}{-0.000}{10.9091}{mi}{80}\RG{7.004}{-2.697}{7.9701}{mi}{27}\RG{8.519}{3.959}{7.9701}{rm}{84}\RG{14.801}{-0.000}{10.9091}{mi}{65}%
}\hskip.1em\allowbreak{} 的第\allowbreak{}\hskip.1em \RB{3.696}{7.711}{0.620}{%
\RG{0.000}{-0.000}{10.9091}{mi}{105}%
}\hskip.1em\allowbreak{} 行是\allowbreak{}\hskip.1em \RB{8.366}{8.311}{0.500}{%
\RG{0.000}{-0.000}{10.9091}{mi}{65}%
}\hskip.1em\allowbreak{} 的第\allowbreak{}\hskip.1em \RB{18.279}{8.682}{3.227}{%
\RG{0.000}{-0.000}{10.9091}{mi}{27}\RG{6.625}{-0.000}{10.9091}{rm}{40}\RG{10.867}{-0.000}{10.9091}{mi}{105}\RG{14.626}{-0.000}{10.9091}{rm}{41}%
}\hskip.1em\allowbreak{} 行。

证明：(1)对换\allowbreak{}\hskip.1em \RB{41.558}{8.687}{3.222}{%
\RG{0.000}{0.005}{10.9091}{mi}{25}\RG{9.640}{0.005}{10.9091}{rm}{61}\RG{21.138}{0.005}{10.9091}{rm}{40}\RG{25.397}{0.005}{10.9091}{mi}{105}\RG{32.783}{0.005}{10.9091}{mi}{106}\RG{37.905}{0.005}{10.9091}{rm}{41}%
}\hskip.1em\allowbreak{} 只交换\allowbreak{}\hskip.1em \RB{13.416}{7.716}{2.731}{%
\RG{0.000}{0.005}{10.9091}{mi}{105}\RG{3.753}{0.005}{10.9091}{mi}{59}\RG{8.586}{0.005}{10.9091}{mi}{106}%
}\hskip.1em\allowbreak{} 而固定其余指标，故\allowbreak{}\hskip.1em \RB{41.132}{7.956}{3.757}{%
\RG{0.000}{0.005}{10.9091}{mi}{80}\RG{7.004}{-1.631}{7.9701}{mi}{105}\RG{9.881}{-1.631}{7.9701}{mi}{106}\RG{17.300}{0.005}{10.9091}{rm}{61}\RG{28.815}{0.005}{10.9091}{mi}{80}\RG{35.818}{-1.631}{7.9701}{mi}{25}%
}\hskip.1em\allowbreak{} 。对基础循环矩阵

\[
\RB{80.120}{13.943}{19.773}{%
\RG{21.644}{13.007}{10.9092}{ex}{32}\RG{66.951}{13.007}{10.9092}{ex}{33}\RG{30.281}{1.633}{10.9091}{rm}{48}\RG{45.698}{1.633}{10.9091}{mi}{73}\RG{50.493}{-0.004}{7.9701}{mi}{110}\RG{55.632}{-0.004}{7.9701}{sy}{0}\RG{62.219}{-0.004}{7.9701}{rm}{49}\RG{0.000}{-5.647}{10.9091}{mi}{74}\RG{10.129}{-5.647}{10.9091}{rm}{61}\RG{77.406}{-5.647}{10.9091}{mi}{59}\RG{30.281}{-13.542}{10.9091}{rm}{49}\RG{53.594}{-13.542}{10.9091}{rm}{48}%
}
\]
当\allowbreak{}\hskip.1em \RB{26.358}{7.770}{0.616}{%
\RG{0.000}{0.004}{10.9091}{mi}{110}\RG{9.680}{0.004}{10.9091}{rm}{61}\RG{21.287}{0.004}{10.9091}{rm}{49}%
}\hskip.1em\allowbreak{} 时取\allowbreak{}\hskip.1em \RB{35.178}{8.686}{3.223}{%
\RG{0.000}{0.004}{10.9091}{mi}{74}\RG{10.230}{0.004}{10.9091}{rm}{61}\RG{21.837}{0.004}{10.9091}{rm}{40}\RG{26.070}{0.004}{10.9091}{rm}{49}\RG{31.525}{0.004}{10.9091}{rm}{41}%
}\hskip.1em\allowbreak{} ；一般地，有\allowbreak{}\hskip.1em \RB{42.117}{7.955}{2.213}{%
\RG{0.000}{0.004}{10.9091}{mi}{74}\RG{7.091}{0.004}{10.9091}{mi}{101}\RG{12.178}{-1.633}{7.9701}{rm}{49}\RG{20.043}{0.004}{10.9091}{rm}{61}\RG{31.660}{0.004}{10.9091}{mi}{101}\RG{36.739}{-1.633}{7.9701}{mi}{110}%
}\hskip.1em\allowbreak{} 、\allowbreak{}\hskip.1em \RB{48.920}{7.955}{2.213}{%
\RG{0.000}{0.004}{10.9091}{mi}{74}\RG{7.091}{0.004}{10.9091}{mi}{101}\RG{12.178}{-1.633}{7.9701}{mi}{105}\RG{18.691}{0.004}{10.9091}{rm}{61}\RG{30.308}{0.004}{10.9091}{mi}{101}\RG{35.388}{-1.633}{7.9701}{mi}{105}\RG{38.271}{-1.633}{7.9701}{sy}{0}\RG{44.857}{-1.633}{7.9701}{rm}{49}%
}\hskip.1em\allowbreak{} （\allowbreak{}\hskip.1em \RB{45.440}{7.770}{1.991}{%
\RG{0.000}{0.004}{10.9091}{rm}{50}\RG{8.586}{0.004}{10.9091}{sy}{20}\RG{20.203}{0.004}{10.9091}{mi}{105}\RG{27.094}{0.004}{10.9091}{sy}{20}\RG{38.711}{0.004}{10.9091}{mi}{110}%
}\hskip.1em\allowbreak{} ），故它对应置换\allowbreak{}\hskip.1em \RB{41.057}{8.686}{3.223}{%
\RG{0.000}{0.004}{10.9091}{mi}{26}\RG{5.640}{0.004}{10.9091}{rm}{40}\RG{9.873}{0.004}{10.9091}{rm}{49}\RG{15.327}{0.004}{10.9091}{rm}{41}\RG{22.691}{0.004}{10.9091}{rm}{61}\RG{34.328}{0.004}{10.9091}{mi}{110}%
}\hskip.1em\allowbreak{} 、\allowbreak{}\hskip.1em \RB{54.591}{8.681}{3.228}{%
\RG{0.000}{-0.001}{10.9091}{mi}{26}\RG{5.640}{-0.001}{10.9091}{rm}{40}\RG{9.882}{-0.001}{10.9091}{mi}{105}\RG{13.641}{-0.001}{10.9091}{rm}{41}\RG{20.896}{-0.001}{10.9091}{rm}{61}\RG{32.429}{-0.001}{10.9091}{mi}{105}\RG{38.611}{-0.001}{10.9091}{sy}{0}\RG{49.520}{-0.001}{10.9091}{rm}{49}%
}\hskip.1em\allowbreak{} 。反单位阵对应置换\allowbreak{}\hskip.1em \RB{75.440}{8.681}{3.228}{%
\RG{0.000}{-0.001}{10.9091}{mi}{20}\RG{6.286}{-0.001}{10.9091}{rm}{40}\RG{10.528}{-0.001}{10.9091}{mi}{105}\RG{14.286}{-0.001}{10.9091}{rm}{41}\RG{21.541}{-0.001}{10.9091}{rm}{61}\RG{33.074}{-0.001}{10.9091}{mi}{110}\RG{42.046}{-0.001}{10.9091}{rm}{43}\RG{52.945}{-0.001}{10.9091}{rm}{49}\RG{60.834}{-0.001}{10.9091}{sy}{0}\RG{71.744}{-0.001}{10.9091}{mi}{105}%
}\hskip.1em\allowbreak{} 。

(2)在行列式的按列置换展开式

\[
\RB{166.812}{17.652}{15.213}{%
\RG{114.016}{13.637}{7.9701}{mi}{110}\RG{46.785}{10.364}{10.9092}{ex}{88}\RG{109.615}{10.364}{10.9092}{ex}{89}\RG{87.611}{4.505}{7.9701}{mi}{78}\RG{95.181}{4.505}{7.9701}{rm}{40}\RG{98.474}{4.505}{7.9701}{mi}{25}\RG{103.623}{4.505}{7.9701}{rm}{41}\RG{0.000}{0.001}{10.9091}{rm}{100}\RG{6.055}{0.001}{10.9091}{rm}{101}\RG{10.898}{0.001}{10.9091}{rm}{116}\RG{16.969}{0.001}{10.9091}{mi}{80}\RG{32.627}{0.001}{10.9091}{rm}{61}\RG{65.186}{0.001}{10.9091}{rm}{40}\RG{69.429}{0.001}{10.9091}{sy}{0}\RG{77.914}{0.001}{10.9091}{rm}{49}\RG{83.368}{0.001}{10.9091}{rm}{41}\RG{123.937}{0.001}{10.9091}{rm}{40}\RG{128.180}{0.001}{10.9091}{mi}{80}\RG{140.807}{0.001}{10.9091}{rm}{41}\RG{23.973}{-1.636}{7.9701}{mi}{27}\RG{135.183}{-1.636}{7.9701}{mi}{27}\RG{145.049}{-2.219}{7.9701}{mi}{25}\RG{150.198}{-2.219}{7.9701}{rm}{40}\RG{153.491}{-2.219}{7.9701}{mi}{106}\RG{157.375}{-2.219}{7.9701}{rm}{41}\RG{160.669}{-2.219}{7.9701}{mi}{59}\RG{163.021}{-2.219}{7.9701}{mi}{106}\RG{109.233}{-13.087}{7.9701}{mi}{106}\RG{113.117}{-13.087}{7.9701}{rm}{61}\RG{119.700}{-13.087}{7.9701}{rm}{49}\RG{44.142}{-13.264}{7.9701}{mi}{25}\RG{49.291}{-13.264}{7.9701}{sy}{50}\RG{54.936}{-13.264}{7.9701}{mi}{83}\RG{60.086}{-14.260}{5.9776}{mi}{110}%
}
\]
中，第\allowbreak{}\hskip.1em \RB{4.831}{7.712}{2.735}{%
\RG{0.000}{0.001}{10.9091}{mi}{106}%
}\hskip.1em\allowbreak{} 列唯一的非零元在第\allowbreak{}\hskip.1em \RB{19.638}{8.683}{3.226}{%
\RG{0.000}{0.001}{10.9091}{mi}{27}\RG{6.625}{0.001}{10.9091}{rm}{40}\RG{10.868}{0.001}{10.9091}{mi}{106}\RG{15.985}{0.001}{10.9091}{rm}{41}%
}\hskip.1em\allowbreak{} 行，值为\allowbreak{}\hskip.1em \RB{5.071}{7.767}{0.500}{%
\RG{0.000}{0.001}{10.9091}{rm}{49}%
}\hskip.1em\allowbreak{} 。因此只有\allowbreak{}\hskip.1em \RB{27.841}{5.203}{0.619}{%
\RG{0.000}{0.001}{10.9091}{mi}{25}\RG{9.640}{0.001}{10.9091}{rm}{61}\RG{21.155}{0.001}{10.9091}{mi}{27}%
}\hskip.1em\allowbreak{} 的项非零，从而\allowbreak{}\hskip.1em \RB{85.519}{10.438}{3.226}{%
\RG{0.000}{0.001}{10.9091}{rm}{100}\RG{6.055}{0.001}{10.9091}{rm}{101}\RG{10.898}{0.001}{10.9091}{rm}{116}\RG{16.970}{0.001}{10.9091}{mi}{80}\RG{23.974}{-1.635}{7.9701}{mi}{27}\RG{32.628}{0.001}{10.9091}{rm}{61}\RG{44.127}{0.001}{10.9091}{rm}{40}\RG{48.385}{0.001}{10.9091}{sy}{0}\RG{56.870}{0.001}{10.9091}{rm}{49}\RG{62.325}{0.001}{10.9091}{rm}{41}\RG{66.568}{3.960}{7.9701}{mi}{78}\RG{74.138}{3.960}{7.9701}{rm}{40}\RG{77.432}{3.960}{7.9701}{mi}{27}\RG{82.557}{3.960}{7.9701}{rm}{41}%
}\hskip.1em\allowbreak{} 。

(3)由定义，对每个\allowbreak{}\hskip.1em \RB{3.696}{7.715}{0.616}{%
\RG{0.000}{0.004}{10.9091}{mi}{105}%
}\hskip.1em\allowbreak{} 有

\[
\RB{169.057}{7.951}{4.713}{%
\RG{0.000}{0.000}{10.9091}{mi}{80}\RG{12.180}{0.000}{10.9091}{mi}{80}\RG{24.808}{0.000}{10.9091}{mi}{101}\RG{36.299}{0.000}{10.9091}{rm}{61}\RG{47.814}{0.000}{10.9091}{mi}{80}\RG{59.994}{0.000}{10.9091}{mi}{101}\RG{83.198}{0.000}{10.9091}{rm}{61}\RG{94.713}{0.000}{10.9091}{mi}{101}\RG{129.181}{0.000}{10.9091}{rm}{61}\RG{140.696}{0.000}{10.9091}{mi}{80}\RG{158.002}{0.000}{10.9091}{mi}{101}\RG{166.462}{0.000}{10.9091}{mi}{58}\RG{7.004}{-1.636}{7.9701}{mi}{28}\RG{19.184}{-1.636}{7.9701}{mi}{27}\RG{29.887}{-1.636}{7.9701}{mi}{105}\RG{54.818}{-1.636}{7.9701}{mi}{28}\RG{147.700}{-1.636}{7.9701}{mi}{28}\RG{152.371}{-1.636}{7.9701}{mi}{27}\RG{163.081}{-1.636}{7.9701}{mi}{105}\RG{65.074}{-2.220}{7.9701}{mi}{27}\RG{70.200}{-2.220}{7.9701}{rm}{40}\RG{73.493}{-2.220}{7.9701}{mi}{105}\RG{76.376}{-2.220}{7.9701}{rm}{41}\RG{99.793}{-2.220}{7.9701}{mi}{28}\RG{104.471}{-2.220}{7.9701}{rm}{40}\RG{107.764}{-2.220}{7.9701}{mi}{27}\RG{112.890}{-2.220}{7.9701}{rm}{40}\RG{116.183}{-2.220}{7.9701}{mi}{105}\RG{119.066}{-2.220}{7.9701}{rm}{41}\RG{122.358}{-2.220}{7.9701}{rm}{41}%
}
\]
这些矩阵在一组基上的作用相同，故\allowbreak{}\hskip.1em \RB{56.325}{7.948}{2.236}{%
\RG{0.000}{-0.003}{10.9091}{mi}{80}\RG{7.004}{-1.640}{7.9701}{mi}{28}\RG{11.675}{-1.640}{7.9701}{mi}{27}\RG{20.336}{-0.003}{10.9091}{rm}{61}\RG{31.851}{-0.003}{10.9091}{mi}{80}\RG{38.855}{-1.640}{7.9701}{mi}{28}\RG{44.031}{-0.003}{10.9091}{mi}{80}\RG{51.035}{-1.640}{7.9701}{mi}{27}%
}\hskip.1em\allowbreak{} 。又\allowbreak{}\hskip.1em \RB{42.543}{7.948}{2.220}{%
\RG{0.000}{-0.003}{10.9091}{mi}{80}\RG{7.004}{-1.640}{7.9701}{mi}{101}\RG{11.431}{-0.003}{10.9091}{mi}{101}\RG{16.510}{-1.640}{7.9701}{mi}{105}\RG{22.922}{-0.003}{10.9091}{rm}{61}\RG{34.437}{-0.003}{10.9091}{mi}{101}\RG{39.516}{-1.640}{7.9701}{mi}{105}%
}\hskip.1em\allowbreak{} ，所以\allowbreak{}\hskip.1em \RB{36.149}{7.948}{2.220}{%
\RG{0.000}{-0.003}{10.9091}{mi}{80}\RG{7.003}{-1.640}{7.9701}{mi}{101}\RG{14.460}{-0.003}{10.9091}{rm}{61}\RG{25.975}{-0.003}{10.9091}{mi}{73}\RG{30.771}{-1.640}{7.9701}{mi}{110}%
}\hskip.1em\allowbreak{} ，进而

\[
\RB{112.441}{7.954}{2.941}{%
\RG{0.000}{0.003}{10.9091}{mi}{80}\RG{22.534}{0.003}{10.9091}{mi}{80}\RG{38.192}{0.003}{10.9091}{rm}{61}\RG{49.707}{0.003}{10.9091}{mi}{80}\RG{62.335}{0.003}{10.9091}{mi}{80}\RG{87.900}{0.003}{10.9091}{rm}{61}\RG{99.415}{0.003}{10.9091}{mi}{73}\RG{109.847}{0.003}{10.9091}{mi}{58}\RG{12.129}{-0.092}{5.9776}{sy}{0}\RG{17.885}{-0.092}{5.9776}{rm}{49}\RG{74.464}{-0.092}{5.9776}{sy}{0}\RG{80.220}{-0.092}{5.9776}{rm}{49}\RG{29.538}{-1.634}{7.9701}{mi}{27}\RG{56.711}{-1.634}{7.9701}{mi}{27}\RG{104.210}{-1.634}{7.9701}{mi}{110}\RG{7.004}{-2.361}{7.9701}{mi}{27}\RG{69.339}{-2.361}{7.9701}{mi}{27}%
}
\]
另一方面，\allowbreak{}\hskip.1em \RB{12.294}{7.950}{2.218}{%
\RG{0.000}{-0.001}{10.9091}{mi}{80}\RG{7.004}{-1.638}{7.9701}{mi}{27}%
}\hskip.1em\allowbreak{} 的各列是标准单位向量的一个排列，故

\[
\RB{128.325}{10.475}{5.716}{%
\RG{12.762}{4.499}{7.9701}{rm}{84}\RG{62.805}{4.499}{7.9701}{rm}{84}\RG{0.000}{-0.005}{10.9091}{rm}{40}\RG{4.243}{-0.005}{10.9091}{mi}{80}\RG{19.044}{-0.005}{10.9091}{mi}{80}\RG{31.671}{-0.005}{10.9091}{rm}{41}\RG{46.209}{-0.005}{10.9091}{rm}{61}\RG{57.724}{-0.005}{10.9091}{mi}{101}\RG{77.898}{-0.005}{10.9091}{mi}{101}\RG{102.101}{-0.005}{10.9091}{rm}{61}\RG{113.617}{-0.005}{10.9091}{mi}{14}\RG{125.730}{-0.005}{10.9091}{mi}{58}\RG{26.047}{-1.642}{7.9701}{mi}{27}\RG{35.913}{-1.642}{7.9701}{mi}{105}\RG{38.791}{-1.642}{7.9701}{mi}{106}\RG{118.465}{-1.642}{7.9701}{mi}{105}\RG{121.342}{-1.642}{7.9701}{mi}{106}\RG{82.977}{-2.225}{7.9701}{mi}{27}\RG{88.102}{-2.225}{7.9701}{rm}{40}\RG{91.396}{-2.225}{7.9701}{mi}{106}\RG{95.280}{-2.225}{7.9701}{rm}{41}\RG{11.246}{-2.702}{7.9701}{mi}{27}\RG{62.805}{-3.224}{7.9701}{mi}{27}\RG{67.930}{-3.224}{7.9701}{rm}{40}\RG{71.223}{-3.224}{7.9701}{mi}{105}\RG{74.106}{-3.224}{7.9701}{rm}{41}%
}
\]
于是\allowbreak{}\hskip.1em \RB{85.958}{9.935}{3.276}{%
\RG{0.000}{0.001}{10.9091}{mi}{80}\RG{7.004}{-2.363}{7.9701}{mi}{27}\RG{12.129}{-0.094}{5.9776}{sy}{0}\RG{17.885}{-0.094}{5.9776}{rm}{49}\RG{25.564}{0.001}{10.9091}{rm}{61}\RG{37.080}{0.001}{10.9091}{mi}{80}\RG{44.083}{-2.696}{7.9701}{mi}{27}\RG{45.599}{3.960}{7.9701}{sy}{0}\RG{52.185}{3.960}{7.9701}{rm}{49}\RG{59.948}{0.001}{10.9091}{rm}{61}\RG{71.464}{0.001}{10.9091}{mi}{80}\RG{78.467}{-2.696}{7.9701}{mi}{27}\RG{79.983}{3.960}{7.9701}{rm}{84}%
}\hskip.1em\allowbreak{} 。

(4)由

\[
\RB{187.282}{10.479}{5.712}{%
\RG{112.402}{4.503}{7.9701}{rm}{84}\RG{127.203}{4.503}{7.9701}{rm}{84}\RG{161.292}{4.503}{7.9701}{rm}{84}\RG{0.000}{-0.001}{10.9091}{rm}{40}\RG{4.242}{-0.001}{10.9091}{mi}{65}\RG{12.424}{-0.001}{10.9091}{mi}{80}\RG{25.051}{-0.001}{10.9091}{rm}{41}\RG{29.294}{-0.001}{10.9091}{mi}{101}\RG{40.785}{-0.001}{10.9091}{rm}{61}\RG{52.300}{-0.001}{10.9091}{mi}{65}\RG{60.482}{-0.001}{10.9091}{mi}{101}\RG{80.656}{-0.001}{10.9091}{mi}{59}\RG{107.307}{-0.001}{10.9091}{mi}{101}\RG{118.684}{-0.001}{10.9091}{mi}{80}\RG{133.485}{-0.001}{10.9091}{mi}{65}\RG{144.697}{-0.001}{10.9091}{rm}{61}\RG{156.212}{-0.001}{10.9091}{mi}{101}\RG{176.385}{-0.001}{10.9091}{mi}{65}\RG{184.567}{-0.001}{10.9091}{mi}{59}\RG{19.428}{-1.637}{7.9701}{mi}{27}\RG{34.373}{-1.637}{7.9701}{mi}{105}\RG{65.562}{-2.221}{7.9701}{mi}{27}\RG{70.688}{-2.221}{7.9701}{rm}{40}\RG{73.981}{-2.221}{7.9701}{mi}{105}\RG{76.864}{-2.221}{7.9701}{rm}{41}\RG{112.402}{-2.698}{7.9701}{mi}{105}\RG{125.688}{-2.698}{7.9701}{mi}{27}\RG{161.292}{-3.220}{7.9701}{mi}{27}\RG{166.417}{-3.220}{7.9701}{rm}{40}\RG{169.710}{-3.220}{7.9701}{mi}{105}\RG{172.593}{-3.220}{7.9701}{rm}{41}%
}
\]
分别得到所述列、行置换结论。\allowbreak{}\hskip.1em \RB{8.366}{7.993}{0.502}{%
\RG{0.000}{-0.002}{10.9091}{am}{3}%
}\hskip.1em\allowbreak{}

小结：恒等式\allowbreak{}\hskip.1em \RB{55.077}{7.949}{4.714}{%
\RG{0.000}{-0.002}{10.9091}{mi}{80}\RG{7.004}{-1.638}{7.9701}{mi}{27}\RG{12.627}{-0.002}{10.9091}{mi}{101}\RG{17.707}{-1.638}{7.9701}{mi}{105}\RG{24.168}{-0.002}{10.9091}{rm}{61}\RG{35.732}{-0.002}{10.9091}{mi}{101}\RG{40.812}{-2.222}{7.9701}{mi}{27}\RG{45.937}{-2.222}{7.9701}{rm}{40}\RG{49.231}{-2.222}{7.9701}{mi}{105}\RG{52.114}{-2.222}{7.9701}{rm}{41}%
}\hskip.1em\allowbreak{} 统一给出乘法、求逆及行列置换规则；映射\allowbreak{}\hskip.1em \RB{35.988}{7.949}{2.218}{%
\RG{0.000}{-0.002}{10.9091}{mi}{27}\RG{9.705}{-0.002}{10.9091}{sy}{55}\RG{9.705}{-0.002}{10.9091}{sy}{33}\RG{23.694}{-0.002}{10.9091}{mi}{80}\RG{30.698}{-1.638}{7.9701}{mi}{27}%
}\hskip.1em\allowbreak{} 因而是\allowbreak{}\hskip.1em \RB{12.067}{8.189}{2.218}{%
\RG{0.000}{-0.002}{10.9091}{mi}{83}\RG{6.689}{-1.638}{7.9701}{mi}{110}%
}\hskip.1em\allowbreak{} 到置换矩阵群的群同构。

题目来源：谢启鸿《复旦大学高等代数习题集（2023版）》，第80–81页。

% Original page 5.
\section*{2019A04}
求下列两类\allowbreak{}\hskip.1em \RB{6.729}{5.322}{0.620}{%
\RG{0.000}{-0.000}{10.9091}{mi}{110}%
}\hskip.1em\allowbreak{} 阶方阵：

(1)与全体\allowbreak{}\hskip.1em \RB{6.729}{5.327}{0.615}{%
\RG{0.000}{0.005}{10.9091}{mi}{110}%
}\hskip.1em\allowbreak{} 阶循环矩阵都乘法可交换的矩阵。(2)与全体\allowbreak{}\hskip.1em \RB{6.729}{5.326}{0.616}{%
\RG{0.000}{0.004}{10.9091}{mi}{110}%
}\hskip.1em\allowbreak{} 阶置换矩阵\allowbreak{}\hskip.1em \RB{64.615}{8.686}{3.223}{%
\RG{0.000}{0.004}{10.9091}{sy}{102}\RG{5.454}{0.004}{10.9091}{mi}{80}\RG{12.458}{-1.633}{7.9701}{mi}{27}\RG{21.112}{0.004}{10.9091}{rm}{58}\RG{27.172}{0.004}{10.9091}{mi}{27}\RG{36.828}{0.004}{10.9091}{sy}{50}\RG{47.131}{0.004}{10.9091}{mi}{83}\RG{53.820}{-1.633}{7.9701}{mi}{110}\RG{59.456}{0.004}{10.9091}{sy}{103}%
}\hskip.1em\allowbreak{} 都乘法可交换的矩阵。

解：以下矩阵的元素均属于数域\allowbreak{}\hskip.1em \RB{8.889}{7.971}{0.502}{%
\RG{0.000}{-0.002}{10.9091}{bb}{75}%
}\hskip.1em\allowbreak{} 。

(1)所求矩阵恰为全体循环矩阵。令

\[
\RB{80.000}{13.952}{19.764}{%
\RG{21.644}{13.016}{10.9092}{ex}{32}\RG{66.951}{13.016}{10.9092}{ex}{33}\RG{30.281}{1.641}{10.9091}{rm}{48}\RG{45.698}{1.641}{10.9091}{mi}{73}\RG{50.493}{0.005}{7.9701}{mi}{110}\RG{55.632}{0.005}{7.9701}{sy}{0}\RG{62.219}{0.005}{7.9701}{rm}{49}\RG{0.000}{-5.639}{10.9091}{mi}{74}\RG{10.129}{-5.639}{10.9091}{rm}{61}\RG{77.406}{-5.639}{10.9091}{mi}{58}\RG{30.281}{-13.534}{10.9091}{rm}{49}\RG{53.594}{-13.534}{10.9091}{rm}{48}%
}
\]
当\allowbreak{}\hskip.1em \RB{26.147}{7.762}{0.624}{%
\RG{0.000}{-0.004}{10.9091}{mi}{110}\RG{9.578}{-0.004}{10.9091}{rm}{61}\RG{21.076}{-0.004}{10.9091}{rm}{49}%
}\hskip.1em\allowbreak{} 时，取\allowbreak{}\hskip.1em \RB{34.967}{8.678}{3.231}{%
\RG{0.000}{-0.004}{10.9091}{mi}{74}\RG{10.129}{-0.004}{10.9091}{rm}{61}\RG{21.627}{-0.004}{10.9091}{rm}{40}\RG{25.860}{-0.004}{10.9091}{rm}{49}\RG{31.315}{-0.004}{10.9091}{rm}{41}%
}\hskip.1em\allowbreak{} 。由各次幂\allowbreak{}\hskip.1em \RB{68.295}{9.755}{2.609}{%
\RG{0.000}{-0.004}{10.9091}{mi}{73}\RG{4.795}{-1.640}{7.9701}{mi}{110}\RG{10.432}{-0.004}{10.9091}{mi}{59}\RG{15.265}{-0.004}{10.9091}{mi}{74}\RG{21.156}{-0.004}{10.9091}{mi}{59}\RG{25.989}{-0.004}{10.9091}{mi}{58}\RG{30.832}{-0.004}{10.9091}{mi}{58}\RG{35.665}{-0.004}{10.9091}{mi}{58}\RG{40.509}{-0.004}{10.9091}{mi}{59}\RG{45.352}{-0.004}{10.9091}{mi}{74}\RG{52.507}{3.955}{7.9701}{mi}{110}\RG{57.645}{3.955}{7.9701}{sy}{0}\RG{64.233}{3.955}{7.9701}{rm}{49}%
}\hskip.1em\allowbreak{} 中非零元的位置，循环矩阵恰可表示为

\[
\RB{239.853}{10.304}{2.605}{%
\RG{117.441}{4.504}{7.9701}{mi}{110}\RG{122.579}{4.504}{7.9701}{sy}{0}\RG{129.166}{4.504}{7.9701}{rm}{49}\RG{0.000}{0.000}{10.9091}{mi}{99}\RG{9.453}{0.000}{10.9091}{mi}{73}\RG{22.309}{0.000}{10.9091}{rm}{43}\RG{33.218}{0.000}{10.9091}{mi}{99}\RG{42.671}{0.000}{10.9091}{mi}{74}\RG{52.194}{0.000}{10.9091}{rm}{43}\RG{63.104}{0.000}{10.9091}{sy}{1}\RG{67.936}{0.000}{10.9091}{sy}{1}\RG{72.780}{0.000}{10.9091}{sy}{1}\RG{78.256}{0.000}{10.9091}{rm}{43}\RG{89.165}{0.000}{10.9091}{mi}{99}\RG{110.343}{0.000}{10.9091}{mi}{74}\RG{133.898}{0.000}{10.9091}{mi}{59}\RG{160.549}{0.000}{10.9091}{mi}{99}\RG{170.018}{0.000}{10.9091}{mi}{59}\RG{174.851}{0.000}{10.9091}{mi}{58}\RG{179.695}{0.000}{10.9091}{mi}{58}\RG{184.538}{0.000}{10.9091}{mi}{58}\RG{189.371}{0.000}{10.9091}{mi}{59}\RG{194.215}{0.000}{10.9091}{mi}{99}\RG{218.470}{0.000}{10.9091}{sy}{50}\RG{228.773}{0.000}{10.9091}{bb}{75}\RG{237.258}{0.000}{10.9091}{mi}{58}\RG{4.721}{-1.636}{7.9701}{rm}{48}\RG{14.249}{-1.636}{7.9701}{mi}{110}\RG{37.939}{-1.636}{7.9701}{rm}{49}\RG{93.886}{-1.636}{7.9701}{mi}{110}\RG{99.024}{-1.636}{7.9701}{sy}{0}\RG{105.610}{-1.636}{7.9701}{rm}{49}\RG{165.286}{-1.636}{7.9701}{rm}{48}\RG{198.982}{-1.636}{7.9701}{mi}{110}\RG{204.121}{-1.636}{7.9701}{sy}{0}\RG{210.707}{-1.636}{7.9701}{rm}{49}%
}
\]
例如其第一行为\allowbreak{}\hskip.1em \RB{77.071}{8.686}{3.223}{%
\RG{0.000}{0.004}{10.9091}{rm}{40}\RG{4.243}{0.004}{10.9091}{mi}{99}\RG{8.964}{-1.633}{7.9701}{rm}{48}\RG{13.696}{0.004}{10.9091}{mi}{59}\RG{18.529}{0.004}{10.9091}{mi}{99}\RG{23.265}{-1.633}{7.9701}{rm}{49}\RG{27.998}{0.004}{10.9091}{mi}{59}\RG{32.831}{0.004}{10.9091}{mi}{58}\RG{37.675}{0.004}{10.9091}{mi}{58}\RG{42.507}{0.004}{10.9091}{mi}{58}\RG{47.351}{0.004}{10.9091}{mi}{59}\RG{52.195}{0.004}{10.9091}{mi}{99}\RG{56.961}{-1.633}{7.9701}{mi}{110}\RG{62.100}{-1.633}{7.9701}{sy}{0}\RG{68.687}{-1.633}{7.9701}{rm}{49}\RG{73.419}{0.004}{10.9091}{rm}{41}%
}\hskip.1em\allowbreak{} ，以后每一行由上一行向右循环移动一位得到。因此，同为\allowbreak{}\hskip.1em \RB{7.406}{7.955}{0.736}{%
\RG{0.000}{0.004}{10.9091}{mi}{74}%
}\hskip.1em\allowbreak{} 的多项式的任意两个循环矩阵可交换。

反之，若\allowbreak{}\hskip.1em \RB{8.366}{8.313}{0.500}{%
\RG{0.000}{0.002}{10.9091}{mi}{65}%
}\hskip.1em\allowbreak{} 与全体循环矩阵可交换，则\allowbreak{}\hskip.1em \RB{45.282}{8.313}{0.738}{%
\RG{0.000}{0.002}{10.9091}{mi}{65}\RG{8.182}{0.002}{10.9091}{mi}{74}\RG{18.310}{0.002}{10.9091}{rm}{61}\RG{29.825}{0.002}{10.9091}{mi}{74}\RG{36.916}{0.002}{10.9091}{mi}{65}%
}\hskip.1em\allowbreak{} 。由于

\[
\RB{88.511}{10.300}{2.609}{%
\RG{62.911}{4.500}{7.9701}{mi}{110}\RG{68.050}{4.500}{7.9701}{sy}{0}\RG{74.637}{4.500}{7.9701}{rm}{49}\RG{0.000}{-0.005}{10.9091}{mi}{101}\RG{9.812}{-0.004}{10.9091}{mi}{59}\RG{14.645}{-0.004}{10.9091}{mi}{74}\RG{21.736}{-0.004}{10.9091}{mi}{101}\RG{31.570}{-0.004}{10.9091}{mi}{59}\RG{36.403}{-0.004}{10.9091}{mi}{58}\RG{41.247}{-0.004}{10.9091}{mi}{58}\RG{46.090}{-0.004}{10.9091}{mi}{58}\RG{50.923}{-0.004}{10.9091}{mi}{59}\RG{55.767}{-0.004}{10.9091}{mi}{74}\RG{79.369}{-0.005}{10.9091}{mi}{101}\RG{5.079}{-1.641}{7.9701}{rm}{49}\RG{26.838}{-1.641}{7.9701}{rm}{49}\RG{84.449}{-1.641}{7.9701}{rm}{49}%
}
\]
是标准基的一个排列，存在次数不超过\allowbreak{}\hskip.1em \RB{24.952}{7.765}{0.621}{%
\RG{0.000}{-0.001}{10.9091}{mi}{110}\RG{8.972}{-0.001}{10.9091}{sy}{0}\RG{19.881}{-0.001}{10.9091}{rm}{49}%
}\hskip.1em\allowbreak{} 的多项式\allowbreak{}\hskip.1em \RB{5.846}{5.321}{2.617}{%
\RG{0.000}{-0.001}{10.9091}{mi}{112}%
}\hskip.1em\allowbreak{} ，使得\allowbreak{}\hskip.1em \RB{62.754}{8.681}{3.228}{%
\RG{0.000}{-0.001}{10.9091}{mi}{65}\RG{8.182}{-0.001}{10.9091}{mi}{101}\RG{13.262}{-1.637}{7.9701}{rm}{49}\RG{21.024}{-0.001}{10.9091}{rm}{61}\RG{32.540}{-0.001}{10.9091}{mi}{112}\RG{38.028}{-0.001}{10.9091}{rm}{40}\RG{42.271}{-0.001}{10.9091}{mi}{74}\RG{49.369}{-0.001}{10.9091}{rm}{41}\RG{53.612}{-0.001}{10.9091}{mi}{101}\RG{58.691}{-1.637}{7.9701}{rm}{49}%
}\hskip.1em\allowbreak{} 。对\allowbreak{}\hskip.1em \RB{64.723}{7.765}{1.996}{%
\RG{0.000}{-0.001}{10.9091}{rm}{48}\RG{8.485}{-0.001}{10.9091}{sy}{20}\RG{20.000}{-0.001}{10.9091}{mi}{114}\RG{28.255}{-0.001}{10.9091}{sy}{20}\RG{39.770}{-0.001}{10.9091}{mi}{110}\RG{48.743}{-0.001}{10.9091}{sy}{0}\RG{59.653}{-0.001}{10.9091}{rm}{49}%
}\hskip.1em\allowbreak{} ，

\[
\RB{190.600}{8.685}{3.224}{%
\RG{15.280}{4.507}{7.9701}{mi}{114}\RG{51.290}{4.507}{7.9701}{mi}{114}\RG{95.484}{4.507}{7.9701}{mi}{114}\RG{173.638}{4.507}{7.9701}{mi}{114}\RG{0.000}{0.003}{10.9091}{mi}{65}\RG{8.182}{0.003}{10.9091}{mi}{74}\RG{19.835}{0.003}{10.9091}{mi}{101}\RG{32.677}{0.003}{10.9091}{rm}{61}\RG{44.192}{0.003}{10.9091}{mi}{74}\RG{55.845}{0.003}{10.9091}{mi}{65}\RG{64.027}{0.003}{10.9091}{mi}{101}\RG{76.870}{0.003}{10.9091}{rm}{61}\RG{88.385}{0.003}{10.9091}{mi}{74}\RG{100.039}{0.003}{10.9091}{mi}{112}\RG{105.527}{0.003}{10.9091}{rm}{40}\RG{109.770}{0.003}{10.9091}{mi}{74}\RG{116.868}{0.003}{10.9091}{rm}{41}\RG{121.111}{0.003}{10.9091}{mi}{101}\RG{133.953}{0.003}{10.9091}{rm}{61}\RG{145.468}{0.003}{10.9091}{mi}{112}\RG{150.956}{0.003}{10.9091}{rm}{40}\RG{155.199}{0.003}{10.9091}{mi}{74}\RG{162.297}{0.003}{10.9091}{rm}{41}\RG{166.540}{0.003}{10.9091}{mi}{74}\RG{178.193}{0.003}{10.9091}{mi}{101}\RG{188.006}{0.003}{10.9091}{mi}{58}\RG{24.914}{-1.634}{7.9701}{rm}{49}\RG{69.108}{-1.634}{7.9701}{rm}{49}\RG{126.190}{-1.634}{7.9701}{rm}{49}\RG{183.273}{-1.634}{7.9701}{rm}{49}%
}
\]
\allowbreak{}\hskip.1em \RB{8.366}{8.307}{0.504}{%
\RG{0.000}{-0.004}{10.9091}{mi}{65}%
}\hskip.1em\allowbreak{} 与\allowbreak{}\hskip.1em \RB{20.483}{8.678}{3.231}{%
\RG{0.000}{-0.004}{10.9091}{mi}{112}\RG{5.489}{-0.004}{10.9091}{rm}{40}\RG{9.731}{-0.004}{10.9091}{mi}{74}\RG{16.830}{-0.004}{10.9091}{rm}{41}%
}\hskip.1em\allowbreak{} 在一组基上作用相同，故\allowbreak{}\hskip.1em \RB{43.209}{8.678}{3.231}{%
\RG{0.000}{-0.004}{10.9091}{mi}{65}\RG{11.212}{-0.004}{10.9091}{rm}{61}\RG{22.727}{-0.004}{10.9091}{mi}{112}\RG{28.216}{-0.004}{10.9091}{rm}{40}\RG{32.458}{-0.004}{10.9091}{mi}{74}\RG{39.556}{-0.004}{10.9091}{rm}{41}%
}\hskip.1em\allowbreak{} ，即\allowbreak{}\hskip.1em \RB{8.366}{8.307}{0.504}{%
\RG{0.000}{-0.004}{10.9091}{mi}{65}%
}\hskip.1em\allowbreak{} 是循环矩阵。

(2)记\allowbreak{}\hskip.1em \RB{54.270}{8.682}{3.227}{%
\RG{0.000}{-0.000}{10.9091}{mi}{69}\RG{11.712}{-0.000}{10.9091}{rm}{61}\RG{23.210}{-0.000}{10.9091}{rm}{40}\RG{27.443}{-0.000}{10.9091}{rm}{49}\RG{32.898}{-0.000}{10.9091}{rm}{41}\RG{37.167}{-1.637}{7.9701}{mi}{110}\RG{42.305}{-1.637}{7.9701}{sy}{2}\RG{48.892}{-1.637}{7.9701}{mi}{110}%
}\hskip.1em\allowbreak{} 为全\allowbreak{}\hskip.1em \RB{5.071}{7.766}{0.500}{%
\RG{0.000}{-0.000}{10.9091}{rm}{49}%
}\hskip.1em\allowbreak{} 矩阵。所求矩阵恰为

\[
\RB{132.000}{8.314}{2.602}{%
\RG{0.000}{0.003}{10.9091}{mi}{65}\RG{11.212}{0.003}{10.9091}{rm}{61}\RG{22.727}{0.003}{10.9091}{mi}{97}\RG{28.487}{0.003}{10.9091}{mi}{73}\RG{41.349}{0.003}{10.9091}{rm}{43}\RG{52.259}{0.003}{10.9091}{mi}{98}\RG{56.939}{0.003}{10.9091}{mi}{69}\RG{65.612}{0.003}{10.9091}{mi}{59}\RG{92.274}{0.003}{10.9091}{mi}{97}\RG{98.034}{0.003}{10.9091}{mi}{59}\RG{102.867}{0.003}{10.9091}{mi}{98}\RG{110.617}{0.003}{10.9091}{sy}{50}\RG{120.920}{0.003}{10.9091}{bb}{75}\RG{129.405}{0.003}{10.9091}{mi}{58}\RG{33.289}{-1.633}{7.9701}{mi}{110}%
}
\]
事实上，由\allowbreak{}\hskip.1em \RB{48.878}{9.932}{3.280}{%
\RG{0.000}{-0.003}{10.9091}{mi}{80}\RG{7.004}{-2.700}{7.9701}{mi}{27}\RG{8.519}{3.956}{7.9701}{sy}{0}\RG{15.106}{3.956}{7.9701}{rm}{49}\RG{22.868}{-0.003}{10.9091}{rm}{61}\RG{34.383}{-0.003}{10.9091}{mi}{80}\RG{41.387}{-2.700}{7.9701}{mi}{27}\RG{42.902}{3.956}{7.9701}{rm}{84}%
}\hskip.1em\allowbreak{} ，条件\allowbreak{}\hskip.1em \RB{56.348}{8.308}{2.219}{%
\RG{0.000}{-0.003}{10.9091}{mi}{65}\RG{8.182}{-0.003}{10.9091}{mi}{80}\RG{15.186}{-1.639}{7.9701}{mi}{27}\RG{23.840}{-0.003}{10.9091}{rm}{61}\RG{35.355}{-0.003}{10.9091}{mi}{80}\RG{42.358}{-1.639}{7.9701}{mi}{27}\RG{47.982}{-0.003}{10.9091}{mi}{65}%
}\hskip.1em\allowbreak{} 等价于

\[
\RB{60.933}{10.481}{3.276}{%
\RG{8.519}{4.505}{7.9701}{rm}{84}\RG{0.000}{0.001}{10.9091}{mi}{80}\RG{14.801}{0.001}{10.9091}{mi}{65}\RG{22.983}{0.001}{10.9091}{mi}{80}\RG{38.640}{0.001}{10.9091}{rm}{61}\RG{50.156}{0.001}{10.9091}{mi}{65}\RG{58.338}{0.001}{10.9091}{mi}{58}\RG{29.987}{-1.636}{7.9701}{mi}{27}\RG{7.004}{-2.696}{7.9701}{mi}{27}%
}
\]
比较第\allowbreak{}\hskip.1em \RB{21.619}{8.686}{3.223}{%
\RG{0.000}{0.004}{10.9091}{rm}{40}\RG{4.242}{0.004}{10.9091}{mi}{105}\RG{7.995}{0.004}{10.9091}{mi}{59}\RG{12.828}{0.004}{10.9091}{mi}{106}\RG{17.966}{0.004}{10.9091}{rm}{41}%
}\hskip.1em\allowbreak{} 个元素，得到

\[
\RB{131.568}{8.681}{4.714}{%
\RG{0.000}{-0.001}{10.9091}{mi}{97}\RG{41.838}{-0.001}{10.9091}{rm}{61}\RG{53.353}{-0.001}{10.9091}{mi}{97}\RG{88.204}{-0.001}{10.9091}{rm}{40}\RG{92.446}{-0.001}{10.9091}{mi}{27}\RG{102.101}{-0.001}{10.9091}{sy}{50}\RG{112.404}{-0.001}{10.9091}{mi}{83}\RG{124.730}{-0.001}{10.9091}{rm}{41}\RG{128.973}{-0.001}{10.9091}{mi}{58}\RG{59.120}{-1.637}{7.9701}{mi}{105}\RG{61.997}{-1.637}{7.9701}{mi}{106}\RG{119.094}{-1.637}{7.9701}{mi}{110}\RG{5.766}{-2.221}{7.9701}{mi}{27}\RG{10.892}{-2.221}{7.9701}{rm}{40}\RG{14.185}{-2.221}{7.9701}{mi}{105}\RG{17.068}{-2.221}{7.9701}{rm}{41}\RG{20.361}{-2.221}{7.9701}{mi}{59}\RG{22.712}{-2.221}{7.9701}{mi}{27}\RG{27.839}{-2.221}{7.9701}{rm}{40}\RG{31.132}{-2.221}{7.9701}{mi}{106}\RG{35.016}{-2.221}{7.9701}{rm}{41}%
}
\]
任意两个对角位置可通过某个置换互相变换，故所有对角元相等，记为\allowbreak{}\hskip.1em \RB{6.129}{8.068}{0.623}{%
\RG{0.000}{-0.003}{10.9091}{mi}{100}%
}\hskip.1em\allowbreak{} 。当\allowbreak{}\hskip.1em \RB{28.199}{7.763}{1.998}{%
\RG{0.000}{-0.003}{10.9091}{mi}{110}\RG{10.432}{-0.003}{10.9091}{sy}{21}\RG{22.801}{-0.003}{10.9091}{rm}{50}%
}\hskip.1em\allowbreak{} 时，对任意\allowbreak{}\hskip.1em \RB{23.135}{8.313}{2.854}{%
\RG{0.000}{0.002}{10.9091}{mi}{105}\RG{6.788}{0.002}{10.9091}{sy}{54}\RG{6.788}{0.002}{10.9091}{rm}{61}\RG{18.304}{0.002}{10.9091}{mi}{106}%
}\hskip.1em\allowbreak{} 、\allowbreak{}\hskip.1em \RB{25.410}{8.313}{2.854}{%
\RG{0.000}{0.002}{10.9091}{mi}{107}\RG{9.053}{0.002}{10.9091}{sy}{54}\RG{9.053}{0.002}{10.9091}{rm}{61}\RG{20.568}{0.002}{10.9091}{mi}{96}%
}\hskip.1em\allowbreak{} ，可将\allowbreak{}\hskip.1em \RB{26.770}{8.073}{0.618}{%
\RG{0.000}{0.002}{10.9091}{mi}{105}\RG{6.789}{0.002}{10.9091}{sy}{55}\RG{6.789}{0.002}{10.9091}{sy}{33}\RG{20.728}{0.002}{10.9091}{mi}{107}%
}\hskip.1em\allowbreak{} 、\allowbreak{}\hskip.1em \RB{26.928}{8.193}{2.734}{%
\RG{0.000}{0.002}{10.9091}{mi}{106}\RG{8.147}{0.002}{10.9091}{sy}{55}\RG{8.147}{0.002}{10.9091}{sy}{33}\RG{22.086}{0.002}{10.9091}{mi}{96}%
}\hskip.1em\allowbreak{} 扩充为一个置换，故所有非对角元也相等，记为\allowbreak{}\hskip.1em \RB{5.027}{8.073}{0.618}{%
\RG{0.000}{0.002}{10.9091}{mi}{98}%
}\hskip.1em\allowbreak{} 。于是

\[
\RB{94.618}{8.678}{3.231}{%
\RG{0.000}{-0.004}{10.9091}{mi}{65}\RG{11.212}{-0.004}{10.9091}{rm}{61}\RG{22.710}{-0.004}{10.9091}{rm}{40}\RG{26.969}{-0.004}{10.9091}{mi}{100}\RG{35.072}{-0.004}{10.9091}{sy}{0}\RG{45.981}{-0.004}{10.9091}{mi}{98}\RG{50.663}{-0.004}{10.9091}{rm}{41}\RG{54.905}{-0.004}{10.9091}{mi}{73}\RG{67.762}{-0.004}{10.9091}{rm}{43}\RG{78.671}{-0.004}{10.9091}{mi}{98}\RG{83.351}{-0.004}{10.9091}{mi}{69}\RG{92.023}{-0.004}{10.9091}{mi}{58}\RG{59.701}{-1.640}{7.9701}{mi}{110}%
}
\]
当\allowbreak{}\hskip.1em \RB{26.147}{7.766}{0.620}{%
\RG{0.000}{0.000}{10.9091}{mi}{110}\RG{9.578}{0.000}{10.9091}{rm}{61}\RG{21.076}{0.000}{10.9091}{rm}{49}%
}\hskip.1em\allowbreak{} 时，上述所求形式仍包含全部一阶方阵。

反之，置换矩阵只改变\allowbreak{}\hskip.1em \RB{8.832}{7.921}{0.500}{%
\RG{0.000}{0.003}{10.9091}{mi}{69}%
}\hskip.1em\allowbreak{} 的行或列的顺序，所以

\[
\RB{82.977}{7.947}{2.220}{%
\RG{0.000}{-0.004}{10.9091}{mi}{69}\RG{8.673}{-0.004}{10.9091}{mi}{80}\RG{24.340}{-0.004}{10.9091}{rm}{61}\RG{35.855}{-0.004}{10.9091}{mi}{80}\RG{48.482}{-0.004}{10.9091}{mi}{69}\RG{60.195}{-0.004}{10.9091}{rm}{61}\RG{71.710}{-0.004}{10.9091}{mi}{69}\RG{80.382}{-0.004}{10.9091}{mi}{58}\RG{15.686}{-1.640}{7.9701}{mi}{27}\RG{42.859}{-1.640}{7.9701}{mi}{27}%
}
\]
结合\allowbreak{}\hskip.1em \RB{60.405}{7.951}{2.216}{%
\RG{0.000}{0.000}{10.9091}{mi}{73}\RG{4.795}{-1.636}{7.9701}{mi}{110}\RG{10.432}{0.000}{10.9091}{mi}{80}\RG{17.435}{-1.636}{7.9701}{mi}{27}\RG{26.089}{0.000}{10.9091}{rm}{61}\RG{37.604}{0.000}{10.9091}{mi}{80}\RG{44.608}{-1.636}{7.9701}{mi}{27}\RG{50.231}{0.000}{10.9091}{mi}{73}\RG{55.027}{-1.636}{7.9701}{mi}{110}%
}\hskip.1em\allowbreak{} ，可知任意\allowbreak{}\hskip.1em \RB{43.044}{8.071}{2.216}{%
\RG{0.000}{0.000}{10.9091}{mi}{97}\RG{5.760}{0.000}{10.9091}{mi}{73}\RG{10.562}{-1.636}{7.9701}{mi}{110}\RG{18.623}{0.000}{10.9091}{rm}{43}\RG{29.532}{0.000}{10.9091}{mi}{98}\RG{34.212}{0.000}{10.9091}{mi}{69}%
}\hskip.1em\allowbreak{} 均满足要求。

小结：与循环移位可交换，迫使矩阵由一个循环向量上的作用决定；与所有置换可交换，则迫使其元素在对角和非对角两类位置上分别相等。

题目来源：谢启鸿《复旦大学高等代数习题集（2023版）》，第81页。

% Original page 6.
\section*{2019A05}
求下列\allowbreak{}\hskip.1em \RB{6.729}{5.322}{0.620}{%
\RG{0.000}{-0.000}{10.9091}{mi}{110}%
}\hskip.1em\allowbreak{} 阶矩阵的行列式，其中第四行的省略号表示同一带状结构继续平移，其余空白处均为零：

\[
\RB{346.769}{41.749}{70.504}{%
\RG{26.318}{40.823}{10.9092}{ex}{48}\RG{332.811}{40.823}{10.9092}{ex}{49}\RG{41.630}{34.312}{7.9701}{rm}{50}\RG{147.033}{34.312}{7.9701}{rm}{50}\RG{35.863}{30.353}{10.9091}{mi}{97}\RG{48.786}{30.353}{10.9091}{rm}{43}\RG{59.695}{30.353}{10.9091}{mi}{98}\RG{64.376}{30.353}{10.9091}{mi}{99}\RG{90.455}{30.353}{10.9091}{rm}{50}\RG{95.910}{30.353}{10.9091}{mi}{97}\RG{101.670}{30.353}{10.9091}{mi}{98}\RG{142.351}{30.353}{10.9091}{mi}{98}\RG{26.318}{21.623}{10.9092}{ex}{66}\RG{332.811}{21.623}{10.9092}{ex}{67}\RG{84.828}{19.137}{7.9701}{rm}{50}\RG{184.312}{19.137}{7.9701}{rm}{50}\RG{44.510}{15.178}{10.9091}{rm}{50}\RG{49.964}{15.178}{10.9091}{mi}{97}\RG{55.724}{15.178}{10.9091}{mi}{99}\RG{79.037}{15.178}{10.9091}{mi}{97}\RG{91.985}{15.178}{10.9091}{rm}{43}\RG{102.883}{15.178}{10.9091}{rm}{50}\RG{108.348}{15.178}{10.9091}{mi}{98}\RG{113.028}{15.178}{10.9091}{mi}{99}\RG{139.107}{15.178}{10.9091}{rm}{50}\RG{144.562}{15.178}{10.9091}{mi}{97}\RG{150.322}{15.178}{10.9091}{mi}{98}\RG{179.613}{15.178}{10.9091}{mi}{98}\RG{26.318}{15.077}{10.9092}{ex}{66}\RG{332.811}{15.077}{10.9092}{ex}{67}\RG{26.318}{8.532}{10.9092}{ex}{66}\RG{332.811}{8.532}{10.9092}{ex}{67}\RG{52.475}{3.962}{7.9701}{rm}{50}\RG{133.480}{3.962}{7.9701}{rm}{50}\RG{221.591}{3.962}{7.9701}{rm}{50}\RG{26.318}{1.986}{10.9092}{ex}{66}\RG{332.811}{1.986}{10.9092}{ex}{67}\RG{47.754}{0.003}{10.9091}{mi}{99}\RG{90.435}{0.003}{10.9091}{rm}{50}\RG{95.890}{0.003}{10.9091}{mi}{97}\RG{101.650}{0.003}{10.9091}{mi}{99}\RG{127.690}{0.003}{10.9091}{mi}{97}\RG{140.637}{0.003}{10.9091}{rm}{43}\RG{151.535}{0.003}{10.9091}{rm}{50}\RG{157.000}{0.003}{10.9091}{mi}{98}\RG{161.680}{0.003}{10.9091}{mi}{99}\RG{176.385}{0.003}{10.9091}{rm}{50}\RG{181.841}{0.003}{10.9091}{mi}{97}\RG{187.601}{0.003}{10.9091}{mi}{98}\RG{216.892}{0.003}{10.9091}{mi}{98}\RG{26.318}{-4.559}{10.9092}{ex}{66}\RG{332.811}{-4.559}{10.9092}{ex}{67}\RG{26.318}{-11.105}{10.9092}{ex}{66}\RG{332.811}{-11.105}{10.9092}{ex}{67}\RT{178.575}{-12.679}{10.9091}{.}\RT{182.821}{-15.668}{10.9091}{.}\RG{0.000}{-17.105}{10.9091}{mi}{77}\RG{14.803}{-17.105}{10.9091}{rm}{61}\RG{344.175}{-17.105}{10.9091}{mi}{58}\RG{26.318}{-17.650}{10.9092}{ex}{66}\RG{332.811}{-17.650}{10.9092}{ex}{67}\RT{187.066}{-18.657}{10.9091}{.}\RG{26.318}{-24.196}{10.9092}{ex}{66}\RG{332.811}{-24.196}{10.9092}{ex}{67}\RG{26.318}{-30.741}{10.9092}{ex}{66}\RG{332.811}{-30.741}{10.9092}{ex}{67}\RG{147.053}{-30.869}{7.9701}{rm}{50}\RG{208.038}{-30.869}{7.9701}{rm}{50}\RG{316.168}{-30.869}{7.9701}{rm}{50}\RG{142.332}{-34.828}{10.9091}{mi}{99}\RG{176.366}{-34.828}{10.9091}{rm}{50}\RG{181.820}{-34.828}{10.9091}{mi}{97}\RG{187.581}{-34.828}{10.9091}{mi}{99}\RG{202.253}{-34.828}{10.9091}{mi}{97}\RG{215.194}{-34.828}{10.9091}{rm}{43}\RG{226.092}{-34.828}{10.9091}{rm}{50}\RG{231.558}{-34.828}{10.9091}{mi}{98}\RG{236.238}{-34.828}{10.9091}{mi}{99}\RG{262.316}{-34.828}{10.9091}{rm}{50}\RG{267.771}{-34.828}{10.9091}{mi}{97}\RG{273.531}{-34.828}{10.9091}{mi}{98}\RG{311.487}{-34.828}{10.9091}{mi}{98}\RG{26.318}{-37.287}{10.9092}{ex}{66}\RG{332.811}{-37.287}{10.9092}{ex}{67}\RG{26.318}{-43.832}{10.9092}{ex}{66}\RG{332.811}{-43.832}{10.9092}{ex}{67}\RG{184.332}{-46.045}{7.9701}{rm}{50}\RG{256.690}{-46.045}{7.9701}{rm}{50}\RG{179.610}{-50.003}{10.9091}{mi}{99}\RG{213.645}{-50.003}{10.9091}{rm}{50}\RG{219.099}{-50.003}{10.9091}{mi}{97}\RG{224.859}{-50.003}{10.9091}{mi}{99}\RG{250.899}{-50.003}{10.9091}{mi}{97}\RG{263.846}{-50.003}{10.9091}{rm}{43}\RG{274.744}{-50.003}{10.9091}{rm}{50}\RG{280.210}{-50.003}{10.9091}{mi}{98}\RG{284.890}{-50.003}{10.9091}{mi}{99}\RG{308.242}{-50.003}{10.9091}{rm}{50}\RG{313.697}{-50.003}{10.9091}{mi}{97}\RG{319.457}{-50.003}{10.9091}{mi}{98}\RG{26.318}{-50.814}{10.9092}{ex}{64}\RG{332.811}{-50.814}{10.9092}{ex}{65}\RG{221.610}{-61.221}{7.9701}{rm}{50}\RG{305.343}{-61.221}{7.9701}{rm}{50}\RG{216.889}{-65.179}{10.9091}{mi}{99}\RG{262.297}{-65.179}{10.9091}{rm}{50}\RG{267.751}{-65.179}{10.9091}{mi}{97}\RG{273.511}{-65.179}{10.9091}{mi}{99}\RG{299.552}{-65.179}{10.9091}{mi}{97}\RG{312.499}{-65.179}{10.9091}{rm}{43}\RG{323.408}{-65.179}{10.9091}{mi}{98}\RG{328.088}{-65.179}{10.9091}{mi}{99}%
}
\]
解：引入三对角矩阵

\[
\RB{139.649}{25.380}{54.145}{%
\RG{26.557}{24.454}{10.9092}{ex}{48}\RG{125.690}{24.454}{10.9092}{ex}{49}\RG{36.102}{15.174}{10.9091}{mi}{97}\RG{55.848}{15.174}{10.9091}{mi}{98}\RG{26.557}{5.254}{10.9092}{ex}{66}\RG{125.690}{5.254}{10.9092}{ex}{67}\RG{36.625}{-0.001}{10.9091}{mi}{99}\RG{55.302}{-0.001}{10.9091}{mi}{97}\RG{78.538}{-0.001}{10.9091}{mi}{98}\RG{26.557}{-1.291}{10.9092}{ex}{66}\RG{125.690}{-1.291}{10.9092}{ex}{67}\RG{26.557}{-7.837}{10.9092}{ex}{66}\RG{125.690}{-7.837}{10.9092}{ex}{67}\RT{52.438}{-12.684}{10.9091}{.}\RT{75.136}{-12.684}{10.9091}{.}\RT{97.833}{-12.684}{10.9091}{.}\RG{26.557}{-14.382}{10.9092}{ex}{66}\RG{125.690}{-14.382}{10.9092}{ex}{67}\RT{56.682}{-15.673}{10.9091}{.}\RT{79.381}{-15.673}{10.9091}{.}\RT{102.078}{-15.673}{10.9091}{.}\RG{0.000}{-17.110}{10.9091}{mi}{84}\RG{15.042}{-17.110}{10.9091}{rm}{61}\RG{137.054}{-17.109}{10.9091}{mi}{58}\RG{6.375}{-18.746}{7.9701}{mi}{110}\RT{60.928}{-18.662}{10.9091}{.}\RT{83.625}{-18.662}{10.9091}{.}\RT{106.322}{-18.662}{10.9091}{.}\RG{26.557}{-20.928}{10.9092}{ex}{66}\RG{125.690}{-20.928}{10.9092}{ex}{67}\RG{26.557}{-27.473}{10.9092}{ex}{66}\RG{125.690}{-27.473}{10.9092}{ex}{67}\RG{26.557}{-34.455}{10.9092}{ex}{64}\RG{125.690}{-34.455}{10.9092}{ex}{65}\RG{78.536}{-34.833}{10.9091}{mi}{99}\RG{100.692}{-34.833}{10.9091}{mi}{97}\RG{120.449}{-34.833}{10.9091}{mi}{98}\RG{101.233}{-50.008}{10.9091}{mi}{99}\RG{119.909}{-50.008}{10.9091}{mi}{97}%
}
\]
直接计算\allowbreak{}\hskip.1em \RB{12.192}{9.761}{3.275}{%
\RG{0.000}{0.002}{10.9091}{mi}{84}\RG{6.375}{-2.695}{7.9701}{mi}{110}\RG{7.890}{3.961}{7.9701}{rm}{50}%
}\hskip.1em\allowbreak{} ：两端的对角元为\allowbreak{}\hskip.1em \RB{31.875}{9.761}{1.403}{%
\RG{0.000}{0.002}{10.9091}{mi}{97}\RG{5.767}{3.961}{7.9701}{rm}{50}\RG{12.009}{0.002}{10.9091}{rm}{43}\RG{22.004}{0.002}{10.9091}{mi}{98}\RG{26.684}{0.002}{10.9091}{mi}{99}%
}\hskip.1em\allowbreak{} ，其余对角元为\allowbreak{}\hskip.1em \RB{37.330}{9.761}{1.403}{%
\RG{0.000}{0.002}{10.9091}{mi}{97}\RG{5.767}{3.961}{7.9701}{rm}{50}\RG{12.009}{0.002}{10.9091}{rm}{43}\RG{21.991}{0.002}{10.9091}{rm}{50}\RG{27.459}{0.002}{10.9091}{mi}{98}\RG{32.139}{0.002}{10.9091}{mi}{99}%
}\hskip.1em\allowbreak{} ；第一条上、下副对角线分别为\allowbreak{}\hskip.1em \RB{16.242}{8.073}{0.618}{%
\RG{0.000}{0.002}{10.9091}{rm}{50}\RG{5.455}{0.002}{10.9091}{mi}{97}\RG{11.215}{0.002}{10.9091}{mi}{98}%
}\hskip.1em\allowbreak{} 、\allowbreak{}\hskip.1em \RB{16.405}{7.763}{0.623}{%
\RG{0.000}{-0.003}{10.9091}{rm}{50}\RG{5.454}{-0.003}{10.9091}{mi}{97}\RG{11.214}{-0.003}{10.9091}{mi}{99}%
}\hskip.1em\allowbreak{} ，第二条上、下副对角线分别为\allowbreak{}\hskip.1em \RB{8.983}{9.756}{0.623}{%
\RG{0.000}{-0.003}{10.9091}{mi}{98}\RG{4.681}{3.956}{7.9701}{rm}{50}%
}\hskip.1em\allowbreak{} 、\allowbreak{}\hskip.1em \RB{9.023}{9.756}{0.623}{%
\RG{0.000}{-0.003}{10.9091}{mi}{99}\RG{4.721}{3.956}{7.9701}{rm}{50}%
}\hskip.1em\allowbreak{} ，其余元素为零。因此

\[
\RB{153.689}{10.308}{3.274}{%
\RG{34.208}{4.508}{7.9701}{rm}{50}\RG{146.362}{4.508}{7.9701}{rm}{50}\RG{0.000}{0.003}{10.9091}{mi}{77}\RG{14.803}{0.003}{10.9091}{rm}{61}\RG{26.318}{0.003}{10.9091}{mi}{84}\RG{38.941}{0.003}{10.9091}{mi}{59}\RG{65.608}{0.003}{10.9091}{rm}{100}\RG{71.662}{0.003}{10.9091}{rm}{101}\RG{76.506}{0.003}{10.9091}{rm}{116}\RG{82.578}{0.003}{10.9091}{mi}{77}\RG{97.381}{0.003}{10.9091}{rm}{61}\RG{108.879}{0.003}{10.9091}{rm}{40}\RG{113.112}{0.003}{10.9091}{rm}{100}\RG{119.166}{0.003}{10.9091}{rm}{101}\RG{124.010}{0.003}{10.9091}{rm}{116}\RG{130.108}{0.003}{10.9091}{mi}{84}\RG{142.119}{0.003}{10.9091}{rm}{41}\RG{151.094}{0.003}{10.9091}{mi}{58}\RG{136.483}{-1.633}{7.9701}{mi}{110}\RG{32.693}{-2.694}{7.9701}{mi}{110}%
}
\]
记\allowbreak{}\hskip.1em \RB{54.582}{8.069}{2.219}{%
\RG{0.000}{-0.002}{10.9091}{mi}{100}\RG{5.678}{-1.639}{7.9701}{mi}{110}\RG{14.344}{-0.002}{10.9091}{rm}{61}\RG{25.842}{-0.002}{10.9091}{rm}{100}\RG{31.897}{-0.002}{10.9091}{rm}{101}\RG{36.741}{-0.002}{10.9091}{rm}{116}\RG{42.829}{-0.002}{10.9091}{mi}{84}\RG{49.204}{-1.639}{7.9701}{mi}{110}%
}\hskip.1em\allowbreak{} ，并置\allowbreak{}\hskip.1em \RB{30.009}{8.069}{2.306}{%
\RG{0.000}{-0.002}{10.9091}{mi}{100}\RG{5.678}{-1.639}{7.9701}{rm}{48}\RG{13.440}{-0.002}{10.9091}{rm}{61}\RG{24.938}{-0.002}{10.9091}{rm}{49}%
}\hskip.1em\allowbreak{} 。沿第一行展开，再沿第二项的余子式的第一列展开，得到

\[
\RB{263.455}{8.686}{3.223}{%
\RG{0.000}{0.004}{10.9091}{mi}{100}\RG{13.440}{0.004}{10.9091}{rm}{61}\RG{24.938}{0.004}{10.9091}{rm}{49}\RG{30.410}{0.004}{10.9091}{mi}{59}\RG{57.061}{0.004}{10.9091}{mi}{100}\RG{70.518}{0.004}{10.9091}{rm}{61}\RG{82.033}{0.004}{10.9091}{mi}{97}\RG{87.793}{0.004}{10.9091}{mi}{59}\RG{114.444}{0.004}{10.9091}{mi}{100}\RG{128.811}{0.004}{10.9091}{rm}{61}\RG{140.326}{0.004}{10.9091}{mi}{97}\RG{146.086}{0.004}{10.9091}{mi}{100}\RG{170.652}{0.004}{10.9091}{sy}{0}\RG{181.561}{0.004}{10.9091}{mi}{98}\RG{186.242}{0.004}{10.9091}{mi}{99}\RG{192.765}{0.004}{10.9091}{mi}{100}\RG{225.827}{0.004}{10.9091}{rm}{40}\RG{230.069}{0.004}{10.9091}{mi}{110}\RG{239.648}{0.004}{10.9091}{sy}{21}\RG{251.163}{0.004}{10.9091}{rm}{50}\RG{256.618}{0.004}{10.9091}{rm}{41}\RG{260.860}{0.004}{10.9091}{mi}{58}\RG{5.678}{-1.633}{7.9701}{rm}{48}\RG{62.756}{-1.633}{7.9701}{rm}{49}\RG{120.144}{-1.633}{7.9701}{mi}{110}\RG{151.770}{-1.633}{7.9701}{mi}{110}\RG{156.909}{-1.633}{7.9701}{sy}{0}\RG{163.495}{-1.633}{7.9701}{rm}{49}\RG{198.461}{-1.633}{7.9701}{mi}{110}\RG{203.599}{-1.633}{7.9701}{sy}{0}\RG{210.185}{-1.633}{7.9701}{rm}{50}%
}
\]
由此可得

\[
\RB{149.891}{21.136}{14.021}{%
\RG{78.487}{15.381}{10.9092}{ex}{18}\RG{112.421}{15.381}{10.9092}{ex}{19}\RG{25.859}{14.658}{7.9701}{sy}{98}\RG{29.623}{14.658}{7.9701}{mi}{110}\RG{34.756}{14.658}{7.9701}{mi}{61}\RG{38.996}{14.658}{7.9701}{rm}{50}\RG{43.230}{14.658}{7.9701}{sy}{99}\RG{28.548}{10.362}{10.9092}{ex}{88}\RG{86.517}{7.379}{10.9091}{mi}{110}\RG{95.489}{7.379}{10.9091}{sy}{0}\RG{106.398}{7.379}{10.9091}{mi}{107}\RG{73.367}{4.503}{7.9701}{mi}{107}\RG{126.218}{4.503}{7.9701}{mi}{110}\RG{131.356}{4.503}{7.9701}{sy}{0}\RG{137.943}{4.503}{7.9701}{rm}{50}\RG{142.177}{4.503}{7.9701}{mi}{107}\RG{0.000}{-0.001}{10.9091}{mi}{100}\RG{14.344}{-0.001}{10.9091}{rm}{61}\RG{46.993}{-0.001}{10.9091}{rm}{40}\RG{51.236}{-0.001}{10.9091}{sy}{0}\RG{59.721}{-0.001}{10.9091}{mi}{98}\RG{64.402}{-0.001}{10.9091}{mi}{99}\RG{69.125}{-0.001}{10.9091}{rm}{41}\RG{120.451}{-0.001}{10.9091}{mi}{97}\RG{147.296}{-0.001}{10.9091}{mi}{58}\RG{5.678}{-1.638}{7.9701}{mi}{110}\RG{96.458}{-7.484}{10.9091}{mi}{107}\RG{28.705}{-13.354}{7.9701}{mi}{107}\RG{33.327}{-13.354}{7.9701}{rm}{61}\RG{39.911}{-13.354}{7.9701}{rm}{48}%
}
\]
为核对这一表达式，记右端为\allowbreak{}\hskip.1em \RB{10.491}{5.321}{2.217}{%
\RG{0.000}{-0.001}{10.9091}{mi}{115}\RG{5.113}{-1.637}{7.9701}{mi}{110}%
}\hskip.1em\allowbreak{} ，则\allowbreak{}\hskip.1em \RB{30.295}{7.765}{2.304}{%
\RG{0.000}{-0.001}{10.9091}{mi}{115}\RG{5.114}{-1.637}{7.9701}{rm}{48}\RG{13.300}{-0.001}{10.9091}{rm}{61}\RG{25.224}{-0.001}{10.9091}{rm}{49}%
}\hskip.1em\allowbreak{} 、\allowbreak{}\hskip.1em \RB{31.171}{5.321}{2.137}{%
\RG{0.000}{-0.001}{10.9091}{mi}{115}\RG{5.113}{-1.637}{7.9701}{rm}{49}\RG{13.299}{-0.001}{10.9091}{rm}{61}\RG{25.238}{-0.001}{10.9091}{mi}{97}%
}\hskip.1em\allowbreak{} 。在\allowbreak{}\hskip.1em \RB{73.371}{8.070}{2.217}{%
\RG{0.000}{-0.001}{10.9091}{mi}{97}\RG{5.760}{-0.001}{10.9091}{mi}{115}\RG{10.880}{-1.637}{7.9701}{mi}{110}\RG{16.018}{-1.637}{7.9701}{sy}{0}\RG{22.605}{-1.637}{7.9701}{rm}{49}\RG{29.931}{-0.001}{10.9091}{sy}{0}\RG{41.010}{-0.001}{10.9091}{mi}{98}\RG{45.690}{-0.001}{10.9091}{mi}{99}\RG{52.214}{-0.001}{10.9091}{mi}{115}\RG{57.344}{-1.637}{7.9701}{mi}{110}\RG{62.483}{-1.637}{7.9701}{sy}{0}\RG{69.070}{-1.637}{7.9701}{rm}{50}%
}\hskip.1em\allowbreak{} 中，\allowbreak{}\hskip.1em \RB{57.967}{9.989}{3.228}{%
\RG{0.000}{-0.001}{10.9091}{rm}{40}\RG{4.242}{-0.001}{10.9091}{sy}{0}\RG{12.727}{-0.001}{10.9091}{mi}{98}\RG{17.407}{-0.001}{10.9091}{mi}{99}\RG{22.130}{-0.001}{10.9091}{rm}{41}\RG{26.373}{3.958}{7.9701}{mi}{107}\RG{31.492}{-0.001}{10.9091}{mi}{97}\RG{37.259}{3.958}{7.9701}{mi}{110}\RG{42.397}{3.958}{7.9701}{sy}{0}\RG{48.983}{3.958}{7.9701}{rm}{50}\RG{53.218}{3.958}{7.9701}{mi}{107}%
}\hskip.1em\allowbreak{} 的系数为

\[
\RB{194.065}{8.942}{18.229}{%
\RG{0.000}{8.006}{10.9092}{ex}{18}\RG{52.723}{8.006}{10.9092}{ex}{19}\RG{74.087}{8.006}{10.9092}{ex}{18}\RG{126.809}{8.006}{10.9092}{ex}{19}\RG{149.385}{8.006}{10.9092}{ex}{18}\RG{183.320}{8.006}{10.9092}{ex}{19}\RG{8.031}{0.004}{10.9091}{mi}{110}\RG{17.003}{0.004}{10.9091}{sy}{0}\RG{27.912}{0.004}{10.9091}{rm}{49}\RG{35.791}{0.004}{10.9091}{sy}{0}\RG{46.700}{0.004}{10.9091}{mi}{107}\RG{82.118}{0.004}{10.9091}{mi}{110}\RG{91.090}{0.004}{10.9091}{sy}{0}\RG{101.999}{0.004}{10.9091}{rm}{49}\RG{109.878}{0.004}{10.9091}{sy}{0}\RG{120.787}{0.004}{10.9091}{mi}{107}\RG{157.415}{0.004}{10.9091}{mi}{110}\RG{166.388}{0.004}{10.9091}{sy}{0}\RG{177.297}{0.004}{10.9091}{mi}{107}\RG{63.177}{-7.376}{10.9091}{rm}{43}\RG{137.870}{-7.376}{10.9091}{rm}{61}\RG{191.351}{-7.376}{10.9091}{mi}{59}\RG{27.365}{-14.859}{10.9091}{mi}{107}\RG{92.058}{-14.859}{10.9091}{mi}{107}\RG{100.505}{-14.859}{10.9091}{sy}{0}\RG{111.414}{-14.859}{10.9091}{rm}{49}\RG{167.356}{-14.859}{10.9091}{mi}{107}%
}
\]
其中超出范围的二项式系数取零。故\allowbreak{}\hskip.1em \RB{10.492}{5.322}{2.216}{%
\RG{0.000}{0.001}{10.9091}{mi}{115}\RG{5.114}{-1.636}{7.9701}{mi}{110}%
}\hskip.1em\allowbreak{} 满足同一递推关系和初值，由归纳法知\allowbreak{}\hskip.1em \RB{36.351}{8.071}{2.216}{%
\RG{0.000}{-0.000}{10.9091}{mi}{115}\RG{5.113}{-1.636}{7.9701}{mi}{110}\RG{13.780}{0.001}{10.9091}{rm}{61}\RG{25.295}{-0.000}{10.9091}{mi}{100}\RG{30.973}{-1.636}{7.9701}{mi}{110}%
}\hskip.1em\allowbreak{} 。

\noindent \RB{304.614}{36.557}{17.396}{%
\RT{0.000}{27.329}{10.9091}{因}\RT{10.909}{27.329}{10.9091}{此}\RT{21.818}{27.329}{10.9091}{所}\RT{32.728}{27.329}{10.9091}{求}\RT{43.637}{27.329}{10.9091}{行}\RT{54.546}{27.329}{10.9091}{列}\RT{65.455}{27.329}{10.9091}{式}\RT{76.364}{27.329}{10.9091}{为}\RG{159.486}{21.929}{10.9092}{ex}{50}\RG{288.196}{21.929}{10.9092}{ex}{51}\RG{295.469}{19.211}{7.9701}{rm}{50}\RG{219.387}{15.384}{10.9092}{ex}{18}\RG{253.321}{15.384}{10.9092}{ex}{19}\RG{166.759}{14.661}{7.9701}{sy}{98}\RG{170.523}{14.661}{7.9701}{mi}{110}\RG{175.656}{14.661}{7.9701}{mi}{61}\RG{179.895}{14.661}{7.9701}{rm}{50}\RG{184.130}{14.661}{7.9701}{sy}{99}\RG{169.447}{10.365}{10.9092}{ex}{88}\RG{227.417}{7.382}{10.9091}{mi}{110}\RG{236.389}{7.382}{10.9091}{sy}{0}\RG{247.298}{7.382}{10.9091}{mi}{107}\RG{214.267}{4.506}{7.9701}{mi}{107}\RG{267.118}{4.506}{7.9701}{mi}{110}\RG{272.256}{4.506}{7.9701}{sy}{0}\RG{278.842}{4.506}{7.9701}{rm}{50}\RG{283.076}{4.506}{7.9701}{mi}{107}\RG{159.486}{2.293}{10.9092}{ex}{52}\RG{288.196}{2.293}{10.9092}{ex}{53}\RG{116.199}{0.002}{10.9091}{rm}{100}\RG{122.254}{0.002}{10.9091}{rm}{101}\RG{127.097}{0.002}{10.9091}{rm}{116}\RG{133.168}{0.002}{10.9091}{mi}{77}\RG{147.971}{0.002}{10.9091}{rm}{61}\RG{187.894}{0.002}{10.9091}{rm}{40}\RG{192.137}{0.002}{10.9091}{sy}{0}\RG{200.622}{0.002}{10.9091}{mi}{98}\RG{205.302}{0.002}{10.9091}{mi}{99}\RG{210.024}{0.002}{10.9091}{rm}{41}\RG{261.351}{0.002}{10.9091}{mi}{97}\RG{302.019}{0.002}{10.9091}{mi}{58}\RG{237.357}{-7.481}{10.9091}{mi}{107}\RG{169.605}{-13.351}{7.9701}{mi}{107}\RG{174.227}{-13.351}{7.9701}{rm}{61}\RG{180.811}{-13.351}{7.9701}{rm}{48}%
}

这一多项式公式对\allowbreak{}\hskip.1em \RB{25.307}{8.069}{2.608}{%
\RG{0.000}{-0.002}{10.9091}{mi}{97}\RG{5.760}{-0.002}{10.9091}{mi}{59}\RG{10.593}{-0.002}{10.9091}{mi}{98}\RG{15.273}{-0.002}{10.9091}{mi}{59}\RG{20.117}{-0.002}{10.9091}{mi}{99}%
}\hskip.1em\allowbreak{} 的退化取值也成立。

小结：先识别五对角矩阵为三对角矩阵的平方，再将后者的行列式化为二阶递推。用有限和表示结果，可避免特征方程重根时另作讨论。

题目来源：谢启鸿《复旦大学高等代数习题集（2023版）》，第81页。

% Original page 7.
\section*{2019A06}
证明：任意\allowbreak{}\hskip.1em \RB{6.729}{5.322}{0.620}{%
\RG{0.000}{-0.000}{10.9091}{mi}{110}%
}\hskip.1em\allowbreak{} 阶非异矩阵\allowbreak{}\hskip.1em \RB{8.366}{8.311}{0.500}{%
\RG{0.000}{-0.000}{10.9091}{mi}{65}%
}\hskip.1em\allowbreak{} ，仅通过第三类初等行变换，即把某一行的倍数加到另一行，就可变为

\[
\RB{96.383}{8.678}{3.231}{%
\RG{0.000}{-0.004}{10.9091}{rm}{100}\RG{6.055}{-0.004}{10.9091}{rm}{105}\RG{9.076}{-0.004}{10.9091}{rm}{97}\RG{14.531}{-0.004}{10.9091}{rm}{103}\RG{20.127}{-0.004}{10.9091}{rm}{40}\RG{24.360}{-0.004}{10.9091}{rm}{49}\RG{29.848}{-0.004}{10.9091}{mi}{59}\RG{34.681}{-0.004}{10.9091}{mi}{58}\RG{39.525}{-0.004}{10.9091}{mi}{58}\RG{44.368}{-0.004}{10.9091}{mi}{58}\RG{49.201}{-0.004}{10.9091}{mi}{59}\RG{54.091}{-0.004}{10.9091}{rm}{49}\RG{59.545}{-0.004}{10.9091}{mi}{59}\RG{64.395}{-0.004}{10.9091}{rm}{100}\RG{70.449}{-0.004}{10.9091}{rm}{101}\RG{75.293}{-0.004}{10.9091}{rm}{116}\RG{81.364}{-0.004}{10.9091}{mi}{65}\RG{89.546}{-0.004}{10.9091}{rm}{41}\RG{93.789}{-0.004}{10.9091}{mi}{58}%
}
\]
证明：\allowbreak{}\hskip.1em \RB{26.147}{7.766}{0.620}{%
\RG{0.000}{0.000}{10.9091}{mi}{110}\RG{9.578}{0.000}{10.9091}{rm}{61}\RG{21.076}{0.000}{10.9091}{rm}{49}%
}\hskip.1em\allowbreak{} 时无须变换。以下设\allowbreak{}\hskip.1em \RB{26.492}{7.766}{1.995}{%
\RG{0.000}{0.000}{10.9091}{mi}{110}\RG{9.579}{0.000}{10.9091}{sy}{21}\RG{21.094}{0.000}{10.9091}{rm}{50}%
}\hskip.1em\allowbreak{} 。

先说明仅用第三类初等行变换即可把\allowbreak{}\hskip.1em \RB{8.366}{8.314}{0.500}{%
\RG{0.000}{0.003}{10.9091}{mi}{65}%
}\hskip.1em\allowbreak{} 化为非异对角矩阵。按列进行消元：假设前\allowbreak{}\hskip.1em \RB{24.719}{8.074}{0.617}{%
\RG{0.000}{0.003}{10.9091}{mi}{107}\RG{8.593}{0.003}{10.9091}{sy}{0}\RG{19.648}{0.003}{10.9091}{rm}{49}%
}\hskip.1em\allowbreak{} 列的主元均非零，且其下方元素已消为零，此时矩阵具有分块形式

\[
\RB{48.984}{0.937}{32.779}{%
\RG{0.000}{0.001}{10.9092}{ex}{32}\RG{35.814}{0.001}{10.9092}{ex}{33}\RG{8.637}{-11.373}{10.9091}{mi}{85}\RG{28.799}{-11.373}{10.9091}{sy}{3}\RG{46.269}{-18.653}{10.9091}{mi}{59}\RG{10.228}{-26.548}{10.9091}{rm}{48}\RG{27.237}{-26.548}{10.9091}{mi}{67}%
}
\]
其中\allowbreak{}\hskip.1em \RB{8.791}{7.947}{0.744}{%
\RG{0.000}{-0.004}{10.9091}{mi}{85}%
}\hskip.1em\allowbreak{} 是对角元非零的上三角矩阵。由于整个矩阵非异，\allowbreak{}\hskip.1em \RB{8.791}{8.187}{0.744}{%
\RG{0.000}{-0.004}{10.9091}{mi}{67}%
}\hskip.1em\allowbreak{} 也非异，故\allowbreak{}\hskip.1em \RB{8.791}{8.187}{0.744}{%
\RG{0.000}{-0.004}{10.9091}{mi}{67}%
}\hskip.1em\allowbreak{} 的第一列不全为零。若当前第\allowbreak{}\hskip.1em \RB{24.790}{8.683}{3.226}{%
\RG{0.000}{0.001}{10.9091}{rm}{40}\RG{4.243}{0.001}{10.9091}{mi}{107}\RG{10.254}{0.001}{10.9091}{mi}{59}\RG{15.098}{0.001}{10.9091}{mi}{107}\RG{21.137}{0.001}{10.9091}{rm}{41}%
}\hskip.1em\allowbreak{} 元为零，就从第\allowbreak{}\hskip.1em \RB{6.042}{8.072}{0.619}{%
\RG{0.000}{0.001}{10.9091}{mi}{107}%
}\hskip.1em\allowbreak{} 行以下选取该列非零的一行加到第\allowbreak{}\hskip.1em \RB{6.042}{8.072}{0.619}{%
\RG{0.000}{0.001}{10.9091}{mi}{107}%
}\hskip.1em\allowbreak{} 行，使主元非零；随后把第\allowbreak{}\hskip.1em \RB{6.042}{8.067}{0.624}{%
\RG{0.000}{-0.004}{10.9091}{mi}{107}%
}\hskip.1em\allowbreak{} 行的适当倍数加到下面各行，消去该主元下方的元素。这些行的前\allowbreak{}\hskip.1em \RB{23.978}{8.067}{0.624}{%
\RG{0.000}{-0.004}{10.9091}{mi}{107}\RG{8.222}{-0.004}{10.9091}{sy}{0}\RG{18.907}{-0.004}{10.9091}{rm}{49}%
}\hskip.1em\allowbreak{} 个元素均为零，因而不破坏此前的消元结果。

依次进行，得到对角元均非零的上三角矩阵；再从最后一列开始，用相应主元行消去主元上方的元素，得到

\[
\RB{156.805}{8.686}{3.223}{%
\RG{0.000}{0.004}{10.9091}{mi}{68}\RG{12.365}{0.004}{10.9091}{rm}{61}\RG{23.863}{0.004}{10.9091}{rm}{100}\RG{29.918}{0.004}{10.9091}{rm}{105}\RG{32.940}{0.004}{10.9091}{rm}{97}\RG{38.394}{0.004}{10.9091}{rm}{103}\RG{44.002}{0.004}{10.9091}{rm}{40}\RG{48.274}{0.004}{10.9091}{mi}{100}\RG{58.684}{0.004}{10.9091}{mi}{59}\RG{63.517}{0.004}{10.9091}{mi}{58}\RG{68.361}{0.004}{10.9091}{mi}{58}\RG{73.205}{0.004}{10.9091}{mi}{58}\RG{78.037}{0.004}{10.9091}{mi}{59}\RG{82.881}{0.004}{10.9091}{mi}{100}\RG{94.242}{0.004}{10.9091}{rm}{41}\RG{98.485}{0.004}{10.9091}{mi}{59}\RG{125.136}{0.004}{10.9091}{mi}{100}\RG{137.241}{0.004}{10.9091}{sy}{54}\RG{137.241}{0.004}{10.9091}{rm}{61}\RG{148.739}{0.004}{10.9091}{rm}{48}\RG{154.210}{0.004}{10.9091}{mi}{58}\RG{53.952}{-1.632}{7.9701}{rm}{49}\RG{88.606}{-1.632}{7.9701}{mi}{110}\RG{130.829}{-1.632}{7.9701}{mi}{105}%
}
\]
所用行变换均不改变行列式，故

\[
\RB{80.380}{8.309}{2.218}{%
\RG{0.000}{-0.002}{10.9091}{mi}{100}\RG{12.229}{-0.002}{10.9091}{sy}{1}\RG{17.062}{-0.002}{10.9091}{sy}{1}\RG{21.906}{-0.002}{10.9091}{sy}{1}\RG{26.774}{-0.002}{10.9091}{mi}{100}\RG{41.119}{-0.002}{10.9091}{rm}{61}\RG{52.618}{-0.002}{10.9091}{rm}{100}\RG{58.672}{-0.002}{10.9091}{rm}{101}\RG{63.516}{-0.002}{10.9091}{rm}{116}\RG{69.604}{-0.002}{10.9091}{mi}{65}\RG{77.785}{-0.002}{10.9091}{mi}{58}\RG{5.678}{-1.638}{7.9701}{rm}{49}\RG{32.452}{-1.638}{7.9701}{mi}{110}%
}
\]
下面仍只用第三类初等行变换，把前\allowbreak{}\hskip.1em \RB{24.986}{7.762}{0.624}{%
\RG{0.000}{-0.004}{10.9091}{mi}{110}\RG{8.989}{-0.004}{10.9091}{sy}{0}\RG{19.915}{-0.004}{10.9091}{rm}{49}%
}\hskip.1em\allowbreak{} 个对角元依次化为\allowbreak{}\hskip.1em \RB{5.071}{7.762}{0.504}{%
\RG{0.000}{-0.004}{10.9091}{rm}{49}%
}\hskip.1em\allowbreak{} 。对任意非零\allowbreak{}\hskip.1em \RB{15.697}{8.067}{2.609}{%
\RG{0.000}{-0.004}{10.9091}{mi}{100}\RG{5.673}{-0.004}{10.9091}{mi}{59}\RG{10.506}{-0.004}{10.9091}{mi}{101}%
}\hskip.1em\allowbreak{} ，二阶对角块可作如下变换：

\[
\RB{203.811}{135.711}{14.123}{%
\RG{0.000}{134.774}{10.9092}{ex}{32}\RG{29.732}{134.774}{10.9092}{ex}{33}\RG{119.188}{134.774}{10.9092}{ex}{32}\RG{157.181}{134.774}{10.9092}{ex}{33}\RG{88.089}{125.253}{5.9776}{sy}{0}\RG{93.846}{125.253}{5.9776}{rm}{49}\RG{8.636}{123.399}{10.9091}{mi}{100}\RG{24.277}{123.399}{10.9091}{rm}{48}\RG{131.955}{123.399}{10.9091}{mi}{100}\RG{151.726}{123.399}{10.9091}{rm}{48}\RG{47.618}{122.440}{7.9701}{mi}{82}\RG{58.147}{122.440}{7.9701}{sy}{32}\RG{66.617}{122.440}{7.9701}{mi}{82}\RG{77.146}{122.440}{7.9701}{sy}{0}\RG{83.732}{122.440}{7.9701}{mi}{100}\RG{97.997}{122.440}{7.9701}{mi}{82}\RG{53.996}{121.334}{5.9776}{rm}{50}\RG{72.995}{121.334}{5.9776}{rm}{50}\RG{104.375}{121.334}{5.9776}{rm}{49}\RG{41.398}{116.119}{10.9091}{sy}{0}\RG{46.046}{116.119}{10.9091}{sy}{0}\RG{52.100}{116.119}{10.9091}{sy}{0}\RG{58.155}{116.119}{10.9091}{sy}{0}\RG{64.199}{116.119}{10.9091}{sy}{0}\RG{70.253}{116.119}{10.9091}{sy}{0}\RG{76.308}{116.119}{10.9091}{sy}{0}\RG{82.362}{116.119}{10.9091}{sy}{0}\RG{88.417}{116.119}{10.9091}{sy}{0}\RG{94.461}{116.119}{10.9091}{sy}{0}\RG{100.515}{116.119}{10.9091}{sy}{0}\RG{105.163}{116.119}{10.9091}{sy}{33}\RG{8.748}{108.224}{10.9091}{rm}{48}\RG{24.464}{108.224}{10.9091}{mi}{101}\RG{127.824}{108.224}{10.9091}{sy}{0}\RG{136.309}{108.224}{10.9091}{rm}{49}\RG{151.914}{108.224}{10.9091}{mi}{101}\RG{126.688}{96.069}{10.9092}{ex}{32}\RG{197.257}{96.069}{10.9092}{ex}{33}\RG{47.618}{84.621}{7.9701}{mi}{82}\RG{58.148}{84.621}{7.9701}{sy}{32}\RG{66.617}{84.621}{7.9701}{mi}{82}\RG{77.146}{84.621}{7.9701}{rm}{43}\RG{83.729}{84.621}{7.9701}{rm}{40}\RG{87.025}{84.621}{7.9701}{mi}{100}\RG{91.383}{84.621}{7.9701}{sy}{0}\RG{97.969}{84.621}{7.9701}{rm}{49}\RG{102.201}{84.621}{7.9701}{rm}{41}\RG{105.497}{84.621}{7.9701}{mi}{82}\RG{139.566}{84.694}{10.9091}{rm}{49}\RG{159.225}{84.694}{10.9091}{rm}{40}\RG{163.468}{84.694}{10.9091}{mi}{100}\RG{171.571}{84.694}{10.9091}{sy}{0}\RG{182.481}{84.694}{10.9091}{rm}{49}\RG{187.935}{84.694}{10.9091}{rm}{41}\RG{192.178}{84.694}{10.9091}{mi}{101}\RG{53.996}{83.514}{5.9776}{rm}{49}\RG{72.995}{83.514}{5.9776}{rm}{49}\RG{111.875}{83.514}{5.9776}{rm}{50}\RG{41.398}{77.414}{10.9091}{sy}{0}\RG{46.766}{77.414}{10.9091}{sy}{0}\RG{52.820}{77.414}{10.9091}{sy}{0}\RG{58.875}{77.414}{10.9091}{sy}{0}\RG{64.919}{77.414}{10.9091}{sy}{0}\RG{70.973}{77.414}{10.9091}{sy}{0}\RG{77.028}{77.414}{10.9091}{sy}{0}\RG{83.082}{77.414}{10.9091}{sy}{0}\RG{89.137}{77.414}{10.9091}{sy}{0}\RG{95.181}{77.414}{10.9091}{sy}{0}\RG{101.235}{77.414}{10.9091}{sy}{0}\RG{107.290}{77.414}{10.9091}{sy}{0}\RG{112.646}{77.414}{10.9091}{sy}{33}\RG{135.324}{69.519}{10.9091}{sy}{0}\RG{143.809}{69.519}{10.9091}{rm}{49}\RG{175.701}{69.519}{10.9091}{mi}{101}\RG{104.923}{57.364}{10.9092}{ex}{32}\RG{167.007}{57.364}{10.9092}{ex}{33}\RG{113.559}{45.988}{10.9091}{rm}{49}\RG{128.974}{45.988}{10.9091}{rm}{40}\RG{133.219}{45.988}{10.9091}{mi}{100}\RG{141.321}{45.988}{10.9091}{sy}{0}\RG{152.230}{45.988}{10.9091}{rm}{49}\RG{157.685}{45.988}{10.9091}{rm}{41}\RG{161.927}{45.988}{10.9091}{mi}{101}\RG{47.618}{45.029}{7.9701}{mi}{82}\RG{58.147}{45.029}{7.9701}{sy}{32}\RG{66.617}{45.029}{7.9701}{mi}{82}\RG{77.146}{45.029}{7.9701}{rm}{43}\RG{83.732}{45.029}{7.9701}{mi}{82}\RG{53.996}{43.922}{5.9776}{rm}{50}\RG{72.995}{43.922}{5.9776}{rm}{50}\RG{90.110}{43.922}{5.9776}{rm}{49}\RG{41.398}{38.708}{10.9091}{sy}{0}\RG{44.977}{38.708}{10.9091}{sy}{0}\RG{51.031}{38.708}{10.9091}{sy}{0}\RG{57.075}{38.708}{10.9091}{sy}{0}\RG{63.129}{38.708}{10.9091}{sy}{0}\RG{69.184}{38.708}{10.9091}{sy}{0}\RG{75.239}{38.708}{10.9091}{sy}{0}\RG{81.293}{38.708}{10.9091}{sy}{0}\RG{87.337}{38.708}{10.9091}{sy}{0}\RG{90.915}{38.708}{10.9091}{sy}{33}\RG{113.559}{30.813}{10.9091}{rm}{48}\RG{142.613}{30.813}{10.9091}{mi}{100}\RG{148.286}{30.813}{10.9091}{mi}{101}\RG{136.595}{18.657}{10.9092}{ex}{32}\RG{171.406}{18.657}{10.9092}{ex}{33}\RG{91.383}{10.022}{5.9776}{sy}{0}\RG{97.139}{10.022}{5.9776}{rm}{49}\RG{145.231}{7.282}{10.9091}{rm}{49}\RG{163.297}{7.282}{10.9091}{rm}{48}\RG{47.618}{7.210}{7.9701}{mi}{82}\RG{58.148}{7.210}{7.9701}{sy}{32}\RG{66.617}{7.210}{7.9701}{mi}{82}\RG{77.146}{7.210}{7.9701}{rm}{43}\RG{83.729}{7.210}{7.9701}{rm}{40}\RG{87.025}{7.210}{7.9701}{mi}{100}\RG{101.290}{7.210}{7.9701}{sy}{0}\RG{107.876}{7.210}{7.9701}{rm}{49}\RG{112.108}{7.210}{7.9701}{rm}{41}\RG{115.404}{7.210}{7.9701}{mi}{82}\RG{53.996}{6.103}{5.9776}{rm}{49}\RG{72.995}{6.103}{5.9776}{rm}{49}\RG{121.782}{6.103}{5.9776}{rm}{50}\RG{41.398}{0.003}{10.9091}{sy}{0}\RG{45.664}{0.003}{10.9091}{sy}{0}\RG{51.708}{0.003}{10.9091}{sy}{0}\RG{57.762}{0.003}{10.9091}{sy}{0}\RG{63.817}{0.003}{10.9091}{sy}{0}\RG{69.871}{0.003}{10.9091}{sy}{0}\RG{75.926}{0.003}{10.9091}{sy}{0}\RG{81.970}{0.003}{10.9091}{sy}{0}\RG{88.024}{0.003}{10.9091}{sy}{0}\RG{94.079}{0.003}{10.9091}{sy}{0}\RG{100.133}{0.003}{10.9091}{sy}{0}\RG{106.177}{0.003}{10.9091}{sy}{0}\RG{112.232}{0.003}{10.9091}{sy}{0}\RG{118.286}{0.003}{10.9091}{sy}{0}\RG{122.541}{0.003}{10.9091}{sy}{33}\RG{181.861}{0.003}{10.9091}{mi}{58}\RG{145.231}{-7.893}{10.9091}{rm}{48}\RG{160.648}{-7.893}{10.9091}{mi}{100}\RG{166.321}{-7.893}{10.9091}{mi}{101}%
}
\]
若某一步的系数为零，直接略去该步即可。

对\allowbreak{}\hskip.1em \RB{74.551}{7.770}{2.601}{%
\RG{0.000}{0.004}{10.9091}{mi}{105}\RG{7.359}{0.004}{10.9091}{rm}{61}\RG{19.436}{0.004}{10.9091}{rm}{49}\RG{24.899}{0.004}{10.9091}{mi}{59}\RG{29.732}{0.004}{10.9091}{mi}{58}\RG{34.576}{0.004}{10.9091}{mi}{58}\RG{39.408}{0.004}{10.9091}{mi}{58}\RG{44.252}{0.004}{10.9091}{mi}{59}\RG{49.096}{0.004}{10.9091}{mi}{110}\RG{58.341}{0.004}{10.9091}{sy}{0}\RG{69.480}{0.004}{10.9091}{rm}{49}%
}\hskip.1em\allowbreak{} ，把上述四步依次用于第\allowbreak{}\hskip.1em \RB{15.315}{7.715}{2.601}{%
\RG{0.000}{0.004}{10.9091}{mi}{105}\RG{3.753}{0.004}{10.9091}{mi}{59}\RG{8.586}{0.004}{10.9091}{mi}{110}%
}\hskip.1em\allowbreak{} 两行以及相应对角块，就把第\allowbreak{}\hskip.1em \RB{3.696}{7.715}{0.616}{%
\RG{0.000}{0.004}{10.9091}{mi}{105}%
}\hskip.1em\allowbreak{} 个对角元变为\allowbreak{}\hskip.1em \RB{5.071}{7.765}{0.501}{%
\RG{0.000}{-0.001}{10.9091}{rm}{49}%
}\hskip.1em\allowbreak{} ，并将原来的\allowbreak{}\hskip.1em \RB{8.706}{8.070}{2.217}{%
\RG{0.000}{-0.001}{10.9091}{mi}{100}\RG{5.679}{-1.637}{7.9701}{mi}{105}%
}\hskip.1em\allowbreak{} 乘入最后一个对角元。其他行不受影响。因此最终得到

\[
\RB{217.654}{8.677}{3.232}{%
\RG{0.000}{-0.005}{10.9091}{rm}{100}\RG{6.055}{-0.005}{10.9091}{rm}{105}\RG{9.076}{-0.005}{10.9091}{rm}{97}\RG{14.531}{-0.005}{10.9091}{rm}{103}\RG{20.127}{-0.005}{10.9091}{rm}{40}\RG{24.360}{-0.005}{10.9091}{rm}{49}\RG{29.848}{-0.005}{10.9091}{mi}{59}\RG{34.681}{-0.005}{10.9091}{mi}{58}\RG{39.525}{-0.005}{10.9091}{mi}{58}\RG{44.368}{-0.005}{10.9091}{mi}{58}\RG{49.201}{-0.005}{10.9091}{mi}{59}\RG{54.091}{-0.005}{10.9091}{rm}{49}\RG{59.545}{-0.005}{10.9091}{mi}{59}\RG{64.378}{-0.005}{10.9091}{mi}{100}\RG{76.623}{-0.005}{10.9091}{sy}{1}\RG{81.455}{-0.005}{10.9091}{sy}{1}\RG{86.299}{-0.005}{10.9091}{sy}{1}\RG{91.169}{-0.005}{10.9091}{mi}{100}\RG{102.483}{-0.005}{10.9091}{rm}{41}\RG{109.737}{-0.005}{10.9091}{rm}{61}\RG{121.247}{-0.005}{10.9091}{rm}{100}\RG{127.301}{-0.005}{10.9091}{rm}{105}\RG{130.323}{-0.005}{10.9091}{rm}{97}\RG{135.778}{-0.005}{10.9091}{rm}{103}\RG{141.385}{-0.005}{10.9091}{rm}{40}\RG{145.618}{-0.005}{10.9091}{rm}{49}\RG{151.119}{-0.005}{10.9091}{mi}{59}\RG{155.952}{-0.005}{10.9091}{mi}{58}\RG{160.796}{-0.005}{10.9091}{mi}{58}\RG{165.639}{-0.005}{10.9091}{mi}{58}\RG{170.472}{-0.005}{10.9091}{mi}{59}\RG{175.362}{-0.005}{10.9091}{rm}{49}\RG{180.816}{-0.005}{10.9091}{mi}{59}\RG{185.666}{-0.005}{10.9091}{rm}{100}\RG{191.720}{-0.005}{10.9091}{rm}{101}\RG{196.564}{-0.005}{10.9091}{rm}{116}\RG{202.635}{-0.005}{10.9091}{mi}{65}\RG{210.817}{-0.005}{10.9091}{rm}{41}\RG{215.060}{-0.005}{10.9091}{mi}{58}\RG{70.073}{-1.641}{7.9701}{rm}{49}\RG{96.847}{-1.641}{7.9701}{mi}{110}%
}
\]
全过程只使用第三类初等行变换。\allowbreak{}\hskip.1em \RB{8.366}{7.996}{0.500}{%
\RG{0.000}{0.001}{10.9091}{am}{3}%
}\hskip.1em\allowbreak{}

小结：零主元可通过加行消除；对角化后，四次加行把一个非零对角元归一，并把其因子转移到最后一个对角元。

题目来源：谢启鸿《复旦大学高等代数习题集（2023版）》，第81页。


\end{document}
